% =============================================================================
% GLOSSARY OF TERMS — Alert Aggregation in Microservices
% Standalone file for % =============================================================================
% GLOSSARY OF TERMS — Alert Aggregation in Microservices
% Standalone file for % =============================================================================
% GLOSSARY OF TERMS — Alert Aggregation in Microservices
% Standalone file for % =============================================================================
% GLOSSARY OF TERMS — Alert Aggregation in Microservices
% Standalone file for \input{glossary} in main IEEEtran document
% Compiles with: pdflatex main.tex (where main.tex includes this file)
% =============================================================================

% -----------------------------------------------------------------------------
% USAGE IN MAIN DOCUMENT:
% \section*{Glossary of Terms}
% \addcontentsline{toc}{section}{Glossary of Terms}
% \input{glossary}
% -----------------------------------------------------------------------------

\begin{description}

% =============================================================================
% CORE PROBLEM TERMS
% =============================================================================

\item[Alert Fatigue]
A cognitive state experienced by on-call engineers when the volume of alerts exceeds human cognitive capacity, leading to desensitization, missed critical alerts, and delayed incident response~\cite{pagerduty2020, grafana2025}.

\item[Alert Storm]
A sudden surge of alerts triggered by a single root cause propagating through dependent microservices, resulting in hundreds to thousands of alerts within minutes~\cite{kuang2024cola}.

\item[Noisy Alert Warning]
A warning-level alert generated by the monitoring system that is technically correct but does not require meaningful remediation action from an on-call engineer, classified into four categories: \textit{Duplicate}, \textit{Transient}, \textit{Unactionable}, and \textit{Cascading} (defined below).

\item[Early Warning Phase]
The pre-failure period during which a microservice exhibits performance degradation (e.g., increased latency, resource saturation) but has not yet violated Service Level Objectives (SLOs) or caused user-facing outages.

\item[Actionable Alert]
An alert that satisfies the Google SRE principle: ``every alert sent to an on-call engineer must require immediate human intervention''~\cite{google-sre-ch10, google-sre-ch11}.

% =============================================================================
% FOUR CATEGORIES OF NOISY ALERT WARNINGS
% =============================================================================

\item[Duplicate Warning Alert]
Multiple alerts originating from the same service with identical or near-identical message templates within a short time window ($\Delta t \le 60$~seconds), representing repeated observations of the same underlying condition rather than distinct incidents.

\item[Transient Warning Alert]
A warning alert that occurs during normal system operation (i.e., before fault injection time $T_{\text{inject}}$ in RCAEval) and resolves spontaneously without human intervention, typically caused by temporary network fluctuations or brief resource spikes.

\item[Unactionable Warning Alert]
A technically accurate warning alert for which no Standard Operating Procedure (SOP) or runbook exists, leaving the on-call engineer without a defined remediation path; the alert consumes attention but yields no operational value.

\item[Cascading Warning Alert]
A warning alert emitted by a downstream service that is a symptom of an upstream root cause; remediating the alerting service does not resolve the underlying issue, which originates in a different service (the root cause service).

% =============================================================================
% METHODOLOGY: FOUR EXPERIMENTAL ARMS
% =============================================================================

\item[Arm 1: Rule-Based Aggregation (Baseline)]
An alert grouping approach that clusters alerts based on static label matching (e.g., service name, alert name, severity) and fixed time windows (\texttt{group\_wait}, \texttt{group\_interval}), as implemented in Prometheus Alertmanager~\cite{prometheus-alertmanager}. Lacks awareness of service dependencies and semantic content.

\item[Arm 2: Semantic Similarity-Based Aggregation]
A method that parses raw logs into templates (via Drain3~\cite{drain3}), encodes templates into dense vector embeddings using Sentence-BERT~\cite{sbert}, and clusters alerts using density-based algorithms (HDBSCAN/DBSCAN) based on cosine similarity. Operates without topology knowledge.

\item[Arm 3: Temporal-Spatial Correlation Aggregation]
A topology-aware method that constructs a static service dependency graph (caller-callee adjacency matrix) from microservice architecture, then correlates alerts occurring within a sliding time window ($\Delta t \in \{30, 60, 120\}$~seconds) if a dependency path of length $k \le 2$ exists between the alerting services.

\item[Arm 4: LLM-Based Agentic Aggregation (Proposed)]
A hybrid two-stage approach: (1) a fast filter removes obvious duplicates and performs time-window pre-grouping; (2) an LLM agent (GPT-4o-mini / DeepSeek-V3) reasons over the remaining clusters using Chain-of-Thought prompting and three tools: \texttt{get\_related\_alerts}, \texttt{get\_service\_dependency}, and optionally \texttt{get\_runbook}. Limited to a maximum of two tool calls per grouping decision.

% =============================================================================
% TECHNICAL CONCEPTS & ARCHITECTURE
% =============================================================================

\item[Service Topology (Service Dependency Graph)]
A directed graph $G = (V, E)$ where vertices $V$ represent microservices and edges $E$ represent synchronous (HTTP/gRPC) or asynchronous (message queue) call relationships. Used to distinguish root cause services from downstream symptom services.

\item[Semantic Similarity]
A measure of meaning-based resemblance between two log messages, computed as the cosine similarity between their dense vector embeddings (typically 384-dimensional from \texttt{sentence-transformers/all-MiniLM-L6-v2}). Captures paraphrases and synonyms that lexical matching misses.

\item[Runbook Reasoning]
The process by which an LLM agent consults a synthetic runbook (a structured troubleshooting guide for a specific fault type) to determine whether an alert is actionable, requires escalation, or can be safely aggregated as noise.

\item[Tool-Calling Agent]
An LLM agent augmented with a predefined set of functions (tools) it can invoke during reasoning: \texttt{get\_related\_alerts(service, time\_window)}, \texttt{get\_service\_dependency(service)}, and \texttt{get\_runbook(fault\_type)}. Tool outputs are fed back into the agent's context for subsequent reasoning steps.

\item[Chain-of-Thought (CoT) Prompting]
A prompting technique that instructs the LLM to generate intermediate reasoning steps before producing a final answer, improving accuracy on multi-step causal inference tasks~\cite{wei2022cot}.

\item[Fast Filter (Deduplication + Time-Window Pre-Grouping)]
A lightweight deterministic preprocessing stage that removes exact duplicate alerts and performs initial grouping by service and time window, reducing the alert volume passed to the LLM agent by an estimated 70--80\%.

\item[Synthetic Runbook]
A researcher-authored troubleshooting document for each fault type in the benchmark dataset (e.g., CPU stress, memory leak, network latency), containing symptoms, diagnosis steps, and remediation actions. Not a production SOP; used solely for experimental evaluation.

% =============================================================================
% EVALUATION METRICS
% =============================================================================

\item[Alert Reduction Ratio (ARR)]
The primary metric quantifying alert volume reduction:
\begin{equation}
\text{ARR} = \left(1 - \frac{N_{\text{groups}}}{N_{\text{raw\_alerts}}}\right) \times 100\%
\end{equation}
where $N_{\text{raw\_alerts}}$ is the number of warning alerts before aggregation and $N_{\text{groups}}$ is the number of aggregated groups presented to the engineer. Target: $\text{ARR} \ge 60\%$.

\item[Root Cause Preservation Rate (RCPR)]
The critical safety metric ensuring the root cause evidence is not lost during aggregation. It is defined at two granularities. At service granularity, which is the primary metric and spans every case that exposes logs ($M = 359$):
\begin{equation}
\text{RCPR}_{\text{svc}} = \frac{1}{M} \sum_{i=1}^{M} \mathbb{I}(s_{\text{root}}^{(i)} \in \text{PrimaryGroup}_i) \times 100\%
\end{equation}
At alert granularity, a secondary metric restricted to the $M'$ cases whose ground-truth \texttt{root\_cause.txt} indicator line is itself a warning-tier log ($M' = 4$):
\begin{equation}
\text{RCPR}_{\text{alert}} = \frac{1}{M'} \sum_{i=1}^{M'} \mathbb{I}(\text{RootCauseAlert}_i \in \text{PrimaryGroup}_i) \times 100\%
\end{equation}
where $\mathbb{I}$ is the indicator function. Target: $\text{RCPR} \ge 95\%$ (mandatory).

\item[Pairwise Grouping Precision ($P_p$)]
The fraction of alert pairs placed in the same group that truly share the same root cause:
\begin{equation}
P_p = \frac{|\{(a_x, a_y) : \text{same group} \land \text{same root cause}\}|}{|\{(a_x, a_y) : \text{same group}\}|}
\end{equation}

\item[Pairwise Grouping Recall ($R_p$)]
The fraction of alert pairs sharing the same root cause that are successfully grouped together:
\begin{equation}
R_p = \frac{|\{(a_x, a_y) : \text{same group} \land \text{same root cause}\}|}{|\{(a_x, a_y) : \text{same root cause}\}|}
\end{equation}

\item[Pairwise Grouping F1-Score ($F1_p$)]
The harmonic mean of pairwise precision and recall:
\begin{equation}
F1_p = 2 \times \frac{P_p \times R_p}{P_p + R_p}
\end{equation}
Target: $F1_p \ge 0.85$.

\item[Group Purity]
The average homogeneity of root causes within each generated group:
\begin{equation}
\text{Purity} = \frac{1}{N_{\text{groups}}} \sum_{g=1}^{N_{\text{groups}}} \frac{\max_{c} |\{a \in g : \text{root\_cause}(a) = c\}|}{|g|}
\end{equation}
Target: $\text{Purity} \ge 0.80$.

\item[Processing Latency]
The wall-clock time required to process 100 warning alerts through the aggregation pipeline, measured in milliseconds. For Arm 4, includes LLM inference time and token overhead.

% =============================================================================
% DATASET & BENCHMARK TERMS
% =============================================================================

\item[RCAEval]
A benchmark for root cause analysis in microservice systems comprising 735 failure cases across three systems (Online Boutique, Sock Shop, Train Ticket) with logs, metrics, traces, and root cause labels~\cite{pham2025rcaeval}. Licensed under MIT. Direct inspection shows that only \textbf{359 of the 735 cases ship raw log telemetry} (suite RE1 is metric-only) and that the log schema contains \texttt{timestamp}, \texttt{container\_name}, and \texttt{message} but \textbf{no severity column}.

\item[Severity Tier Assignment]
The three-tier procedure by which this work assigns the severity tier $\ell_i$, since RCAEval supplies none: \emph{Tier 1 (explicit)} parses a level token from the message text (Sock Shop, Train Ticket); \emph{Tier 2 (frequency-inferred)} flags a template as anomalous when its fault-window frequency exceeds its normal-window frequency by a ratio $R$ (Online Boutique); \emph{Tier 3 (ground-truth marker)} tags the \texttt{root\_cause.txt} indicator line as \texttt{ROOT\_CAUSE} for the eight annotated cases.

\item[LEMMA-RCA]
A large-scale multi-modal dataset for root cause analysis containing hundreds of microservice entities and four fault types (DDoS, storage failure, resource contention, noisy neighbor)~\cite{zheng2024lemma}. Licensed under CC-BY-NC-4.0 (non-commercial). Its log pipeline aggregates raw log lines into three-dimensional time series using golden-signal keywords that cover only \texttt{error}, \texttt{exception}, and \texttt{critical} --- notably \textbf{not} \texttt{warning} --- so it offers no independent warning tier.

\item[Fault Injection]
The controlled introduction of a failure (e.g., CPU stress, network latency, packet loss) into a microservice system at a known timestamp $T_{\text{inject}}$ to generate labeled failure cases for benchmarking.

\item[Normal Window]
The interval $[T_{\text{start}}, T_{\text{inject}})$ preceding fault injection, during which the system operates normally; alerts occurring here are candidates for \textit{Transient} noise. Note that the dataset fields \texttt{normal\_timesteps} and \texttt{faulty\_timesteps} are integer counts of sample points, not timestamp lists, and therefore define these windows only indirectly.

\item[Fault Window]
The interval $[T_{\text{inject}}, T_{\text{end}}]$ following fault injection, during which the system exhibits degradation; alerts here are labeled as \textit{Signal} (if from the root cause service) or \textit{Cascading} (if from a service reachable from the root cause service in the call graph).

\item[Ground Truth Root Cause Label]
The authoritative annotation of which service is the root cause of a failure case, provided by the RCAEval benchmark. Used to compute RCPR and pairwise metrics.

% =============================================================================
% BASELINE & INFRASTRUCTURE TERMS
% =============================================================================

\item[Prometheus Alertmanager]
The de facto standard alert routing and grouping component in the Cloud Native Computing Foundation (CNCF) ecosystem. Groups alerts by label matchers and time windows (\texttt{group\_wait}, \texttt{group\_interval}, \texttt{group\_by})~\cite{prometheus-alertmanager}.

\item[Group Wait (\texttt{group\_wait})]
The initial delay after the first alert in a group before sending a notification, allowing additional alerts to arrive and be grouped. Default: 30~seconds.

\item[Group Interval (\texttt{group\_interval})]
The minimum interval between subsequent notifications for the same group. Default: 5~minutes.

\item[Drain3]
An online log parsing algorithm that extracts log templates from raw log streams by clustering log messages based on token similarity, producing parameterized templates with variable placeholders~\cite{drain3}.

\item[Sentence-BERT (S-BERT)]
A modification of BERT that uses siamese and triplet network structures to derive semantically meaningful sentence embeddings, optimized for cosine similarity comparison~\cite{sbert}.

\item[HDBSCAN]
Hierarchical Density-Based Spatial Clustering of Applications with Noise; a density-based clustering algorithm that does not require specifying the number of clusters and can identify outliers (noise points)~\cite{hdbscan}.

% =============================================================================
% EXPERIMENTAL DESIGN TERMS
% =============================================================================

\item[Ablation Study]
An experimental procedure where individual components (e.g., \texttt{get\_service\_dependency} tool, fast filter, runbook) are removed from the full system to measure their individual contribution to performance metrics.

\item[Cross-System Generalization]
The ability of an aggregation method trained or tuned on one microservice system (e.g., Online Boutique) to maintain performance on unseen systems (e.g., Train Ticket) without retraining.

\item[Statistical Significance Test]
Paired Wilcoxon signed-rank test (non-parametric) used to compare metric distributions across the four arms on the same failure cases, with significance level $\alpha = 0.05$.

\item[Metric Gaming]
The phenomenon where a method optimizes a single metric (e.g., ARR) at the expense of operational utility (e.g., grouping all alerts into one group achieves ARR $\approx$ 100\% but RCPR = 0\%). Mitigated by the multi-metric framework (ARR, RCPR, F1, Purity).

% =============================================================================
% INDUSTRY STANDARDS & REFERENCES
% =============================================================================

\item[Google SRE Principles]
Foundational Site Reliability Engineering practices documented in \textit{Site Reliability Engineering} (Ch.~6, 10, 11) and \textit{The Site Reliability Workbook} (Ch.~5), emphasizing actionable alerts, maximum 2 incidents per 12-hour shift, and alerting on SLOs rather than raw metrics~\cite{google-sre-book, google-sre-workbook}.

\item[Observability Pillars]
The three canonical telemetry types: \textit{Logs} (immutable event records), \textit{Metrics} (aggregated numerical measurements), and \textit{Traces} (distributed request flow records)~\cite{grafana2025}.

\end{description}

% =============================================================================
% LOCAL BIBLIOGRAPHY MACROS (if not using main document's .bib)
% Include these in your main .bib or define via IEEEabrv.bib/IEEEfull.bib
% =============================================================================
% @string{srebook = "Site Reliability Engineering: How Google Runs Production Systems"}
% @string{sreworkbook = "The Site Reliability Workbook: Practical Ways to Implement SRE"}
% @article{pagerduty2020, ...}
% @article{grafana2025, ...}
% @inproceedings{kuang2024cola, ...}
% @inproceedings{pham2025rcaeval, ...}
% @article{zheng2024lemma, ...}
% @inproceedings{yu2025alertguardian, ...}
% @manual{prometheus-alertmanager, ...}
% @article{sbert, ...}
% @article{drain3, ...}
% @article{hdbscan, ...}
% @article{wei2022cot, ...} in main IEEEtran document
% Compiles with: pdflatex main.tex (where main.tex includes this file)
% =============================================================================

% -----------------------------------------------------------------------------
% USAGE IN MAIN DOCUMENT:
% \section*{Glossary of Terms}
% \addcontentsline{toc}{section}{Glossary of Terms}
% % =============================================================================
% GLOSSARY OF TERMS — Alert Aggregation in Microservices
% Standalone file for \input{glossary} in main IEEEtran document
% Compiles with: pdflatex main.tex (where main.tex includes this file)
% =============================================================================

% -----------------------------------------------------------------------------
% USAGE IN MAIN DOCUMENT:
% \section*{Glossary of Terms}
% \addcontentsline{toc}{section}{Glossary of Terms}
% \input{glossary}
% -----------------------------------------------------------------------------

\begin{description}

% =============================================================================
% CORE PROBLEM TERMS
% =============================================================================

\item[Alert Fatigue]
A cognitive state experienced by on-call engineers when the volume of alerts exceeds human cognitive capacity, leading to desensitization, missed critical alerts, and delayed incident response~\cite{pagerduty2020, grafana2025}.

\item[Alert Storm]
A sudden surge of alerts triggered by a single root cause propagating through dependent microservices, resulting in hundreds to thousands of alerts within minutes~\cite{kuang2024cola}.

\item[Noisy Alert Warning]
A warning-level alert generated by the monitoring system that is technically correct but does not require meaningful remediation action from an on-call engineer, classified into four categories: \textit{Duplicate}, \textit{Transient}, \textit{Unactionable}, and \textit{Cascading} (defined below).

\item[Early Warning Phase]
The pre-failure period during which a microservice exhibits performance degradation (e.g., increased latency, resource saturation) but has not yet violated Service Level Objectives (SLOs) or caused user-facing outages.

\item[Actionable Alert]
An alert that satisfies the Google SRE principle: ``every alert sent to an on-call engineer must require immediate human intervention''~\cite{google-sre-ch10, google-sre-ch11}.

% =============================================================================
% FOUR CATEGORIES OF NOISY ALERT WARNINGS
% =============================================================================

\item[Duplicate Warning Alert]
Multiple alerts originating from the same service with identical or near-identical message templates within a short time window ($\Delta t \le 60$~seconds), representing repeated observations of the same underlying condition rather than distinct incidents.

\item[Transient Warning Alert]
A warning alert that occurs during normal system operation (i.e., before fault injection time $T_{\text{inject}}$ in RCAEval) and resolves spontaneously without human intervention, typically caused by temporary network fluctuations or brief resource spikes.

\item[Unactionable Warning Alert]
A technically accurate warning alert for which no Standard Operating Procedure (SOP) or runbook exists, leaving the on-call engineer without a defined remediation path; the alert consumes attention but yields no operational value.

\item[Cascading Warning Alert]
A warning alert emitted by a downstream service that is a symptom of an upstream root cause; remediating the alerting service does not resolve the underlying issue, which originates in a different service (the root cause service).

% =============================================================================
% METHODOLOGY: FOUR EXPERIMENTAL ARMS
% =============================================================================

\item[Arm 1: Rule-Based Aggregation (Baseline)]
An alert grouping approach that clusters alerts based on static label matching (e.g., service name, alert name, severity) and fixed time windows (\texttt{group\_wait}, \texttt{group\_interval}), as implemented in Prometheus Alertmanager~\cite{prometheus-alertmanager}. Lacks awareness of service dependencies and semantic content.

\item[Arm 2: Semantic Similarity-Based Aggregation]
A method that parses raw logs into templates (via Drain3~\cite{drain3}), encodes templates into dense vector embeddings using Sentence-BERT~\cite{sbert}, and clusters alerts using density-based algorithms (HDBSCAN/DBSCAN) based on cosine similarity. Operates without topology knowledge.

\item[Arm 3: Temporal-Spatial Correlation Aggregation]
A topology-aware method that constructs a static service dependency graph (caller-callee adjacency matrix) from microservice architecture, then correlates alerts occurring within a sliding time window ($\Delta t \in \{30, 60, 120\}$~seconds) if a dependency path of length $k \le 2$ exists between the alerting services.

\item[Arm 4: LLM-Based Agentic Aggregation (Proposed)]
A hybrid two-stage approach: (1) a fast filter removes obvious duplicates and performs time-window pre-grouping; (2) an LLM agent (GPT-4o-mini / DeepSeek-V3) reasons over the remaining clusters using Chain-of-Thought prompting and three tools: \texttt{get\_related\_alerts}, \texttt{get\_service\_dependency}, and optionally \texttt{get\_runbook}. Limited to a maximum of two tool calls per grouping decision.

% =============================================================================
% TECHNICAL CONCEPTS & ARCHITECTURE
% =============================================================================

\item[Service Topology (Service Dependency Graph)]
A directed graph $G = (V, E)$ where vertices $V$ represent microservices and edges $E$ represent synchronous (HTTP/gRPC) or asynchronous (message queue) call relationships. Used to distinguish root cause services from downstream symptom services.

\item[Semantic Similarity]
A measure of meaning-based resemblance between two log messages, computed as the cosine similarity between their dense vector embeddings (typically 384-dimensional from \texttt{sentence-transformers/all-MiniLM-L6-v2}). Captures paraphrases and synonyms that lexical matching misses.

\item[Runbook Reasoning]
The process by which an LLM agent consults a synthetic runbook (a structured troubleshooting guide for a specific fault type) to determine whether an alert is actionable, requires escalation, or can be safely aggregated as noise.

\item[Tool-Calling Agent]
An LLM agent augmented with a predefined set of functions (tools) it can invoke during reasoning: \texttt{get\_related\_alerts(service, time\_window)}, \texttt{get\_service\_dependency(service)}, and \texttt{get\_runbook(fault\_type)}. Tool outputs are fed back into the agent's context for subsequent reasoning steps.

\item[Chain-of-Thought (CoT) Prompting]
A prompting technique that instructs the LLM to generate intermediate reasoning steps before producing a final answer, improving accuracy on multi-step causal inference tasks~\cite{wei2022cot}.

\item[Fast Filter (Deduplication + Time-Window Pre-Grouping)]
A lightweight deterministic preprocessing stage that removes exact duplicate alerts and performs initial grouping by service and time window, reducing the alert volume passed to the LLM agent by an estimated 70--80\%.

\item[Synthetic Runbook]
A researcher-authored troubleshooting document for each fault type in the benchmark dataset (e.g., CPU stress, memory leak, network latency), containing symptoms, diagnosis steps, and remediation actions. Not a production SOP; used solely for experimental evaluation.

% =============================================================================
% EVALUATION METRICS
% =============================================================================

\item[Alert Reduction Ratio (ARR)]
The primary metric quantifying alert volume reduction:
\begin{equation}
\text{ARR} = \left(1 - \frac{N_{\text{groups}}}{N_{\text{raw\_alerts}}}\right) \times 100\%
\end{equation}
where $N_{\text{raw\_alerts}}$ is the number of warning alerts before aggregation and $N_{\text{groups}}$ is the number of aggregated groups presented to the engineer. Target: $\text{ARR} \ge 60\%$.

\item[Root Cause Preservation Rate (RCPR)]
The critical safety metric ensuring the root cause evidence is not lost during aggregation. It is defined at two granularities. At service granularity, which is the primary metric and spans every case that exposes logs ($M = 359$):
\begin{equation}
\text{RCPR}_{\text{svc}} = \frac{1}{M} \sum_{i=1}^{M} \mathbb{I}(s_{\text{root}}^{(i)} \in \text{PrimaryGroup}_i) \times 100\%
\end{equation}
At alert granularity, a secondary metric restricted to the $M'$ cases whose ground-truth \texttt{root\_cause.txt} indicator line is itself a warning-tier log ($M' = 4$):
\begin{equation}
\text{RCPR}_{\text{alert}} = \frac{1}{M'} \sum_{i=1}^{M'} \mathbb{I}(\text{RootCauseAlert}_i \in \text{PrimaryGroup}_i) \times 100\%
\end{equation}
where $\mathbb{I}$ is the indicator function. Target: $\text{RCPR} \ge 95\%$ (mandatory).

\item[Pairwise Grouping Precision ($P_p$)]
The fraction of alert pairs placed in the same group that truly share the same root cause:
\begin{equation}
P_p = \frac{|\{(a_x, a_y) : \text{same group} \land \text{same root cause}\}|}{|\{(a_x, a_y) : \text{same group}\}|}
\end{equation}

\item[Pairwise Grouping Recall ($R_p$)]
The fraction of alert pairs sharing the same root cause that are successfully grouped together:
\begin{equation}
R_p = \frac{|\{(a_x, a_y) : \text{same group} \land \text{same root cause}\}|}{|\{(a_x, a_y) : \text{same root cause}\}|}
\end{equation}

\item[Pairwise Grouping F1-Score ($F1_p$)]
The harmonic mean of pairwise precision and recall:
\begin{equation}
F1_p = 2 \times \frac{P_p \times R_p}{P_p + R_p}
\end{equation}
Target: $F1_p \ge 0.85$.

\item[Group Purity]
The average homogeneity of root causes within each generated group:
\begin{equation}
\text{Purity} = \frac{1}{N_{\text{groups}}} \sum_{g=1}^{N_{\text{groups}}} \frac{\max_{c} |\{a \in g : \text{root\_cause}(a) = c\}|}{|g|}
\end{equation}
Target: $\text{Purity} \ge 0.80$.

\item[Processing Latency]
The wall-clock time required to process 100 warning alerts through the aggregation pipeline, measured in milliseconds. For Arm 4, includes LLM inference time and token overhead.

% =============================================================================
% DATASET & BENCHMARK TERMS
% =============================================================================

\item[RCAEval]
A benchmark for root cause analysis in microservice systems comprising 735 failure cases across three systems (Online Boutique, Sock Shop, Train Ticket) with logs, metrics, traces, and root cause labels~\cite{pham2025rcaeval}. Licensed under MIT. Direct inspection shows that only \textbf{359 of the 735 cases ship raw log telemetry} (suite RE1 is metric-only) and that the log schema contains \texttt{timestamp}, \texttt{container\_name}, and \texttt{message} but \textbf{no severity column}.

\item[Severity Tier Assignment]
The three-tier procedure by which this work assigns the severity tier $\ell_i$, since RCAEval supplies none: \emph{Tier 1 (explicit)} parses a level token from the message text (Sock Shop, Train Ticket); \emph{Tier 2 (frequency-inferred)} flags a template as anomalous when its fault-window frequency exceeds its normal-window frequency by a ratio $R$ (Online Boutique); \emph{Tier 3 (ground-truth marker)} tags the \texttt{root\_cause.txt} indicator line as \texttt{ROOT\_CAUSE} for the eight annotated cases.

\item[LEMMA-RCA]
A large-scale multi-modal dataset for root cause analysis containing hundreds of microservice entities and four fault types (DDoS, storage failure, resource contention, noisy neighbor)~\cite{zheng2024lemma}. Licensed under CC-BY-NC-4.0 (non-commercial). Its log pipeline aggregates raw log lines into three-dimensional time series using golden-signal keywords that cover only \texttt{error}, \texttt{exception}, and \texttt{critical} --- notably \textbf{not} \texttt{warning} --- so it offers no independent warning tier.

\item[Fault Injection]
The controlled introduction of a failure (e.g., CPU stress, network latency, packet loss) into a microservice system at a known timestamp $T_{\text{inject}}$ to generate labeled failure cases for benchmarking.

\item[Normal Window]
The interval $[T_{\text{start}}, T_{\text{inject}})$ preceding fault injection, during which the system operates normally; alerts occurring here are candidates for \textit{Transient} noise. Note that the dataset fields \texttt{normal\_timesteps} and \texttt{faulty\_timesteps} are integer counts of sample points, not timestamp lists, and therefore define these windows only indirectly.

\item[Fault Window]
The interval $[T_{\text{inject}}, T_{\text{end}}]$ following fault injection, during which the system exhibits degradation; alerts here are labeled as \textit{Signal} (if from the root cause service) or \textit{Cascading} (if from a service reachable from the root cause service in the call graph).

\item[Ground Truth Root Cause Label]
The authoritative annotation of which service is the root cause of a failure case, provided by the RCAEval benchmark. Used to compute RCPR and pairwise metrics.

% =============================================================================
% BASELINE & INFRASTRUCTURE TERMS
% =============================================================================

\item[Prometheus Alertmanager]
The de facto standard alert routing and grouping component in the Cloud Native Computing Foundation (CNCF) ecosystem. Groups alerts by label matchers and time windows (\texttt{group\_wait}, \texttt{group\_interval}, \texttt{group\_by})~\cite{prometheus-alertmanager}.

\item[Group Wait (\texttt{group\_wait})]
The initial delay after the first alert in a group before sending a notification, allowing additional alerts to arrive and be grouped. Default: 30~seconds.

\item[Group Interval (\texttt{group\_interval})]
The minimum interval between subsequent notifications for the same group. Default: 5~minutes.

\item[Drain3]
An online log parsing algorithm that extracts log templates from raw log streams by clustering log messages based on token similarity, producing parameterized templates with variable placeholders~\cite{drain3}.

\item[Sentence-BERT (S-BERT)]
A modification of BERT that uses siamese and triplet network structures to derive semantically meaningful sentence embeddings, optimized for cosine similarity comparison~\cite{sbert}.

\item[HDBSCAN]
Hierarchical Density-Based Spatial Clustering of Applications with Noise; a density-based clustering algorithm that does not require specifying the number of clusters and can identify outliers (noise points)~\cite{hdbscan}.

% =============================================================================
% EXPERIMENTAL DESIGN TERMS
% =============================================================================

\item[Ablation Study]
An experimental procedure where individual components (e.g., \texttt{get\_service\_dependency} tool, fast filter, runbook) are removed from the full system to measure their individual contribution to performance metrics.

\item[Cross-System Generalization]
The ability of an aggregation method trained or tuned on one microservice system (e.g., Online Boutique) to maintain performance on unseen systems (e.g., Train Ticket) without retraining.

\item[Statistical Significance Test]
Paired Wilcoxon signed-rank test (non-parametric) used to compare metric distributions across the four arms on the same failure cases, with significance level $\alpha = 0.05$.

\item[Metric Gaming]
The phenomenon where a method optimizes a single metric (e.g., ARR) at the expense of operational utility (e.g., grouping all alerts into one group achieves ARR $\approx$ 100\% but RCPR = 0\%). Mitigated by the multi-metric framework (ARR, RCPR, F1, Purity).

% =============================================================================
% INDUSTRY STANDARDS & REFERENCES
% =============================================================================

\item[Google SRE Principles]
Foundational Site Reliability Engineering practices documented in \textit{Site Reliability Engineering} (Ch.~6, 10, 11) and \textit{The Site Reliability Workbook} (Ch.~5), emphasizing actionable alerts, maximum 2 incidents per 12-hour shift, and alerting on SLOs rather than raw metrics~\cite{google-sre-book, google-sre-workbook}.

\item[Observability Pillars]
The three canonical telemetry types: \textit{Logs} (immutable event records), \textit{Metrics} (aggregated numerical measurements), and \textit{Traces} (distributed request flow records)~\cite{grafana2025}.

\end{description}

% =============================================================================
% LOCAL BIBLIOGRAPHY MACROS (if not using main document's .bib)
% Include these in your main .bib or define via IEEEabrv.bib/IEEEfull.bib
% =============================================================================
% @string{srebook = "Site Reliability Engineering: How Google Runs Production Systems"}
% @string{sreworkbook = "The Site Reliability Workbook: Practical Ways to Implement SRE"}
% @article{pagerduty2020, ...}
% @article{grafana2025, ...}
% @inproceedings{kuang2024cola, ...}
% @inproceedings{pham2025rcaeval, ...}
% @article{zheng2024lemma, ...}
% @inproceedings{yu2025alertguardian, ...}
% @manual{prometheus-alertmanager, ...}
% @article{sbert, ...}
% @article{drain3, ...}
% @article{hdbscan, ...}
% @article{wei2022cot, ...}
% -----------------------------------------------------------------------------

\begin{description}

% =============================================================================
% CORE PROBLEM TERMS
% =============================================================================

\item[Alert Fatigue]
A cognitive state experienced by on-call engineers when the volume of alerts exceeds human cognitive capacity, leading to desensitization, missed critical alerts, and delayed incident response~\cite{pagerduty2020, grafana2025}.

\item[Alert Storm]
A sudden surge of alerts triggered by a single root cause propagating through dependent microservices, resulting in hundreds to thousands of alerts within minutes~\cite{kuang2024cola}.

\item[Noisy Alert Warning]
A warning-level alert generated by the monitoring system that is technically correct but does not require meaningful remediation action from an on-call engineer, classified into four categories: \textit{Duplicate}, \textit{Transient}, \textit{Unactionable}, and \textit{Cascading} (defined below).

\item[Early Warning Phase]
The pre-failure period during which a microservice exhibits performance degradation (e.g., increased latency, resource saturation) but has not yet violated Service Level Objectives (SLOs) or caused user-facing outages.

\item[Actionable Alert]
An alert that satisfies the Google SRE principle: ``every alert sent to an on-call engineer must require immediate human intervention''~\cite{google-sre-ch10, google-sre-ch11}.

% =============================================================================
% FOUR CATEGORIES OF NOISY ALERT WARNINGS
% =============================================================================

\item[Duplicate Warning Alert]
Multiple alerts originating from the same service with identical or near-identical message templates within a short time window ($\Delta t \le 60$~seconds), representing repeated observations of the same underlying condition rather than distinct incidents.

\item[Transient Warning Alert]
A warning alert that occurs during normal system operation (i.e., before fault injection time $T_{\text{inject}}$ in RCAEval) and resolves spontaneously without human intervention, typically caused by temporary network fluctuations or brief resource spikes.

\item[Unactionable Warning Alert]
A technically accurate warning alert for which no Standard Operating Procedure (SOP) or runbook exists, leaving the on-call engineer without a defined remediation path; the alert consumes attention but yields no operational value.

\item[Cascading Warning Alert]
A warning alert emitted by a downstream service that is a symptom of an upstream root cause; remediating the alerting service does not resolve the underlying issue, which originates in a different service (the root cause service).

% =============================================================================
% METHODOLOGY: FOUR EXPERIMENTAL ARMS
% =============================================================================

\item[Arm 1: Rule-Based Aggregation (Baseline)]
An alert grouping approach that clusters alerts based on static label matching (e.g., service name, alert name, severity) and fixed time windows (\texttt{group\_wait}, \texttt{group\_interval}), as implemented in Prometheus Alertmanager~\cite{prometheus-alertmanager}. Lacks awareness of service dependencies and semantic content.

\item[Arm 2: Semantic Similarity-Based Aggregation]
A method that parses raw logs into templates (via Drain3~\cite{drain3}), encodes templates into dense vector embeddings using Sentence-BERT~\cite{sbert}, and clusters alerts using density-based algorithms (HDBSCAN/DBSCAN) based on cosine similarity. Operates without topology knowledge.

\item[Arm 3: Temporal-Spatial Correlation Aggregation]
A topology-aware method that constructs a static service dependency graph (caller-callee adjacency matrix) from microservice architecture, then correlates alerts occurring within a sliding time window ($\Delta t \in \{30, 60, 120\}$~seconds) if a dependency path of length $k \le 2$ exists between the alerting services.

\item[Arm 4: LLM-Based Agentic Aggregation (Proposed)]
A hybrid two-stage approach: (1) a fast filter removes obvious duplicates and performs time-window pre-grouping; (2) an LLM agent (GPT-4o-mini / DeepSeek-V3) reasons over the remaining clusters using Chain-of-Thought prompting and three tools: \texttt{get\_related\_alerts}, \texttt{get\_service\_dependency}, and optionally \texttt{get\_runbook}. Limited to a maximum of two tool calls per grouping decision.

% =============================================================================
% TECHNICAL CONCEPTS & ARCHITECTURE
% =============================================================================

\item[Service Topology (Service Dependency Graph)]
A directed graph $G = (V, E)$ where vertices $V$ represent microservices and edges $E$ represent synchronous (HTTP/gRPC) or asynchronous (message queue) call relationships. Used to distinguish root cause services from downstream symptom services.

\item[Semantic Similarity]
A measure of meaning-based resemblance between two log messages, computed as the cosine similarity between their dense vector embeddings (typically 384-dimensional from \texttt{sentence-transformers/all-MiniLM-L6-v2}). Captures paraphrases and synonyms that lexical matching misses.

\item[Runbook Reasoning]
The process by which an LLM agent consults a synthetic runbook (a structured troubleshooting guide for a specific fault type) to determine whether an alert is actionable, requires escalation, or can be safely aggregated as noise.

\item[Tool-Calling Agent]
An LLM agent augmented with a predefined set of functions (tools) it can invoke during reasoning: \texttt{get\_related\_alerts(service, time\_window)}, \texttt{get\_service\_dependency(service)}, and \texttt{get\_runbook(fault\_type)}. Tool outputs are fed back into the agent's context for subsequent reasoning steps.

\item[Chain-of-Thought (CoT) Prompting]
A prompting technique that instructs the LLM to generate intermediate reasoning steps before producing a final answer, improving accuracy on multi-step causal inference tasks~\cite{wei2022cot}.

\item[Fast Filter (Deduplication + Time-Window Pre-Grouping)]
A lightweight deterministic preprocessing stage that removes exact duplicate alerts and performs initial grouping by service and time window, reducing the alert volume passed to the LLM agent by an estimated 70--80\%.

\item[Synthetic Runbook]
A researcher-authored troubleshooting document for each fault type in the benchmark dataset (e.g., CPU stress, memory leak, network latency), containing symptoms, diagnosis steps, and remediation actions. Not a production SOP; used solely for experimental evaluation.

% =============================================================================
% EVALUATION METRICS
% =============================================================================

\item[Alert Reduction Ratio (ARR)]
The primary metric quantifying alert volume reduction:
\begin{equation}
\text{ARR} = \left(1 - \frac{N_{\text{groups}}}{N_{\text{raw\_alerts}}}\right) \times 100\%
\end{equation}
where $N_{\text{raw\_alerts}}$ is the number of warning alerts before aggregation and $N_{\text{groups}}$ is the number of aggregated groups presented to the engineer. Target: $\text{ARR} \ge 60\%$.

\item[Root Cause Preservation Rate (RCPR)]
The critical safety metric ensuring the root cause evidence is not lost during aggregation. It is defined at two granularities. At service granularity, which is the primary metric and spans every case that exposes logs ($M = 359$):
\begin{equation}
\text{RCPR}_{\text{svc}} = \frac{1}{M} \sum_{i=1}^{M} \mathbb{I}(s_{\text{root}}^{(i)} \in \text{PrimaryGroup}_i) \times 100\%
\end{equation}
At alert granularity, a secondary metric restricted to the $M'$ cases whose ground-truth \texttt{root\_cause.txt} indicator line is itself a warning-tier log ($M' = 4$):
\begin{equation}
\text{RCPR}_{\text{alert}} = \frac{1}{M'} \sum_{i=1}^{M'} \mathbb{I}(\text{RootCauseAlert}_i \in \text{PrimaryGroup}_i) \times 100\%
\end{equation}
where $\mathbb{I}$ is the indicator function. Target: $\text{RCPR} \ge 95\%$ (mandatory).

\item[Pairwise Grouping Precision ($P_p$)]
The fraction of alert pairs placed in the same group that truly share the same root cause:
\begin{equation}
P_p = \frac{|\{(a_x, a_y) : \text{same group} \land \text{same root cause}\}|}{|\{(a_x, a_y) : \text{same group}\}|}
\end{equation}

\item[Pairwise Grouping Recall ($R_p$)]
The fraction of alert pairs sharing the same root cause that are successfully grouped together:
\begin{equation}
R_p = \frac{|\{(a_x, a_y) : \text{same group} \land \text{same root cause}\}|}{|\{(a_x, a_y) : \text{same root cause}\}|}
\end{equation}

\item[Pairwise Grouping F1-Score ($F1_p$)]
The harmonic mean of pairwise precision and recall:
\begin{equation}
F1_p = 2 \times \frac{P_p \times R_p}{P_p + R_p}
\end{equation}
Target: $F1_p \ge 0.85$.

\item[Group Purity]
The average homogeneity of root causes within each generated group:
\begin{equation}
\text{Purity} = \frac{1}{N_{\text{groups}}} \sum_{g=1}^{N_{\text{groups}}} \frac{\max_{c} |\{a \in g : \text{root\_cause}(a) = c\}|}{|g|}
\end{equation}
Target: $\text{Purity} \ge 0.80$.

\item[Processing Latency]
The wall-clock time required to process 100 warning alerts through the aggregation pipeline, measured in milliseconds. For Arm 4, includes LLM inference time and token overhead.

% =============================================================================
% DATASET & BENCHMARK TERMS
% =============================================================================

\item[RCAEval]
A benchmark for root cause analysis in microservice systems comprising 735 failure cases across three systems (Online Boutique, Sock Shop, Train Ticket) with logs, metrics, traces, and root cause labels~\cite{pham2025rcaeval}. Licensed under MIT. Direct inspection shows that only \textbf{359 of the 735 cases ship raw log telemetry} (suite RE1 is metric-only) and that the log schema contains \texttt{timestamp}, \texttt{container\_name}, and \texttt{message} but \textbf{no severity column}.

\item[Severity Tier Assignment]
The three-tier procedure by which this work assigns the severity tier $\ell_i$, since RCAEval supplies none: \emph{Tier 1 (explicit)} parses a level token from the message text (Sock Shop, Train Ticket); \emph{Tier 2 (frequency-inferred)} flags a template as anomalous when its fault-window frequency exceeds its normal-window frequency by a ratio $R$ (Online Boutique); \emph{Tier 3 (ground-truth marker)} tags the \texttt{root\_cause.txt} indicator line as \texttt{ROOT\_CAUSE} for the eight annotated cases.

\item[LEMMA-RCA]
A large-scale multi-modal dataset for root cause analysis containing hundreds of microservice entities and four fault types (DDoS, storage failure, resource contention, noisy neighbor)~\cite{zheng2024lemma}. Licensed under CC-BY-NC-4.0 (non-commercial). Its log pipeline aggregates raw log lines into three-dimensional time series using golden-signal keywords that cover only \texttt{error}, \texttt{exception}, and \texttt{critical} --- notably \textbf{not} \texttt{warning} --- so it offers no independent warning tier.

\item[Fault Injection]
The controlled introduction of a failure (e.g., CPU stress, network latency, packet loss) into a microservice system at a known timestamp $T_{\text{inject}}$ to generate labeled failure cases for benchmarking.

\item[Normal Window]
The interval $[T_{\text{start}}, T_{\text{inject}})$ preceding fault injection, during which the system operates normally; alerts occurring here are candidates for \textit{Transient} noise. Note that the dataset fields \texttt{normal\_timesteps} and \texttt{faulty\_timesteps} are integer counts of sample points, not timestamp lists, and therefore define these windows only indirectly.

\item[Fault Window]
The interval $[T_{\text{inject}}, T_{\text{end}}]$ following fault injection, during which the system exhibits degradation; alerts here are labeled as \textit{Signal} (if from the root cause service) or \textit{Cascading} (if from a service reachable from the root cause service in the call graph).

\item[Ground Truth Root Cause Label]
The authoritative annotation of which service is the root cause of a failure case, provided by the RCAEval benchmark. Used to compute RCPR and pairwise metrics.

% =============================================================================
% BASELINE & INFRASTRUCTURE TERMS
% =============================================================================

\item[Prometheus Alertmanager]
The de facto standard alert routing and grouping component in the Cloud Native Computing Foundation (CNCF) ecosystem. Groups alerts by label matchers and time windows (\texttt{group\_wait}, \texttt{group\_interval}, \texttt{group\_by})~\cite{prometheus-alertmanager}.

\item[Group Wait (\texttt{group\_wait})]
The initial delay after the first alert in a group before sending a notification, allowing additional alerts to arrive and be grouped. Default: 30~seconds.

\item[Group Interval (\texttt{group\_interval})]
The minimum interval between subsequent notifications for the same group. Default: 5~minutes.

\item[Drain3]
An online log parsing algorithm that extracts log templates from raw log streams by clustering log messages based on token similarity, producing parameterized templates with variable placeholders~\cite{drain3}.

\item[Sentence-BERT (S-BERT)]
A modification of BERT that uses siamese and triplet network structures to derive semantically meaningful sentence embeddings, optimized for cosine similarity comparison~\cite{sbert}.

\item[HDBSCAN]
Hierarchical Density-Based Spatial Clustering of Applications with Noise; a density-based clustering algorithm that does not require specifying the number of clusters and can identify outliers (noise points)~\cite{hdbscan}.

% =============================================================================
% EXPERIMENTAL DESIGN TERMS
% =============================================================================

\item[Ablation Study]
An experimental procedure where individual components (e.g., \texttt{get\_service\_dependency} tool, fast filter, runbook) are removed from the full system to measure their individual contribution to performance metrics.

\item[Cross-System Generalization]
The ability of an aggregation method trained or tuned on one microservice system (e.g., Online Boutique) to maintain performance on unseen systems (e.g., Train Ticket) without retraining.

\item[Statistical Significance Test]
Paired Wilcoxon signed-rank test (non-parametric) used to compare metric distributions across the four arms on the same failure cases, with significance level $\alpha = 0.05$.

\item[Metric Gaming]
The phenomenon where a method optimizes a single metric (e.g., ARR) at the expense of operational utility (e.g., grouping all alerts into one group achieves ARR $\approx$ 100\% but RCPR = 0\%). Mitigated by the multi-metric framework (ARR, RCPR, F1, Purity).

% =============================================================================
% INDUSTRY STANDARDS & REFERENCES
% =============================================================================

\item[Google SRE Principles]
Foundational Site Reliability Engineering practices documented in \textit{Site Reliability Engineering} (Ch.~6, 10, 11) and \textit{The Site Reliability Workbook} (Ch.~5), emphasizing actionable alerts, maximum 2 incidents per 12-hour shift, and alerting on SLOs rather than raw metrics~\cite{google-sre-book, google-sre-workbook}.

\item[Observability Pillars]
The three canonical telemetry types: \textit{Logs} (immutable event records), \textit{Metrics} (aggregated numerical measurements), and \textit{Traces} (distributed request flow records)~\cite{grafana2025}.

\end{description}

% =============================================================================
% LOCAL BIBLIOGRAPHY MACROS (if not using main document's .bib)
% Include these in your main .bib or define via IEEEabrv.bib/IEEEfull.bib
% =============================================================================
% @string{srebook = "Site Reliability Engineering: How Google Runs Production Systems"}
% @string{sreworkbook = "The Site Reliability Workbook: Practical Ways to Implement SRE"}
% @article{pagerduty2020, ...}
% @article{grafana2025, ...}
% @inproceedings{kuang2024cola, ...}
% @inproceedings{pham2025rcaeval, ...}
% @article{zheng2024lemma, ...}
% @inproceedings{yu2025alertguardian, ...}
% @manual{prometheus-alertmanager, ...}
% @article{sbert, ...}
% @article{drain3, ...}
% @article{hdbscan, ...}
% @article{wei2022cot, ...} in main IEEEtran document
% Compiles with: pdflatex main.tex (where main.tex includes this file)
% =============================================================================

% -----------------------------------------------------------------------------
% USAGE IN MAIN DOCUMENT:
% \section*{Glossary of Terms}
% \addcontentsline{toc}{section}{Glossary of Terms}
% % =============================================================================
% GLOSSARY OF TERMS — Alert Aggregation in Microservices
% Standalone file for % =============================================================================
% GLOSSARY OF TERMS — Alert Aggregation in Microservices
% Standalone file for \input{glossary} in main IEEEtran document
% Compiles with: pdflatex main.tex (where main.tex includes this file)
% =============================================================================

% -----------------------------------------------------------------------------
% USAGE IN MAIN DOCUMENT:
% \section*{Glossary of Terms}
% \addcontentsline{toc}{section}{Glossary of Terms}
% \input{glossary}
% -----------------------------------------------------------------------------

\begin{description}

% =============================================================================
% CORE PROBLEM TERMS
% =============================================================================

\item[Alert Fatigue]
A cognitive state experienced by on-call engineers when the volume of alerts exceeds human cognitive capacity, leading to desensitization, missed critical alerts, and delayed incident response~\cite{pagerduty2020, grafana2025}.

\item[Alert Storm]
A sudden surge of alerts triggered by a single root cause propagating through dependent microservices, resulting in hundreds to thousands of alerts within minutes~\cite{kuang2024cola}.

\item[Noisy Alert Warning]
A warning-level alert generated by the monitoring system that is technically correct but does not require meaningful remediation action from an on-call engineer, classified into four categories: \textit{Duplicate}, \textit{Transient}, \textit{Unactionable}, and \textit{Cascading} (defined below).

\item[Early Warning Phase]
The pre-failure period during which a microservice exhibits performance degradation (e.g., increased latency, resource saturation) but has not yet violated Service Level Objectives (SLOs) or caused user-facing outages.

\item[Actionable Alert]
An alert that satisfies the Google SRE principle: ``every alert sent to an on-call engineer must require immediate human intervention''~\cite{google-sre-ch10, google-sre-ch11}.

% =============================================================================
% FOUR CATEGORIES OF NOISY ALERT WARNINGS
% =============================================================================

\item[Duplicate Warning Alert]
Multiple alerts originating from the same service with identical or near-identical message templates within a short time window ($\Delta t \le 60$~seconds), representing repeated observations of the same underlying condition rather than distinct incidents.

\item[Transient Warning Alert]
A warning alert that occurs during normal system operation (i.e., before fault injection time $T_{\text{inject}}$ in RCAEval) and resolves spontaneously without human intervention, typically caused by temporary network fluctuations or brief resource spikes.

\item[Unactionable Warning Alert]
A technically accurate warning alert for which no Standard Operating Procedure (SOP) or runbook exists, leaving the on-call engineer without a defined remediation path; the alert consumes attention but yields no operational value.

\item[Cascading Warning Alert]
A warning alert emitted by a downstream service that is a symptom of an upstream root cause; remediating the alerting service does not resolve the underlying issue, which originates in a different service (the root cause service).

% =============================================================================
% METHODOLOGY: FOUR EXPERIMENTAL ARMS
% =============================================================================

\item[Arm 1: Rule-Based Aggregation (Baseline)]
An alert grouping approach that clusters alerts based on static label matching (e.g., service name, alert name, severity) and fixed time windows (\texttt{group\_wait}, \texttt{group\_interval}), as implemented in Prometheus Alertmanager~\cite{prometheus-alertmanager}. Lacks awareness of service dependencies and semantic content.

\item[Arm 2: Semantic Similarity-Based Aggregation]
A method that parses raw logs into templates (via Drain3~\cite{drain3}), encodes templates into dense vector embeddings using Sentence-BERT~\cite{sbert}, and clusters alerts using density-based algorithms (HDBSCAN/DBSCAN) based on cosine similarity. Operates without topology knowledge.

\item[Arm 3: Temporal-Spatial Correlation Aggregation]
A topology-aware method that constructs a static service dependency graph (caller-callee adjacency matrix) from microservice architecture, then correlates alerts occurring within a sliding time window ($\Delta t \in \{30, 60, 120\}$~seconds) if a dependency path of length $k \le 2$ exists between the alerting services.

\item[Arm 4: LLM-Based Agentic Aggregation (Proposed)]
A hybrid two-stage approach: (1) a fast filter removes obvious duplicates and performs time-window pre-grouping; (2) an LLM agent (GPT-4o-mini / DeepSeek-V3) reasons over the remaining clusters using Chain-of-Thought prompting and three tools: \texttt{get\_related\_alerts}, \texttt{get\_service\_dependency}, and optionally \texttt{get\_runbook}. Limited to a maximum of two tool calls per grouping decision.

% =============================================================================
% TECHNICAL CONCEPTS & ARCHITECTURE
% =============================================================================

\item[Service Topology (Service Dependency Graph)]
A directed graph $G = (V, E)$ where vertices $V$ represent microservices and edges $E$ represent synchronous (HTTP/gRPC) or asynchronous (message queue) call relationships. Used to distinguish root cause services from downstream symptom services.

\item[Semantic Similarity]
A measure of meaning-based resemblance between two log messages, computed as the cosine similarity between their dense vector embeddings (typically 384-dimensional from \texttt{sentence-transformers/all-MiniLM-L6-v2}). Captures paraphrases and synonyms that lexical matching misses.

\item[Runbook Reasoning]
The process by which an LLM agent consults a synthetic runbook (a structured troubleshooting guide for a specific fault type) to determine whether an alert is actionable, requires escalation, or can be safely aggregated as noise.

\item[Tool-Calling Agent]
An LLM agent augmented with a predefined set of functions (tools) it can invoke during reasoning: \texttt{get\_related\_alerts(service, time\_window)}, \texttt{get\_service\_dependency(service)}, and \texttt{get\_runbook(fault\_type)}. Tool outputs are fed back into the agent's context for subsequent reasoning steps.

\item[Chain-of-Thought (CoT) Prompting]
A prompting technique that instructs the LLM to generate intermediate reasoning steps before producing a final answer, improving accuracy on multi-step causal inference tasks~\cite{wei2022cot}.

\item[Fast Filter (Deduplication + Time-Window Pre-Grouping)]
A lightweight deterministic preprocessing stage that removes exact duplicate alerts and performs initial grouping by service and time window, reducing the alert volume passed to the LLM agent by an estimated 70--80\%.

\item[Synthetic Runbook]
A researcher-authored troubleshooting document for each fault type in the benchmark dataset (e.g., CPU stress, memory leak, network latency), containing symptoms, diagnosis steps, and remediation actions. Not a production SOP; used solely for experimental evaluation.

% =============================================================================
% EVALUATION METRICS
% =============================================================================

\item[Alert Reduction Ratio (ARR)]
The primary metric quantifying alert volume reduction:
\begin{equation}
\text{ARR} = \left(1 - \frac{N_{\text{groups}}}{N_{\text{raw\_alerts}}}\right) \times 100\%
\end{equation}
where $N_{\text{raw\_alerts}}$ is the number of warning alerts before aggregation and $N_{\text{groups}}$ is the number of aggregated groups presented to the engineer. Target: $\text{ARR} \ge 60\%$.

\item[Root Cause Preservation Rate (RCPR)]
The critical safety metric ensuring the root cause evidence is not lost during aggregation. It is defined at two granularities. At service granularity, which is the primary metric and spans every case that exposes logs ($M = 359$):
\begin{equation}
\text{RCPR}_{\text{svc}} = \frac{1}{M} \sum_{i=1}^{M} \mathbb{I}(s_{\text{root}}^{(i)} \in \text{PrimaryGroup}_i) \times 100\%
\end{equation}
At alert granularity, a secondary metric restricted to the $M'$ cases whose ground-truth \texttt{root\_cause.txt} indicator line is itself a warning-tier log ($M' = 4$):
\begin{equation}
\text{RCPR}_{\text{alert}} = \frac{1}{M'} \sum_{i=1}^{M'} \mathbb{I}(\text{RootCauseAlert}_i \in \text{PrimaryGroup}_i) \times 100\%
\end{equation}
where $\mathbb{I}$ is the indicator function. Target: $\text{RCPR} \ge 95\%$ (mandatory).

\item[Pairwise Grouping Precision ($P_p$)]
The fraction of alert pairs placed in the same group that truly share the same root cause:
\begin{equation}
P_p = \frac{|\{(a_x, a_y) : \text{same group} \land \text{same root cause}\}|}{|\{(a_x, a_y) : \text{same group}\}|}
\end{equation}

\item[Pairwise Grouping Recall ($R_p$)]
The fraction of alert pairs sharing the same root cause that are successfully grouped together:
\begin{equation}
R_p = \frac{|\{(a_x, a_y) : \text{same group} \land \text{same root cause}\}|}{|\{(a_x, a_y) : \text{same root cause}\}|}
\end{equation}

\item[Pairwise Grouping F1-Score ($F1_p$)]
The harmonic mean of pairwise precision and recall:
\begin{equation}
F1_p = 2 \times \frac{P_p \times R_p}{P_p + R_p}
\end{equation}
Target: $F1_p \ge 0.85$.

\item[Group Purity]
The average homogeneity of root causes within each generated group:
\begin{equation}
\text{Purity} = \frac{1}{N_{\text{groups}}} \sum_{g=1}^{N_{\text{groups}}} \frac{\max_{c} |\{a \in g : \text{root\_cause}(a) = c\}|}{|g|}
\end{equation}
Target: $\text{Purity} \ge 0.80$.

\item[Processing Latency]
The wall-clock time required to process 100 warning alerts through the aggregation pipeline, measured in milliseconds. For Arm 4, includes LLM inference time and token overhead.

% =============================================================================
% DATASET & BENCHMARK TERMS
% =============================================================================

\item[RCAEval]
A benchmark for root cause analysis in microservice systems comprising 735 failure cases across three systems (Online Boutique, Sock Shop, Train Ticket) with logs, metrics, traces, and root cause labels~\cite{pham2025rcaeval}. Licensed under MIT. Direct inspection shows that only \textbf{359 of the 735 cases ship raw log telemetry} (suite RE1 is metric-only) and that the log schema contains \texttt{timestamp}, \texttt{container\_name}, and \texttt{message} but \textbf{no severity column}.

\item[Severity Tier Assignment]
The three-tier procedure by which this work assigns the severity tier $\ell_i$, since RCAEval supplies none: \emph{Tier 1 (explicit)} parses a level token from the message text (Sock Shop, Train Ticket); \emph{Tier 2 (frequency-inferred)} flags a template as anomalous when its fault-window frequency exceeds its normal-window frequency by a ratio $R$ (Online Boutique); \emph{Tier 3 (ground-truth marker)} tags the \texttt{root\_cause.txt} indicator line as \texttt{ROOT\_CAUSE} for the eight annotated cases.

\item[LEMMA-RCA]
A large-scale multi-modal dataset for root cause analysis containing hundreds of microservice entities and four fault types (DDoS, storage failure, resource contention, noisy neighbor)~\cite{zheng2024lemma}. Licensed under CC-BY-NC-4.0 (non-commercial). Its log pipeline aggregates raw log lines into three-dimensional time series using golden-signal keywords that cover only \texttt{error}, \texttt{exception}, and \texttt{critical} --- notably \textbf{not} \texttt{warning} --- so it offers no independent warning tier.

\item[Fault Injection]
The controlled introduction of a failure (e.g., CPU stress, network latency, packet loss) into a microservice system at a known timestamp $T_{\text{inject}}$ to generate labeled failure cases for benchmarking.

\item[Normal Window]
The interval $[T_{\text{start}}, T_{\text{inject}})$ preceding fault injection, during which the system operates normally; alerts occurring here are candidates for \textit{Transient} noise. Note that the dataset fields \texttt{normal\_timesteps} and \texttt{faulty\_timesteps} are integer counts of sample points, not timestamp lists, and therefore define these windows only indirectly.

\item[Fault Window]
The interval $[T_{\text{inject}}, T_{\text{end}}]$ following fault injection, during which the system exhibits degradation; alerts here are labeled as \textit{Signal} (if from the root cause service) or \textit{Cascading} (if from a service reachable from the root cause service in the call graph).

\item[Ground Truth Root Cause Label]
The authoritative annotation of which service is the root cause of a failure case, provided by the RCAEval benchmark. Used to compute RCPR and pairwise metrics.

% =============================================================================
% BASELINE & INFRASTRUCTURE TERMS
% =============================================================================

\item[Prometheus Alertmanager]
The de facto standard alert routing and grouping component in the Cloud Native Computing Foundation (CNCF) ecosystem. Groups alerts by label matchers and time windows (\texttt{group\_wait}, \texttt{group\_interval}, \texttt{group\_by})~\cite{prometheus-alertmanager}.

\item[Group Wait (\texttt{group\_wait})]
The initial delay after the first alert in a group before sending a notification, allowing additional alerts to arrive and be grouped. Default: 30~seconds.

\item[Group Interval (\texttt{group\_interval})]
The minimum interval between subsequent notifications for the same group. Default: 5~minutes.

\item[Drain3]
An online log parsing algorithm that extracts log templates from raw log streams by clustering log messages based on token similarity, producing parameterized templates with variable placeholders~\cite{drain3}.

\item[Sentence-BERT (S-BERT)]
A modification of BERT that uses siamese and triplet network structures to derive semantically meaningful sentence embeddings, optimized for cosine similarity comparison~\cite{sbert}.

\item[HDBSCAN]
Hierarchical Density-Based Spatial Clustering of Applications with Noise; a density-based clustering algorithm that does not require specifying the number of clusters and can identify outliers (noise points)~\cite{hdbscan}.

% =============================================================================
% EXPERIMENTAL DESIGN TERMS
% =============================================================================

\item[Ablation Study]
An experimental procedure where individual components (e.g., \texttt{get\_service\_dependency} tool, fast filter, runbook) are removed from the full system to measure their individual contribution to performance metrics.

\item[Cross-System Generalization]
The ability of an aggregation method trained or tuned on one microservice system (e.g., Online Boutique) to maintain performance on unseen systems (e.g., Train Ticket) without retraining.

\item[Statistical Significance Test]
Paired Wilcoxon signed-rank test (non-parametric) used to compare metric distributions across the four arms on the same failure cases, with significance level $\alpha = 0.05$.

\item[Metric Gaming]
The phenomenon where a method optimizes a single metric (e.g., ARR) at the expense of operational utility (e.g., grouping all alerts into one group achieves ARR $\approx$ 100\% but RCPR = 0\%). Mitigated by the multi-metric framework (ARR, RCPR, F1, Purity).

% =============================================================================
% INDUSTRY STANDARDS & REFERENCES
% =============================================================================

\item[Google SRE Principles]
Foundational Site Reliability Engineering practices documented in \textit{Site Reliability Engineering} (Ch.~6, 10, 11) and \textit{The Site Reliability Workbook} (Ch.~5), emphasizing actionable alerts, maximum 2 incidents per 12-hour shift, and alerting on SLOs rather than raw metrics~\cite{google-sre-book, google-sre-workbook}.

\item[Observability Pillars]
The three canonical telemetry types: \textit{Logs} (immutable event records), \textit{Metrics} (aggregated numerical measurements), and \textit{Traces} (distributed request flow records)~\cite{grafana2025}.

\end{description}

% =============================================================================
% LOCAL BIBLIOGRAPHY MACROS (if not using main document's .bib)
% Include these in your main .bib or define via IEEEabrv.bib/IEEEfull.bib
% =============================================================================
% @string{srebook = "Site Reliability Engineering: How Google Runs Production Systems"}
% @string{sreworkbook = "The Site Reliability Workbook: Practical Ways to Implement SRE"}
% @article{pagerduty2020, ...}
% @article{grafana2025, ...}
% @inproceedings{kuang2024cola, ...}
% @inproceedings{pham2025rcaeval, ...}
% @article{zheng2024lemma, ...}
% @inproceedings{yu2025alertguardian, ...}
% @manual{prometheus-alertmanager, ...}
% @article{sbert, ...}
% @article{drain3, ...}
% @article{hdbscan, ...}
% @article{wei2022cot, ...} in main IEEEtran document
% Compiles with: pdflatex main.tex (where main.tex includes this file)
% =============================================================================

% -----------------------------------------------------------------------------
% USAGE IN MAIN DOCUMENT:
% \section*{Glossary of Terms}
% \addcontentsline{toc}{section}{Glossary of Terms}
% % =============================================================================
% GLOSSARY OF TERMS — Alert Aggregation in Microservices
% Standalone file for \input{glossary} in main IEEEtran document
% Compiles with: pdflatex main.tex (where main.tex includes this file)
% =============================================================================

% -----------------------------------------------------------------------------
% USAGE IN MAIN DOCUMENT:
% \section*{Glossary of Terms}
% \addcontentsline{toc}{section}{Glossary of Terms}
% \input{glossary}
% -----------------------------------------------------------------------------

\begin{description}

% =============================================================================
% CORE PROBLEM TERMS
% =============================================================================

\item[Alert Fatigue]
A cognitive state experienced by on-call engineers when the volume of alerts exceeds human cognitive capacity, leading to desensitization, missed critical alerts, and delayed incident response~\cite{pagerduty2020, grafana2025}.

\item[Alert Storm]
A sudden surge of alerts triggered by a single root cause propagating through dependent microservices, resulting in hundreds to thousands of alerts within minutes~\cite{kuang2024cola}.

\item[Noisy Alert Warning]
A warning-level alert generated by the monitoring system that is technically correct but does not require meaningful remediation action from an on-call engineer, classified into four categories: \textit{Duplicate}, \textit{Transient}, \textit{Unactionable}, and \textit{Cascading} (defined below).

\item[Early Warning Phase]
The pre-failure period during which a microservice exhibits performance degradation (e.g., increased latency, resource saturation) but has not yet violated Service Level Objectives (SLOs) or caused user-facing outages.

\item[Actionable Alert]
An alert that satisfies the Google SRE principle: ``every alert sent to an on-call engineer must require immediate human intervention''~\cite{google-sre-ch10, google-sre-ch11}.

% =============================================================================
% FOUR CATEGORIES OF NOISY ALERT WARNINGS
% =============================================================================

\item[Duplicate Warning Alert]
Multiple alerts originating from the same service with identical or near-identical message templates within a short time window ($\Delta t \le 60$~seconds), representing repeated observations of the same underlying condition rather than distinct incidents.

\item[Transient Warning Alert]
A warning alert that occurs during normal system operation (i.e., before fault injection time $T_{\text{inject}}$ in RCAEval) and resolves spontaneously without human intervention, typically caused by temporary network fluctuations or brief resource spikes.

\item[Unactionable Warning Alert]
A technically accurate warning alert for which no Standard Operating Procedure (SOP) or runbook exists, leaving the on-call engineer without a defined remediation path; the alert consumes attention but yields no operational value.

\item[Cascading Warning Alert]
A warning alert emitted by a downstream service that is a symptom of an upstream root cause; remediating the alerting service does not resolve the underlying issue, which originates in a different service (the root cause service).

% =============================================================================
% METHODOLOGY: FOUR EXPERIMENTAL ARMS
% =============================================================================

\item[Arm 1: Rule-Based Aggregation (Baseline)]
An alert grouping approach that clusters alerts based on static label matching (e.g., service name, alert name, severity) and fixed time windows (\texttt{group\_wait}, \texttt{group\_interval}), as implemented in Prometheus Alertmanager~\cite{prometheus-alertmanager}. Lacks awareness of service dependencies and semantic content.

\item[Arm 2: Semantic Similarity-Based Aggregation]
A method that parses raw logs into templates (via Drain3~\cite{drain3}), encodes templates into dense vector embeddings using Sentence-BERT~\cite{sbert}, and clusters alerts using density-based algorithms (HDBSCAN/DBSCAN) based on cosine similarity. Operates without topology knowledge.

\item[Arm 3: Temporal-Spatial Correlation Aggregation]
A topology-aware method that constructs a static service dependency graph (caller-callee adjacency matrix) from microservice architecture, then correlates alerts occurring within a sliding time window ($\Delta t \in \{30, 60, 120\}$~seconds) if a dependency path of length $k \le 2$ exists between the alerting services.

\item[Arm 4: LLM-Based Agentic Aggregation (Proposed)]
A hybrid two-stage approach: (1) a fast filter removes obvious duplicates and performs time-window pre-grouping; (2) an LLM agent (GPT-4o-mini / DeepSeek-V3) reasons over the remaining clusters using Chain-of-Thought prompting and three tools: \texttt{get\_related\_alerts}, \texttt{get\_service\_dependency}, and optionally \texttt{get\_runbook}. Limited to a maximum of two tool calls per grouping decision.

% =============================================================================
% TECHNICAL CONCEPTS & ARCHITECTURE
% =============================================================================

\item[Service Topology (Service Dependency Graph)]
A directed graph $G = (V, E)$ where vertices $V$ represent microservices and edges $E$ represent synchronous (HTTP/gRPC) or asynchronous (message queue) call relationships. Used to distinguish root cause services from downstream symptom services.

\item[Semantic Similarity]
A measure of meaning-based resemblance between two log messages, computed as the cosine similarity between their dense vector embeddings (typically 384-dimensional from \texttt{sentence-transformers/all-MiniLM-L6-v2}). Captures paraphrases and synonyms that lexical matching misses.

\item[Runbook Reasoning]
The process by which an LLM agent consults a synthetic runbook (a structured troubleshooting guide for a specific fault type) to determine whether an alert is actionable, requires escalation, or can be safely aggregated as noise.

\item[Tool-Calling Agent]
An LLM agent augmented with a predefined set of functions (tools) it can invoke during reasoning: \texttt{get\_related\_alerts(service, time\_window)}, \texttt{get\_service\_dependency(service)}, and \texttt{get\_runbook(fault\_type)}. Tool outputs are fed back into the agent's context for subsequent reasoning steps.

\item[Chain-of-Thought (CoT) Prompting]
A prompting technique that instructs the LLM to generate intermediate reasoning steps before producing a final answer, improving accuracy on multi-step causal inference tasks~\cite{wei2022cot}.

\item[Fast Filter (Deduplication + Time-Window Pre-Grouping)]
A lightweight deterministic preprocessing stage that removes exact duplicate alerts and performs initial grouping by service and time window, reducing the alert volume passed to the LLM agent by an estimated 70--80\%.

\item[Synthetic Runbook]
A researcher-authored troubleshooting document for each fault type in the benchmark dataset (e.g., CPU stress, memory leak, network latency), containing symptoms, diagnosis steps, and remediation actions. Not a production SOP; used solely for experimental evaluation.

% =============================================================================
% EVALUATION METRICS
% =============================================================================

\item[Alert Reduction Ratio (ARR)]
The primary metric quantifying alert volume reduction:
\begin{equation}
\text{ARR} = \left(1 - \frac{N_{\text{groups}}}{N_{\text{raw\_alerts}}}\right) \times 100\%
\end{equation}
where $N_{\text{raw\_alerts}}$ is the number of warning alerts before aggregation and $N_{\text{groups}}$ is the number of aggregated groups presented to the engineer. Target: $\text{ARR} \ge 60\%$.

\item[Root Cause Preservation Rate (RCPR)]
The critical safety metric ensuring the root cause evidence is not lost during aggregation. It is defined at two granularities. At service granularity, which is the primary metric and spans every case that exposes logs ($M = 359$):
\begin{equation}
\text{RCPR}_{\text{svc}} = \frac{1}{M} \sum_{i=1}^{M} \mathbb{I}(s_{\text{root}}^{(i)} \in \text{PrimaryGroup}_i) \times 100\%
\end{equation}
At alert granularity, a secondary metric restricted to the $M'$ cases whose ground-truth \texttt{root\_cause.txt} indicator line is itself a warning-tier log ($M' = 4$):
\begin{equation}
\text{RCPR}_{\text{alert}} = \frac{1}{M'} \sum_{i=1}^{M'} \mathbb{I}(\text{RootCauseAlert}_i \in \text{PrimaryGroup}_i) \times 100\%
\end{equation}
where $\mathbb{I}$ is the indicator function. Target: $\text{RCPR} \ge 95\%$ (mandatory).

\item[Pairwise Grouping Precision ($P_p$)]
The fraction of alert pairs placed in the same group that truly share the same root cause:
\begin{equation}
P_p = \frac{|\{(a_x, a_y) : \text{same group} \land \text{same root cause}\}|}{|\{(a_x, a_y) : \text{same group}\}|}
\end{equation}

\item[Pairwise Grouping Recall ($R_p$)]
The fraction of alert pairs sharing the same root cause that are successfully grouped together:
\begin{equation}
R_p = \frac{|\{(a_x, a_y) : \text{same group} \land \text{same root cause}\}|}{|\{(a_x, a_y) : \text{same root cause}\}|}
\end{equation}

\item[Pairwise Grouping F1-Score ($F1_p$)]
The harmonic mean of pairwise precision and recall:
\begin{equation}
F1_p = 2 \times \frac{P_p \times R_p}{P_p + R_p}
\end{equation}
Target: $F1_p \ge 0.85$.

\item[Group Purity]
The average homogeneity of root causes within each generated group:
\begin{equation}
\text{Purity} = \frac{1}{N_{\text{groups}}} \sum_{g=1}^{N_{\text{groups}}} \frac{\max_{c} |\{a \in g : \text{root\_cause}(a) = c\}|}{|g|}
\end{equation}
Target: $\text{Purity} \ge 0.80$.

\item[Processing Latency]
The wall-clock time required to process 100 warning alerts through the aggregation pipeline, measured in milliseconds. For Arm 4, includes LLM inference time and token overhead.

% =============================================================================
% DATASET & BENCHMARK TERMS
% =============================================================================

\item[RCAEval]
A benchmark for root cause analysis in microservice systems comprising 735 failure cases across three systems (Online Boutique, Sock Shop, Train Ticket) with logs, metrics, traces, and root cause labels~\cite{pham2025rcaeval}. Licensed under MIT. Direct inspection shows that only \textbf{359 of the 735 cases ship raw log telemetry} (suite RE1 is metric-only) and that the log schema contains \texttt{timestamp}, \texttt{container\_name}, and \texttt{message} but \textbf{no severity column}.

\item[Severity Tier Assignment]
The three-tier procedure by which this work assigns the severity tier $\ell_i$, since RCAEval supplies none: \emph{Tier 1 (explicit)} parses a level token from the message text (Sock Shop, Train Ticket); \emph{Tier 2 (frequency-inferred)} flags a template as anomalous when its fault-window frequency exceeds its normal-window frequency by a ratio $R$ (Online Boutique); \emph{Tier 3 (ground-truth marker)} tags the \texttt{root\_cause.txt} indicator line as \texttt{ROOT\_CAUSE} for the eight annotated cases.

\item[LEMMA-RCA]
A large-scale multi-modal dataset for root cause analysis containing hundreds of microservice entities and four fault types (DDoS, storage failure, resource contention, noisy neighbor)~\cite{zheng2024lemma}. Licensed under CC-BY-NC-4.0 (non-commercial). Its log pipeline aggregates raw log lines into three-dimensional time series using golden-signal keywords that cover only \texttt{error}, \texttt{exception}, and \texttt{critical} --- notably \textbf{not} \texttt{warning} --- so it offers no independent warning tier.

\item[Fault Injection]
The controlled introduction of a failure (e.g., CPU stress, network latency, packet loss) into a microservice system at a known timestamp $T_{\text{inject}}$ to generate labeled failure cases for benchmarking.

\item[Normal Window]
The interval $[T_{\text{start}}, T_{\text{inject}})$ preceding fault injection, during which the system operates normally; alerts occurring here are candidates for \textit{Transient} noise. Note that the dataset fields \texttt{normal\_timesteps} and \texttt{faulty\_timesteps} are integer counts of sample points, not timestamp lists, and therefore define these windows only indirectly.

\item[Fault Window]
The interval $[T_{\text{inject}}, T_{\text{end}}]$ following fault injection, during which the system exhibits degradation; alerts here are labeled as \textit{Signal} (if from the root cause service) or \textit{Cascading} (if from a service reachable from the root cause service in the call graph).

\item[Ground Truth Root Cause Label]
The authoritative annotation of which service is the root cause of a failure case, provided by the RCAEval benchmark. Used to compute RCPR and pairwise metrics.

% =============================================================================
% BASELINE & INFRASTRUCTURE TERMS
% =============================================================================

\item[Prometheus Alertmanager]
The de facto standard alert routing and grouping component in the Cloud Native Computing Foundation (CNCF) ecosystem. Groups alerts by label matchers and time windows (\texttt{group\_wait}, \texttt{group\_interval}, \texttt{group\_by})~\cite{prometheus-alertmanager}.

\item[Group Wait (\texttt{group\_wait})]
The initial delay after the first alert in a group before sending a notification, allowing additional alerts to arrive and be grouped. Default: 30~seconds.

\item[Group Interval (\texttt{group\_interval})]
The minimum interval between subsequent notifications for the same group. Default: 5~minutes.

\item[Drain3]
An online log parsing algorithm that extracts log templates from raw log streams by clustering log messages based on token similarity, producing parameterized templates with variable placeholders~\cite{drain3}.

\item[Sentence-BERT (S-BERT)]
A modification of BERT that uses siamese and triplet network structures to derive semantically meaningful sentence embeddings, optimized for cosine similarity comparison~\cite{sbert}.

\item[HDBSCAN]
Hierarchical Density-Based Spatial Clustering of Applications with Noise; a density-based clustering algorithm that does not require specifying the number of clusters and can identify outliers (noise points)~\cite{hdbscan}.

% =============================================================================
% EXPERIMENTAL DESIGN TERMS
% =============================================================================

\item[Ablation Study]
An experimental procedure where individual components (e.g., \texttt{get\_service\_dependency} tool, fast filter, runbook) are removed from the full system to measure their individual contribution to performance metrics.

\item[Cross-System Generalization]
The ability of an aggregation method trained or tuned on one microservice system (e.g., Online Boutique) to maintain performance on unseen systems (e.g., Train Ticket) without retraining.

\item[Statistical Significance Test]
Paired Wilcoxon signed-rank test (non-parametric) used to compare metric distributions across the four arms on the same failure cases, with significance level $\alpha = 0.05$.

\item[Metric Gaming]
The phenomenon where a method optimizes a single metric (e.g., ARR) at the expense of operational utility (e.g., grouping all alerts into one group achieves ARR $\approx$ 100\% but RCPR = 0\%). Mitigated by the multi-metric framework (ARR, RCPR, F1, Purity).

% =============================================================================
% INDUSTRY STANDARDS & REFERENCES
% =============================================================================

\item[Google SRE Principles]
Foundational Site Reliability Engineering practices documented in \textit{Site Reliability Engineering} (Ch.~6, 10, 11) and \textit{The Site Reliability Workbook} (Ch.~5), emphasizing actionable alerts, maximum 2 incidents per 12-hour shift, and alerting on SLOs rather than raw metrics~\cite{google-sre-book, google-sre-workbook}.

\item[Observability Pillars]
The three canonical telemetry types: \textit{Logs} (immutable event records), \textit{Metrics} (aggregated numerical measurements), and \textit{Traces} (distributed request flow records)~\cite{grafana2025}.

\end{description}

% =============================================================================
% LOCAL BIBLIOGRAPHY MACROS (if not using main document's .bib)
% Include these in your main .bib or define via IEEEabrv.bib/IEEEfull.bib
% =============================================================================
% @string{srebook = "Site Reliability Engineering: How Google Runs Production Systems"}
% @string{sreworkbook = "The Site Reliability Workbook: Practical Ways to Implement SRE"}
% @article{pagerduty2020, ...}
% @article{grafana2025, ...}
% @inproceedings{kuang2024cola, ...}
% @inproceedings{pham2025rcaeval, ...}
% @article{zheng2024lemma, ...}
% @inproceedings{yu2025alertguardian, ...}
% @manual{prometheus-alertmanager, ...}
% @article{sbert, ...}
% @article{drain3, ...}
% @article{hdbscan, ...}
% @article{wei2022cot, ...}
% -----------------------------------------------------------------------------

\begin{description}

% =============================================================================
% CORE PROBLEM TERMS
% =============================================================================

\item[Alert Fatigue]
A cognitive state experienced by on-call engineers when the volume of alerts exceeds human cognitive capacity, leading to desensitization, missed critical alerts, and delayed incident response~\cite{pagerduty2020, grafana2025}.

\item[Alert Storm]
A sudden surge of alerts triggered by a single root cause propagating through dependent microservices, resulting in hundreds to thousands of alerts within minutes~\cite{kuang2024cola}.

\item[Noisy Alert Warning]
A warning-level alert generated by the monitoring system that is technically correct but does not require meaningful remediation action from an on-call engineer, classified into four categories: \textit{Duplicate}, \textit{Transient}, \textit{Unactionable}, and \textit{Cascading} (defined below).

\item[Early Warning Phase]
The pre-failure period during which a microservice exhibits performance degradation (e.g., increased latency, resource saturation) but has not yet violated Service Level Objectives (SLOs) or caused user-facing outages.

\item[Actionable Alert]
An alert that satisfies the Google SRE principle: ``every alert sent to an on-call engineer must require immediate human intervention''~\cite{google-sre-ch10, google-sre-ch11}.

% =============================================================================
% FOUR CATEGORIES OF NOISY ALERT WARNINGS
% =============================================================================

\item[Duplicate Warning Alert]
Multiple alerts originating from the same service with identical or near-identical message templates within a short time window ($\Delta t \le 60$~seconds), representing repeated observations of the same underlying condition rather than distinct incidents.

\item[Transient Warning Alert]
A warning alert that occurs during normal system operation (i.e., before fault injection time $T_{\text{inject}}$ in RCAEval) and resolves spontaneously without human intervention, typically caused by temporary network fluctuations or brief resource spikes.

\item[Unactionable Warning Alert]
A technically accurate warning alert for which no Standard Operating Procedure (SOP) or runbook exists, leaving the on-call engineer without a defined remediation path; the alert consumes attention but yields no operational value.

\item[Cascading Warning Alert]
A warning alert emitted by a downstream service that is a symptom of an upstream root cause; remediating the alerting service does not resolve the underlying issue, which originates in a different service (the root cause service).

% =============================================================================
% METHODOLOGY: FOUR EXPERIMENTAL ARMS
% =============================================================================

\item[Arm 1: Rule-Based Aggregation (Baseline)]
An alert grouping approach that clusters alerts based on static label matching (e.g., service name, alert name, severity) and fixed time windows (\texttt{group\_wait}, \texttt{group\_interval}), as implemented in Prometheus Alertmanager~\cite{prometheus-alertmanager}. Lacks awareness of service dependencies and semantic content.

\item[Arm 2: Semantic Similarity-Based Aggregation]
A method that parses raw logs into templates (via Drain3~\cite{drain3}), encodes templates into dense vector embeddings using Sentence-BERT~\cite{sbert}, and clusters alerts using density-based algorithms (HDBSCAN/DBSCAN) based on cosine similarity. Operates without topology knowledge.

\item[Arm 3: Temporal-Spatial Correlation Aggregation]
A topology-aware method that constructs a static service dependency graph (caller-callee adjacency matrix) from microservice architecture, then correlates alerts occurring within a sliding time window ($\Delta t \in \{30, 60, 120\}$~seconds) if a dependency path of length $k \le 2$ exists between the alerting services.

\item[Arm 4: LLM-Based Agentic Aggregation (Proposed)]
A hybrid two-stage approach: (1) a fast filter removes obvious duplicates and performs time-window pre-grouping; (2) an LLM agent (GPT-4o-mini / DeepSeek-V3) reasons over the remaining clusters using Chain-of-Thought prompting and three tools: \texttt{get\_related\_alerts}, \texttt{get\_service\_dependency}, and optionally \texttt{get\_runbook}. Limited to a maximum of two tool calls per grouping decision.

% =============================================================================
% TECHNICAL CONCEPTS & ARCHITECTURE
% =============================================================================

\item[Service Topology (Service Dependency Graph)]
A directed graph $G = (V, E)$ where vertices $V$ represent microservices and edges $E$ represent synchronous (HTTP/gRPC) or asynchronous (message queue) call relationships. Used to distinguish root cause services from downstream symptom services.

\item[Semantic Similarity]
A measure of meaning-based resemblance between two log messages, computed as the cosine similarity between their dense vector embeddings (typically 384-dimensional from \texttt{sentence-transformers/all-MiniLM-L6-v2}). Captures paraphrases and synonyms that lexical matching misses.

\item[Runbook Reasoning]
The process by which an LLM agent consults a synthetic runbook (a structured troubleshooting guide for a specific fault type) to determine whether an alert is actionable, requires escalation, or can be safely aggregated as noise.

\item[Tool-Calling Agent]
An LLM agent augmented with a predefined set of functions (tools) it can invoke during reasoning: \texttt{get\_related\_alerts(service, time\_window)}, \texttt{get\_service\_dependency(service)}, and \texttt{get\_runbook(fault\_type)}. Tool outputs are fed back into the agent's context for subsequent reasoning steps.

\item[Chain-of-Thought (CoT) Prompting]
A prompting technique that instructs the LLM to generate intermediate reasoning steps before producing a final answer, improving accuracy on multi-step causal inference tasks~\cite{wei2022cot}.

\item[Fast Filter (Deduplication + Time-Window Pre-Grouping)]
A lightweight deterministic preprocessing stage that removes exact duplicate alerts and performs initial grouping by service and time window, reducing the alert volume passed to the LLM agent by an estimated 70--80\%.

\item[Synthetic Runbook]
A researcher-authored troubleshooting document for each fault type in the benchmark dataset (e.g., CPU stress, memory leak, network latency), containing symptoms, diagnosis steps, and remediation actions. Not a production SOP; used solely for experimental evaluation.

% =============================================================================
% EVALUATION METRICS
% =============================================================================

\item[Alert Reduction Ratio (ARR)]
The primary metric quantifying alert volume reduction:
\begin{equation}
\text{ARR} = \left(1 - \frac{N_{\text{groups}}}{N_{\text{raw\_alerts}}}\right) \times 100\%
\end{equation}
where $N_{\text{raw\_alerts}}$ is the number of warning alerts before aggregation and $N_{\text{groups}}$ is the number of aggregated groups presented to the engineer. Target: $\text{ARR} \ge 60\%$.

\item[Root Cause Preservation Rate (RCPR)]
The critical safety metric ensuring the root cause evidence is not lost during aggregation. It is defined at two granularities. At service granularity, which is the primary metric and spans every case that exposes logs ($M = 359$):
\begin{equation}
\text{RCPR}_{\text{svc}} = \frac{1}{M} \sum_{i=1}^{M} \mathbb{I}(s_{\text{root}}^{(i)} \in \text{PrimaryGroup}_i) \times 100\%
\end{equation}
At alert granularity, a secondary metric restricted to the $M'$ cases whose ground-truth \texttt{root\_cause.txt} indicator line is itself a warning-tier log ($M' = 4$):
\begin{equation}
\text{RCPR}_{\text{alert}} = \frac{1}{M'} \sum_{i=1}^{M'} \mathbb{I}(\text{RootCauseAlert}_i \in \text{PrimaryGroup}_i) \times 100\%
\end{equation}
where $\mathbb{I}$ is the indicator function. Target: $\text{RCPR} \ge 95\%$ (mandatory).

\item[Pairwise Grouping Precision ($P_p$)]
The fraction of alert pairs placed in the same group that truly share the same root cause:
\begin{equation}
P_p = \frac{|\{(a_x, a_y) : \text{same group} \land \text{same root cause}\}|}{|\{(a_x, a_y) : \text{same group}\}|}
\end{equation}

\item[Pairwise Grouping Recall ($R_p$)]
The fraction of alert pairs sharing the same root cause that are successfully grouped together:
\begin{equation}
R_p = \frac{|\{(a_x, a_y) : \text{same group} \land \text{same root cause}\}|}{|\{(a_x, a_y) : \text{same root cause}\}|}
\end{equation}

\item[Pairwise Grouping F1-Score ($F1_p$)]
The harmonic mean of pairwise precision and recall:
\begin{equation}
F1_p = 2 \times \frac{P_p \times R_p}{P_p + R_p}
\end{equation}
Target: $F1_p \ge 0.85$.

\item[Group Purity]
The average homogeneity of root causes within each generated group:
\begin{equation}
\text{Purity} = \frac{1}{N_{\text{groups}}} \sum_{g=1}^{N_{\text{groups}}} \frac{\max_{c} |\{a \in g : \text{root\_cause}(a) = c\}|}{|g|}
\end{equation}
Target: $\text{Purity} \ge 0.80$.

\item[Processing Latency]
The wall-clock time required to process 100 warning alerts through the aggregation pipeline, measured in milliseconds. For Arm 4, includes LLM inference time and token overhead.

% =============================================================================
% DATASET & BENCHMARK TERMS
% =============================================================================

\item[RCAEval]
A benchmark for root cause analysis in microservice systems comprising 735 failure cases across three systems (Online Boutique, Sock Shop, Train Ticket) with logs, metrics, traces, and root cause labels~\cite{pham2025rcaeval}. Licensed under MIT. Direct inspection shows that only \textbf{359 of the 735 cases ship raw log telemetry} (suite RE1 is metric-only) and that the log schema contains \texttt{timestamp}, \texttt{container\_name}, and \texttt{message} but \textbf{no severity column}.

\item[Severity Tier Assignment]
The three-tier procedure by which this work assigns the severity tier $\ell_i$, since RCAEval supplies none: \emph{Tier 1 (explicit)} parses a level token from the message text (Sock Shop, Train Ticket); \emph{Tier 2 (frequency-inferred)} flags a template as anomalous when its fault-window frequency exceeds its normal-window frequency by a ratio $R$ (Online Boutique); \emph{Tier 3 (ground-truth marker)} tags the \texttt{root\_cause.txt} indicator line as \texttt{ROOT\_CAUSE} for the eight annotated cases.

\item[LEMMA-RCA]
A large-scale multi-modal dataset for root cause analysis containing hundreds of microservice entities and four fault types (DDoS, storage failure, resource contention, noisy neighbor)~\cite{zheng2024lemma}. Licensed under CC-BY-NC-4.0 (non-commercial). Its log pipeline aggregates raw log lines into three-dimensional time series using golden-signal keywords that cover only \texttt{error}, \texttt{exception}, and \texttt{critical} --- notably \textbf{not} \texttt{warning} --- so it offers no independent warning tier.

\item[Fault Injection]
The controlled introduction of a failure (e.g., CPU stress, network latency, packet loss) into a microservice system at a known timestamp $T_{\text{inject}}$ to generate labeled failure cases for benchmarking.

\item[Normal Window]
The interval $[T_{\text{start}}, T_{\text{inject}})$ preceding fault injection, during which the system operates normally; alerts occurring here are candidates for \textit{Transient} noise. Note that the dataset fields \texttt{normal\_timesteps} and \texttt{faulty\_timesteps} are integer counts of sample points, not timestamp lists, and therefore define these windows only indirectly.

\item[Fault Window]
The interval $[T_{\text{inject}}, T_{\text{end}}]$ following fault injection, during which the system exhibits degradation; alerts here are labeled as \textit{Signal} (if from the root cause service) or \textit{Cascading} (if from a service reachable from the root cause service in the call graph).

\item[Ground Truth Root Cause Label]
The authoritative annotation of which service is the root cause of a failure case, provided by the RCAEval benchmark. Used to compute RCPR and pairwise metrics.

% =============================================================================
% BASELINE & INFRASTRUCTURE TERMS
% =============================================================================

\item[Prometheus Alertmanager]
The de facto standard alert routing and grouping component in the Cloud Native Computing Foundation (CNCF) ecosystem. Groups alerts by label matchers and time windows (\texttt{group\_wait}, \texttt{group\_interval}, \texttt{group\_by})~\cite{prometheus-alertmanager}.

\item[Group Wait (\texttt{group\_wait})]
The initial delay after the first alert in a group before sending a notification, allowing additional alerts to arrive and be grouped. Default: 30~seconds.

\item[Group Interval (\texttt{group\_interval})]
The minimum interval between subsequent notifications for the same group. Default: 5~minutes.

\item[Drain3]
An online log parsing algorithm that extracts log templates from raw log streams by clustering log messages based on token similarity, producing parameterized templates with variable placeholders~\cite{drain3}.

\item[Sentence-BERT (S-BERT)]
A modification of BERT that uses siamese and triplet network structures to derive semantically meaningful sentence embeddings, optimized for cosine similarity comparison~\cite{sbert}.

\item[HDBSCAN]
Hierarchical Density-Based Spatial Clustering of Applications with Noise; a density-based clustering algorithm that does not require specifying the number of clusters and can identify outliers (noise points)~\cite{hdbscan}.

% =============================================================================
% EXPERIMENTAL DESIGN TERMS
% =============================================================================

\item[Ablation Study]
An experimental procedure where individual components (e.g., \texttt{get\_service\_dependency} tool, fast filter, runbook) are removed from the full system to measure their individual contribution to performance metrics.

\item[Cross-System Generalization]
The ability of an aggregation method trained or tuned on one microservice system (e.g., Online Boutique) to maintain performance on unseen systems (e.g., Train Ticket) without retraining.

\item[Statistical Significance Test]
Paired Wilcoxon signed-rank test (non-parametric) used to compare metric distributions across the four arms on the same failure cases, with significance level $\alpha = 0.05$.

\item[Metric Gaming]
The phenomenon where a method optimizes a single metric (e.g., ARR) at the expense of operational utility (e.g., grouping all alerts into one group achieves ARR $\approx$ 100\% but RCPR = 0\%). Mitigated by the multi-metric framework (ARR, RCPR, F1, Purity).

% =============================================================================
% INDUSTRY STANDARDS & REFERENCES
% =============================================================================

\item[Google SRE Principles]
Foundational Site Reliability Engineering practices documented in \textit{Site Reliability Engineering} (Ch.~6, 10, 11) and \textit{The Site Reliability Workbook} (Ch.~5), emphasizing actionable alerts, maximum 2 incidents per 12-hour shift, and alerting on SLOs rather than raw metrics~\cite{google-sre-book, google-sre-workbook}.

\item[Observability Pillars]
The three canonical telemetry types: \textit{Logs} (immutable event records), \textit{Metrics} (aggregated numerical measurements), and \textit{Traces} (distributed request flow records)~\cite{grafana2025}.

\end{description}

% =============================================================================
% LOCAL BIBLIOGRAPHY MACROS (if not using main document's .bib)
% Include these in your main .bib or define via IEEEabrv.bib/IEEEfull.bib
% =============================================================================
% @string{srebook = "Site Reliability Engineering: How Google Runs Production Systems"}
% @string{sreworkbook = "The Site Reliability Workbook: Practical Ways to Implement SRE"}
% @article{pagerduty2020, ...}
% @article{grafana2025, ...}
% @inproceedings{kuang2024cola, ...}
% @inproceedings{pham2025rcaeval, ...}
% @article{zheng2024lemma, ...}
% @inproceedings{yu2025alertguardian, ...}
% @manual{prometheus-alertmanager, ...}
% @article{sbert, ...}
% @article{drain3, ...}
% @article{hdbscan, ...}
% @article{wei2022cot, ...}
% -----------------------------------------------------------------------------

\begin{description}

% =============================================================================
% CORE PROBLEM TERMS
% =============================================================================

\item[Alert Fatigue]
A cognitive state experienced by on-call engineers when the volume of alerts exceeds human cognitive capacity, leading to desensitization, missed critical alerts, and delayed incident response~\cite{pagerduty2020, grafana2025}.

\item[Alert Storm]
A sudden surge of alerts triggered by a single root cause propagating through dependent microservices, resulting in hundreds to thousands of alerts within minutes~\cite{kuang2024cola}.

\item[Noisy Alert Warning]
A warning-level alert generated by the monitoring system that is technically correct but does not require meaningful remediation action from an on-call engineer, classified into four categories: \textit{Duplicate}, \textit{Transient}, \textit{Unactionable}, and \textit{Cascading} (defined below).

\item[Early Warning Phase]
The pre-failure period during which a microservice exhibits performance degradation (e.g., increased latency, resource saturation) but has not yet violated Service Level Objectives (SLOs) or caused user-facing outages.

\item[Actionable Alert]
An alert that satisfies the Google SRE principle: ``every alert sent to an on-call engineer must require immediate human intervention''~\cite{google-sre-ch10, google-sre-ch11}.

% =============================================================================
% FOUR CATEGORIES OF NOISY ALERT WARNINGS
% =============================================================================

\item[Duplicate Warning Alert]
Multiple alerts originating from the same service with identical or near-identical message templates within a short time window ($\Delta t \le 60$~seconds), representing repeated observations of the same underlying condition rather than distinct incidents.

\item[Transient Warning Alert]
A warning alert that occurs during normal system operation (i.e., before fault injection time $T_{\text{inject}}$ in RCAEval) and resolves spontaneously without human intervention, typically caused by temporary network fluctuations or brief resource spikes.

\item[Unactionable Warning Alert]
A technically accurate warning alert for which no Standard Operating Procedure (SOP) or runbook exists, leaving the on-call engineer without a defined remediation path; the alert consumes attention but yields no operational value.

\item[Cascading Warning Alert]
A warning alert emitted by a downstream service that is a symptom of an upstream root cause; remediating the alerting service does not resolve the underlying issue, which originates in a different service (the root cause service).

% =============================================================================
% METHODOLOGY: FOUR EXPERIMENTAL ARMS
% =============================================================================

\item[Arm 1: Rule-Based Aggregation (Baseline)]
An alert grouping approach that clusters alerts based on static label matching (e.g., service name, alert name, severity) and fixed time windows (\texttt{group\_wait}, \texttt{group\_interval}), as implemented in Prometheus Alertmanager~\cite{prometheus-alertmanager}. Lacks awareness of service dependencies and semantic content.

\item[Arm 2: Semantic Similarity-Based Aggregation]
A method that parses raw logs into templates (via Drain3~\cite{drain3}), encodes templates into dense vector embeddings using Sentence-BERT~\cite{sbert}, and clusters alerts using density-based algorithms (HDBSCAN/DBSCAN) based on cosine similarity. Operates without topology knowledge.

\item[Arm 3: Temporal-Spatial Correlation Aggregation]
A topology-aware method that constructs a static service dependency graph (caller-callee adjacency matrix) from microservice architecture, then correlates alerts occurring within a sliding time window ($\Delta t \in \{30, 60, 120\}$~seconds) if a dependency path of length $k \le 2$ exists between the alerting services.

\item[Arm 4: LLM-Based Agentic Aggregation (Proposed)]
A hybrid two-stage approach: (1) a fast filter removes obvious duplicates and performs time-window pre-grouping; (2) an LLM agent (GPT-4o-mini / DeepSeek-V3) reasons over the remaining clusters using Chain-of-Thought prompting and three tools: \texttt{get\_related\_alerts}, \texttt{get\_service\_dependency}, and optionally \texttt{get\_runbook}. Limited to a maximum of two tool calls per grouping decision.

% =============================================================================
% TECHNICAL CONCEPTS & ARCHITECTURE
% =============================================================================

\item[Service Topology (Service Dependency Graph)]
A directed graph $G = (V, E)$ where vertices $V$ represent microservices and edges $E$ represent synchronous (HTTP/gRPC) or asynchronous (message queue) call relationships. Used to distinguish root cause services from downstream symptom services.

\item[Semantic Similarity]
A measure of meaning-based resemblance between two log messages, computed as the cosine similarity between their dense vector embeddings (typically 384-dimensional from \texttt{sentence-transformers/all-MiniLM-L6-v2}). Captures paraphrases and synonyms that lexical matching misses.

\item[Runbook Reasoning]
The process by which an LLM agent consults a synthetic runbook (a structured troubleshooting guide for a specific fault type) to determine whether an alert is actionable, requires escalation, or can be safely aggregated as noise.

\item[Tool-Calling Agent]
An LLM agent augmented with a predefined set of functions (tools) it can invoke during reasoning: \texttt{get\_related\_alerts(service, time\_window)}, \texttt{get\_service\_dependency(service)}, and \texttt{get\_runbook(fault\_type)}. Tool outputs are fed back into the agent's context for subsequent reasoning steps.

\item[Chain-of-Thought (CoT) Prompting]
A prompting technique that instructs the LLM to generate intermediate reasoning steps before producing a final answer, improving accuracy on multi-step causal inference tasks~\cite{wei2022cot}.

\item[Fast Filter (Deduplication + Time-Window Pre-Grouping)]
A lightweight deterministic preprocessing stage that removes exact duplicate alerts and performs initial grouping by service and time window, reducing the alert volume passed to the LLM agent by an estimated 70--80\%.

\item[Synthetic Runbook]
A researcher-authored troubleshooting document for each fault type in the benchmark dataset (e.g., CPU stress, memory leak, network latency), containing symptoms, diagnosis steps, and remediation actions. Not a production SOP; used solely for experimental evaluation.

% =============================================================================
% EVALUATION METRICS
% =============================================================================

\item[Alert Reduction Ratio (ARR)]
The primary metric quantifying alert volume reduction:
\begin{equation}
\text{ARR} = \left(1 - \frac{N_{\text{groups}}}{N_{\text{raw\_alerts}}}\right) \times 100\%
\end{equation}
where $N_{\text{raw\_alerts}}$ is the number of warning alerts before aggregation and $N_{\text{groups}}$ is the number of aggregated groups presented to the engineer. Target: $\text{ARR} \ge 60\%$.

\item[Root Cause Preservation Rate (RCPR)]
The critical safety metric ensuring the root cause evidence is not lost during aggregation. It is defined at two granularities. At service granularity, which is the primary metric and spans every case that exposes logs ($M = 359$):
\begin{equation}
\text{RCPR}_{\text{svc}} = \frac{1}{M} \sum_{i=1}^{M} \mathbb{I}(s_{\text{root}}^{(i)} \in \text{PrimaryGroup}_i) \times 100\%
\end{equation}
At alert granularity, a secondary metric restricted to the $M'$ cases whose ground-truth \texttt{root\_cause.txt} indicator line is itself a warning-tier log ($M' = 4$):
\begin{equation}
\text{RCPR}_{\text{alert}} = \frac{1}{M'} \sum_{i=1}^{M'} \mathbb{I}(\text{RootCauseAlert}_i \in \text{PrimaryGroup}_i) \times 100\%
\end{equation}
where $\mathbb{I}$ is the indicator function. Target: $\text{RCPR} \ge 95\%$ (mandatory).

\item[Pairwise Grouping Precision ($P_p$)]
The fraction of alert pairs placed in the same group that truly share the same root cause:
\begin{equation}
P_p = \frac{|\{(a_x, a_y) : \text{same group} \land \text{same root cause}\}|}{|\{(a_x, a_y) : \text{same group}\}|}
\end{equation}

\item[Pairwise Grouping Recall ($R_p$)]
The fraction of alert pairs sharing the same root cause that are successfully grouped together:
\begin{equation}
R_p = \frac{|\{(a_x, a_y) : \text{same group} \land \text{same root cause}\}|}{|\{(a_x, a_y) : \text{same root cause}\}|}
\end{equation}

\item[Pairwise Grouping F1-Score ($F1_p$)]
The harmonic mean of pairwise precision and recall:
\begin{equation}
F1_p = 2 \times \frac{P_p \times R_p}{P_p + R_p}
\end{equation}
Target: $F1_p \ge 0.85$.

\item[Group Purity]
The average homogeneity of root causes within each generated group:
\begin{equation}
\text{Purity} = \frac{1}{N_{\text{groups}}} \sum_{g=1}^{N_{\text{groups}}} \frac{\max_{c} |\{a \in g : \text{root\_cause}(a) = c\}|}{|g|}
\end{equation}
Target: $\text{Purity} \ge 0.80$.

\item[Processing Latency]
The wall-clock time required to process 100 warning alerts through the aggregation pipeline, measured in milliseconds. For Arm 4, includes LLM inference time and token overhead.

% =============================================================================
% DATASET & BENCHMARK TERMS
% =============================================================================

\item[RCAEval]
A benchmark for root cause analysis in microservice systems comprising 735 failure cases across three systems (Online Boutique, Sock Shop, Train Ticket) with logs, metrics, traces, and root cause labels~\cite{pham2025rcaeval}. Licensed under MIT. Direct inspection shows that only \textbf{359 of the 735 cases ship raw log telemetry} (suite RE1 is metric-only) and that the log schema contains \texttt{timestamp}, \texttt{container\_name}, and \texttt{message} but \textbf{no severity column}.

\item[Severity Tier Assignment]
The three-tier procedure by which this work assigns the severity tier $\ell_i$, since RCAEval supplies none: \emph{Tier 1 (explicit)} parses a level token from the message text (Sock Shop, Train Ticket); \emph{Tier 2 (frequency-inferred)} flags a template as anomalous when its fault-window frequency exceeds its normal-window frequency by a ratio $R$ (Online Boutique); \emph{Tier 3 (ground-truth marker)} tags the \texttt{root\_cause.txt} indicator line as \texttt{ROOT\_CAUSE} for the eight annotated cases.

\item[LEMMA-RCA]
A large-scale multi-modal dataset for root cause analysis containing hundreds of microservice entities and four fault types (DDoS, storage failure, resource contention, noisy neighbor)~\cite{zheng2024lemma}. Licensed under CC-BY-NC-4.0 (non-commercial). Its log pipeline aggregates raw log lines into three-dimensional time series using golden-signal keywords that cover only \texttt{error}, \texttt{exception}, and \texttt{critical} --- notably \textbf{not} \texttt{warning} --- so it offers no independent warning tier.

\item[Fault Injection]
The controlled introduction of a failure (e.g., CPU stress, network latency, packet loss) into a microservice system at a known timestamp $T_{\text{inject}}$ to generate labeled failure cases for benchmarking.

\item[Normal Window]
The interval $[T_{\text{start}}, T_{\text{inject}})$ preceding fault injection, during which the system operates normally; alerts occurring here are candidates for \textit{Transient} noise. Note that the dataset fields \texttt{normal\_timesteps} and \texttt{faulty\_timesteps} are integer counts of sample points, not timestamp lists, and therefore define these windows only indirectly.

\item[Fault Window]
The interval $[T_{\text{inject}}, T_{\text{end}}]$ following fault injection, during which the system exhibits degradation; alerts here are labeled as \textit{Signal} (if from the root cause service) or \textit{Cascading} (if from a service reachable from the root cause service in the call graph).

\item[Ground Truth Root Cause Label]
The authoritative annotation of which service is the root cause of a failure case, provided by the RCAEval benchmark. Used to compute RCPR and pairwise metrics.

% =============================================================================
% BASELINE & INFRASTRUCTURE TERMS
% =============================================================================

\item[Prometheus Alertmanager]
The de facto standard alert routing and grouping component in the Cloud Native Computing Foundation (CNCF) ecosystem. Groups alerts by label matchers and time windows (\texttt{group\_wait}, \texttt{group\_interval}, \texttt{group\_by})~\cite{prometheus-alertmanager}.

\item[Group Wait (\texttt{group\_wait})]
The initial delay after the first alert in a group before sending a notification, allowing additional alerts to arrive and be grouped. Default: 30~seconds.

\item[Group Interval (\texttt{group\_interval})]
The minimum interval between subsequent notifications for the same group. Default: 5~minutes.

\item[Drain3]
An online log parsing algorithm that extracts log templates from raw log streams by clustering log messages based on token similarity, producing parameterized templates with variable placeholders~\cite{drain3}.

\item[Sentence-BERT (S-BERT)]
A modification of BERT that uses siamese and triplet network structures to derive semantically meaningful sentence embeddings, optimized for cosine similarity comparison~\cite{sbert}.

\item[HDBSCAN]
Hierarchical Density-Based Spatial Clustering of Applications with Noise; a density-based clustering algorithm that does not require specifying the number of clusters and can identify outliers (noise points)~\cite{hdbscan}.

% =============================================================================
% EXPERIMENTAL DESIGN TERMS
% =============================================================================

\item[Ablation Study]
An experimental procedure where individual components (e.g., \texttt{get\_service\_dependency} tool, fast filter, runbook) are removed from the full system to measure their individual contribution to performance metrics.

\item[Cross-System Generalization]
The ability of an aggregation method trained or tuned on one microservice system (e.g., Online Boutique) to maintain performance on unseen systems (e.g., Train Ticket) without retraining.

\item[Statistical Significance Test]
Paired Wilcoxon signed-rank test (non-parametric) used to compare metric distributions across the four arms on the same failure cases, with significance level $\alpha = 0.05$.

\item[Metric Gaming]
The phenomenon where a method optimizes a single metric (e.g., ARR) at the expense of operational utility (e.g., grouping all alerts into one group achieves ARR $\approx$ 100\% but RCPR = 0\%). Mitigated by the multi-metric framework (ARR, RCPR, F1, Purity).

% =============================================================================
% INDUSTRY STANDARDS & REFERENCES
% =============================================================================

\item[Google SRE Principles]
Foundational Site Reliability Engineering practices documented in \textit{Site Reliability Engineering} (Ch.~6, 10, 11) and \textit{The Site Reliability Workbook} (Ch.~5), emphasizing actionable alerts, maximum 2 incidents per 12-hour shift, and alerting on SLOs rather than raw metrics~\cite{google-sre-book, google-sre-workbook}.

\item[Observability Pillars]
The three canonical telemetry types: \textit{Logs} (immutable event records), \textit{Metrics} (aggregated numerical measurements), and \textit{Traces} (distributed request flow records)~\cite{grafana2025}.

\end{description}

% =============================================================================
% LOCAL BIBLIOGRAPHY MACROS (if not using main document's .bib)
% Include these in your main .bib or define via IEEEabrv.bib/IEEEfull.bib
% =============================================================================
% @string{srebook = "Site Reliability Engineering: How Google Runs Production Systems"}
% @string{sreworkbook = "The Site Reliability Workbook: Practical Ways to Implement SRE"}
% @article{pagerduty2020, ...}
% @article{grafana2025, ...}
% @inproceedings{kuang2024cola, ...}
% @inproceedings{pham2025rcaeval, ...}
% @article{zheng2024lemma, ...}
% @inproceedings{yu2025alertguardian, ...}
% @manual{prometheus-alertmanager, ...}
% @article{sbert, ...}
% @article{drain3, ...}
% @article{hdbscan, ...}
% @article{wei2022cot, ...} in main IEEEtran document
% Compiles with: pdflatex main.tex (where main.tex includes this file)
% =============================================================================

% -----------------------------------------------------------------------------
% USAGE IN MAIN DOCUMENT:
% \section*{Glossary of Terms}
% \addcontentsline{toc}{section}{Glossary of Terms}
% % =============================================================================
% GLOSSARY OF TERMS — Alert Aggregation in Microservices
% Standalone file for % =============================================================================
% GLOSSARY OF TERMS — Alert Aggregation in Microservices
% Standalone file for % =============================================================================
% GLOSSARY OF TERMS — Alert Aggregation in Microservices
% Standalone file for \input{glossary} in main IEEEtran document
% Compiles with: pdflatex main.tex (where main.tex includes this file)
% =============================================================================

% -----------------------------------------------------------------------------
% USAGE IN MAIN DOCUMENT:
% \section*{Glossary of Terms}
% \addcontentsline{toc}{section}{Glossary of Terms}
% \input{glossary}
% -----------------------------------------------------------------------------

\begin{description}

% =============================================================================
% CORE PROBLEM TERMS
% =============================================================================

\item[Alert Fatigue]
A cognitive state experienced by on-call engineers when the volume of alerts exceeds human cognitive capacity, leading to desensitization, missed critical alerts, and delayed incident response~\cite{pagerduty2020, grafana2025}.

\item[Alert Storm]
A sudden surge of alerts triggered by a single root cause propagating through dependent microservices, resulting in hundreds to thousands of alerts within minutes~\cite{kuang2024cola}.

\item[Noisy Alert Warning]
A warning-level alert generated by the monitoring system that is technically correct but does not require meaningful remediation action from an on-call engineer, classified into four categories: \textit{Duplicate}, \textit{Transient}, \textit{Unactionable}, and \textit{Cascading} (defined below).

\item[Early Warning Phase]
The pre-failure period during which a microservice exhibits performance degradation (e.g., increased latency, resource saturation) but has not yet violated Service Level Objectives (SLOs) or caused user-facing outages.

\item[Actionable Alert]
An alert that satisfies the Google SRE principle: ``every alert sent to an on-call engineer must require immediate human intervention''~\cite{google-sre-ch10, google-sre-ch11}.

% =============================================================================
% FOUR CATEGORIES OF NOISY ALERT WARNINGS
% =============================================================================

\item[Duplicate Warning Alert]
Multiple alerts originating from the same service with identical or near-identical message templates within a short time window ($\Delta t \le 60$~seconds), representing repeated observations of the same underlying condition rather than distinct incidents.

\item[Transient Warning Alert]
A warning alert that occurs during normal system operation (i.e., before fault injection time $T_{\text{inject}}$ in RCAEval) and resolves spontaneously without human intervention, typically caused by temporary network fluctuations or brief resource spikes.

\item[Unactionable Warning Alert]
A technically accurate warning alert for which no Standard Operating Procedure (SOP) or runbook exists, leaving the on-call engineer without a defined remediation path; the alert consumes attention but yields no operational value.

\item[Cascading Warning Alert]
A warning alert emitted by a downstream service that is a symptom of an upstream root cause; remediating the alerting service does not resolve the underlying issue, which originates in a different service (the root cause service).

% =============================================================================
% METHODOLOGY: FOUR EXPERIMENTAL ARMS
% =============================================================================

\item[Arm 1: Rule-Based Aggregation (Baseline)]
An alert grouping approach that clusters alerts based on static label matching (e.g., service name, alert name, severity) and fixed time windows (\texttt{group\_wait}, \texttt{group\_interval}), as implemented in Prometheus Alertmanager~\cite{prometheus-alertmanager}. Lacks awareness of service dependencies and semantic content.

\item[Arm 2: Semantic Similarity-Based Aggregation]
A method that parses raw logs into templates (via Drain3~\cite{drain3}), encodes templates into dense vector embeddings using Sentence-BERT~\cite{sbert}, and clusters alerts using density-based algorithms (HDBSCAN/DBSCAN) based on cosine similarity. Operates without topology knowledge.

\item[Arm 3: Temporal-Spatial Correlation Aggregation]
A topology-aware method that constructs a static service dependency graph (caller-callee adjacency matrix) from microservice architecture, then correlates alerts occurring within a sliding time window ($\Delta t \in \{30, 60, 120\}$~seconds) if a dependency path of length $k \le 2$ exists between the alerting services.

\item[Arm 4: LLM-Based Agentic Aggregation (Proposed)]
A hybrid two-stage approach: (1) a fast filter removes obvious duplicates and performs time-window pre-grouping; (2) an LLM agent (GPT-4o-mini / DeepSeek-V3) reasons over the remaining clusters using Chain-of-Thought prompting and three tools: \texttt{get\_related\_alerts}, \texttt{get\_service\_dependency}, and optionally \texttt{get\_runbook}. Limited to a maximum of two tool calls per grouping decision.

% =============================================================================
% TECHNICAL CONCEPTS & ARCHITECTURE
% =============================================================================

\item[Service Topology (Service Dependency Graph)]
A directed graph $G = (V, E)$ where vertices $V$ represent microservices and edges $E$ represent synchronous (HTTP/gRPC) or asynchronous (message queue) call relationships. Used to distinguish root cause services from downstream symptom services.

\item[Semantic Similarity]
A measure of meaning-based resemblance between two log messages, computed as the cosine similarity between their dense vector embeddings (typically 384-dimensional from \texttt{sentence-transformers/all-MiniLM-L6-v2}). Captures paraphrases and synonyms that lexical matching misses.

\item[Runbook Reasoning]
The process by which an LLM agent consults a synthetic runbook (a structured troubleshooting guide for a specific fault type) to determine whether an alert is actionable, requires escalation, or can be safely aggregated as noise.

\item[Tool-Calling Agent]
An LLM agent augmented with a predefined set of functions (tools) it can invoke during reasoning: \texttt{get\_related\_alerts(service, time\_window)}, \texttt{get\_service\_dependency(service)}, and \texttt{get\_runbook(fault\_type)}. Tool outputs are fed back into the agent's context for subsequent reasoning steps.

\item[Chain-of-Thought (CoT) Prompting]
A prompting technique that instructs the LLM to generate intermediate reasoning steps before producing a final answer, improving accuracy on multi-step causal inference tasks~\cite{wei2022cot}.

\item[Fast Filter (Deduplication + Time-Window Pre-Grouping)]
A lightweight deterministic preprocessing stage that removes exact duplicate alerts and performs initial grouping by service and time window, reducing the alert volume passed to the LLM agent by an estimated 70--80\%.

\item[Synthetic Runbook]
A researcher-authored troubleshooting document for each fault type in the benchmark dataset (e.g., CPU stress, memory leak, network latency), containing symptoms, diagnosis steps, and remediation actions. Not a production SOP; used solely for experimental evaluation.

% =============================================================================
% EVALUATION METRICS
% =============================================================================

\item[Alert Reduction Ratio (ARR)]
The primary metric quantifying alert volume reduction:
\begin{equation}
\text{ARR} = \left(1 - \frac{N_{\text{groups}}}{N_{\text{raw\_alerts}}}\right) \times 100\%
\end{equation}
where $N_{\text{raw\_alerts}}$ is the number of warning alerts before aggregation and $N_{\text{groups}}$ is the number of aggregated groups presented to the engineer. Target: $\text{ARR} \ge 60\%$.

\item[Root Cause Preservation Rate (RCPR)]
The critical safety metric ensuring the root cause evidence is not lost during aggregation. It is defined at two granularities. At service granularity, which is the primary metric and spans every case that exposes logs ($M = 359$):
\begin{equation}
\text{RCPR}_{\text{svc}} = \frac{1}{M} \sum_{i=1}^{M} \mathbb{I}(s_{\text{root}}^{(i)} \in \text{PrimaryGroup}_i) \times 100\%
\end{equation}
At alert granularity, a secondary metric restricted to the $M'$ cases whose ground-truth \texttt{root\_cause.txt} indicator line is itself a warning-tier log ($M' = 4$):
\begin{equation}
\text{RCPR}_{\text{alert}} = \frac{1}{M'} \sum_{i=1}^{M'} \mathbb{I}(\text{RootCauseAlert}_i \in \text{PrimaryGroup}_i) \times 100\%
\end{equation}
where $\mathbb{I}$ is the indicator function. Target: $\text{RCPR} \ge 95\%$ (mandatory).

\item[Pairwise Grouping Precision ($P_p$)]
The fraction of alert pairs placed in the same group that truly share the same root cause:
\begin{equation}
P_p = \frac{|\{(a_x, a_y) : \text{same group} \land \text{same root cause}\}|}{|\{(a_x, a_y) : \text{same group}\}|}
\end{equation}

\item[Pairwise Grouping Recall ($R_p$)]
The fraction of alert pairs sharing the same root cause that are successfully grouped together:
\begin{equation}
R_p = \frac{|\{(a_x, a_y) : \text{same group} \land \text{same root cause}\}|}{|\{(a_x, a_y) : \text{same root cause}\}|}
\end{equation}

\item[Pairwise Grouping F1-Score ($F1_p$)]
The harmonic mean of pairwise precision and recall:
\begin{equation}
F1_p = 2 \times \frac{P_p \times R_p}{P_p + R_p}
\end{equation}
Target: $F1_p \ge 0.85$.

\item[Group Purity]
The average homogeneity of root causes within each generated group:
\begin{equation}
\text{Purity} = \frac{1}{N_{\text{groups}}} \sum_{g=1}^{N_{\text{groups}}} \frac{\max_{c} |\{a \in g : \text{root\_cause}(a) = c\}|}{|g|}
\end{equation}
Target: $\text{Purity} \ge 0.80$.

\item[Processing Latency]
The wall-clock time required to process 100 warning alerts through the aggregation pipeline, measured in milliseconds. For Arm 4, includes LLM inference time and token overhead.

% =============================================================================
% DATASET & BENCHMARK TERMS
% =============================================================================

\item[RCAEval]
A benchmark for root cause analysis in microservice systems comprising 735 failure cases across three systems (Online Boutique, Sock Shop, Train Ticket) with logs, metrics, traces, and root cause labels~\cite{pham2025rcaeval}. Licensed under MIT. Direct inspection shows that only \textbf{359 of the 735 cases ship raw log telemetry} (suite RE1 is metric-only) and that the log schema contains \texttt{timestamp}, \texttt{container\_name}, and \texttt{message} but \textbf{no severity column}.

\item[Severity Tier Assignment]
The three-tier procedure by which this work assigns the severity tier $\ell_i$, since RCAEval supplies none: \emph{Tier 1 (explicit)} parses a level token from the message text (Sock Shop, Train Ticket); \emph{Tier 2 (frequency-inferred)} flags a template as anomalous when its fault-window frequency exceeds its normal-window frequency by a ratio $R$ (Online Boutique); \emph{Tier 3 (ground-truth marker)} tags the \texttt{root\_cause.txt} indicator line as \texttt{ROOT\_CAUSE} for the eight annotated cases.

\item[LEMMA-RCA]
A large-scale multi-modal dataset for root cause analysis containing hundreds of microservice entities and four fault types (DDoS, storage failure, resource contention, noisy neighbor)~\cite{zheng2024lemma}. Licensed under CC-BY-NC-4.0 (non-commercial). Its log pipeline aggregates raw log lines into three-dimensional time series using golden-signal keywords that cover only \texttt{error}, \texttt{exception}, and \texttt{critical} --- notably \textbf{not} \texttt{warning} --- so it offers no independent warning tier.

\item[Fault Injection]
The controlled introduction of a failure (e.g., CPU stress, network latency, packet loss) into a microservice system at a known timestamp $T_{\text{inject}}$ to generate labeled failure cases for benchmarking.

\item[Normal Window]
The interval $[T_{\text{start}}, T_{\text{inject}})$ preceding fault injection, during which the system operates normally; alerts occurring here are candidates for \textit{Transient} noise. Note that the dataset fields \texttt{normal\_timesteps} and \texttt{faulty\_timesteps} are integer counts of sample points, not timestamp lists, and therefore define these windows only indirectly.

\item[Fault Window]
The interval $[T_{\text{inject}}, T_{\text{end}}]$ following fault injection, during which the system exhibits degradation; alerts here are labeled as \textit{Signal} (if from the root cause service) or \textit{Cascading} (if from a service reachable from the root cause service in the call graph).

\item[Ground Truth Root Cause Label]
The authoritative annotation of which service is the root cause of a failure case, provided by the RCAEval benchmark. Used to compute RCPR and pairwise metrics.

% =============================================================================
% BASELINE & INFRASTRUCTURE TERMS
% =============================================================================

\item[Prometheus Alertmanager]
The de facto standard alert routing and grouping component in the Cloud Native Computing Foundation (CNCF) ecosystem. Groups alerts by label matchers and time windows (\texttt{group\_wait}, \texttt{group\_interval}, \texttt{group\_by})~\cite{prometheus-alertmanager}.

\item[Group Wait (\texttt{group\_wait})]
The initial delay after the first alert in a group before sending a notification, allowing additional alerts to arrive and be grouped. Default: 30~seconds.

\item[Group Interval (\texttt{group\_interval})]
The minimum interval between subsequent notifications for the same group. Default: 5~minutes.

\item[Drain3]
An online log parsing algorithm that extracts log templates from raw log streams by clustering log messages based on token similarity, producing parameterized templates with variable placeholders~\cite{drain3}.

\item[Sentence-BERT (S-BERT)]
A modification of BERT that uses siamese and triplet network structures to derive semantically meaningful sentence embeddings, optimized for cosine similarity comparison~\cite{sbert}.

\item[HDBSCAN]
Hierarchical Density-Based Spatial Clustering of Applications with Noise; a density-based clustering algorithm that does not require specifying the number of clusters and can identify outliers (noise points)~\cite{hdbscan}.

% =============================================================================
% EXPERIMENTAL DESIGN TERMS
% =============================================================================

\item[Ablation Study]
An experimental procedure where individual components (e.g., \texttt{get\_service\_dependency} tool, fast filter, runbook) are removed from the full system to measure their individual contribution to performance metrics.

\item[Cross-System Generalization]
The ability of an aggregation method trained or tuned on one microservice system (e.g., Online Boutique) to maintain performance on unseen systems (e.g., Train Ticket) without retraining.

\item[Statistical Significance Test]
Paired Wilcoxon signed-rank test (non-parametric) used to compare metric distributions across the four arms on the same failure cases, with significance level $\alpha = 0.05$.

\item[Metric Gaming]
The phenomenon where a method optimizes a single metric (e.g., ARR) at the expense of operational utility (e.g., grouping all alerts into one group achieves ARR $\approx$ 100\% but RCPR = 0\%). Mitigated by the multi-metric framework (ARR, RCPR, F1, Purity).

% =============================================================================
% INDUSTRY STANDARDS & REFERENCES
% =============================================================================

\item[Google SRE Principles]
Foundational Site Reliability Engineering practices documented in \textit{Site Reliability Engineering} (Ch.~6, 10, 11) and \textit{The Site Reliability Workbook} (Ch.~5), emphasizing actionable alerts, maximum 2 incidents per 12-hour shift, and alerting on SLOs rather than raw metrics~\cite{google-sre-book, google-sre-workbook}.

\item[Observability Pillars]
The three canonical telemetry types: \textit{Logs} (immutable event records), \textit{Metrics} (aggregated numerical measurements), and \textit{Traces} (distributed request flow records)~\cite{grafana2025}.

\end{description}

% =============================================================================
% LOCAL BIBLIOGRAPHY MACROS (if not using main document's .bib)
% Include these in your main .bib or define via IEEEabrv.bib/IEEEfull.bib
% =============================================================================
% @string{srebook = "Site Reliability Engineering: How Google Runs Production Systems"}
% @string{sreworkbook = "The Site Reliability Workbook: Practical Ways to Implement SRE"}
% @article{pagerduty2020, ...}
% @article{grafana2025, ...}
% @inproceedings{kuang2024cola, ...}
% @inproceedings{pham2025rcaeval, ...}
% @article{zheng2024lemma, ...}
% @inproceedings{yu2025alertguardian, ...}
% @manual{prometheus-alertmanager, ...}
% @article{sbert, ...}
% @article{drain3, ...}
% @article{hdbscan, ...}
% @article{wei2022cot, ...} in main IEEEtran document
% Compiles with: pdflatex main.tex (where main.tex includes this file)
% =============================================================================

% -----------------------------------------------------------------------------
% USAGE IN MAIN DOCUMENT:
% \section*{Glossary of Terms}
% \addcontentsline{toc}{section}{Glossary of Terms}
% % =============================================================================
% GLOSSARY OF TERMS — Alert Aggregation in Microservices
% Standalone file for \input{glossary} in main IEEEtran document
% Compiles with: pdflatex main.tex (where main.tex includes this file)
% =============================================================================

% -----------------------------------------------------------------------------
% USAGE IN MAIN DOCUMENT:
% \section*{Glossary of Terms}
% \addcontentsline{toc}{section}{Glossary of Terms}
% \input{glossary}
% -----------------------------------------------------------------------------

\begin{description}

% =============================================================================
% CORE PROBLEM TERMS
% =============================================================================

\item[Alert Fatigue]
A cognitive state experienced by on-call engineers when the volume of alerts exceeds human cognitive capacity, leading to desensitization, missed critical alerts, and delayed incident response~\cite{pagerduty2020, grafana2025}.

\item[Alert Storm]
A sudden surge of alerts triggered by a single root cause propagating through dependent microservices, resulting in hundreds to thousands of alerts within minutes~\cite{kuang2024cola}.

\item[Noisy Alert Warning]
A warning-level alert generated by the monitoring system that is technically correct but does not require meaningful remediation action from an on-call engineer, classified into four categories: \textit{Duplicate}, \textit{Transient}, \textit{Unactionable}, and \textit{Cascading} (defined below).

\item[Early Warning Phase]
The pre-failure period during which a microservice exhibits performance degradation (e.g., increased latency, resource saturation) but has not yet violated Service Level Objectives (SLOs) or caused user-facing outages.

\item[Actionable Alert]
An alert that satisfies the Google SRE principle: ``every alert sent to an on-call engineer must require immediate human intervention''~\cite{google-sre-ch10, google-sre-ch11}.

% =============================================================================
% FOUR CATEGORIES OF NOISY ALERT WARNINGS
% =============================================================================

\item[Duplicate Warning Alert]
Multiple alerts originating from the same service with identical or near-identical message templates within a short time window ($\Delta t \le 60$~seconds), representing repeated observations of the same underlying condition rather than distinct incidents.

\item[Transient Warning Alert]
A warning alert that occurs during normal system operation (i.e., before fault injection time $T_{\text{inject}}$ in RCAEval) and resolves spontaneously without human intervention, typically caused by temporary network fluctuations or brief resource spikes.

\item[Unactionable Warning Alert]
A technically accurate warning alert for which no Standard Operating Procedure (SOP) or runbook exists, leaving the on-call engineer without a defined remediation path; the alert consumes attention but yields no operational value.

\item[Cascading Warning Alert]
A warning alert emitted by a downstream service that is a symptom of an upstream root cause; remediating the alerting service does not resolve the underlying issue, which originates in a different service (the root cause service).

% =============================================================================
% METHODOLOGY: FOUR EXPERIMENTAL ARMS
% =============================================================================

\item[Arm 1: Rule-Based Aggregation (Baseline)]
An alert grouping approach that clusters alerts based on static label matching (e.g., service name, alert name, severity) and fixed time windows (\texttt{group\_wait}, \texttt{group\_interval}), as implemented in Prometheus Alertmanager~\cite{prometheus-alertmanager}. Lacks awareness of service dependencies and semantic content.

\item[Arm 2: Semantic Similarity-Based Aggregation]
A method that parses raw logs into templates (via Drain3~\cite{drain3}), encodes templates into dense vector embeddings using Sentence-BERT~\cite{sbert}, and clusters alerts using density-based algorithms (HDBSCAN/DBSCAN) based on cosine similarity. Operates without topology knowledge.

\item[Arm 3: Temporal-Spatial Correlation Aggregation]
A topology-aware method that constructs a static service dependency graph (caller-callee adjacency matrix) from microservice architecture, then correlates alerts occurring within a sliding time window ($\Delta t \in \{30, 60, 120\}$~seconds) if a dependency path of length $k \le 2$ exists between the alerting services.

\item[Arm 4: LLM-Based Agentic Aggregation (Proposed)]
A hybrid two-stage approach: (1) a fast filter removes obvious duplicates and performs time-window pre-grouping; (2) an LLM agent (GPT-4o-mini / DeepSeek-V3) reasons over the remaining clusters using Chain-of-Thought prompting and three tools: \texttt{get\_related\_alerts}, \texttt{get\_service\_dependency}, and optionally \texttt{get\_runbook}. Limited to a maximum of two tool calls per grouping decision.

% =============================================================================
% TECHNICAL CONCEPTS & ARCHITECTURE
% =============================================================================

\item[Service Topology (Service Dependency Graph)]
A directed graph $G = (V, E)$ where vertices $V$ represent microservices and edges $E$ represent synchronous (HTTP/gRPC) or asynchronous (message queue) call relationships. Used to distinguish root cause services from downstream symptom services.

\item[Semantic Similarity]
A measure of meaning-based resemblance between two log messages, computed as the cosine similarity between their dense vector embeddings (typically 384-dimensional from \texttt{sentence-transformers/all-MiniLM-L6-v2}). Captures paraphrases and synonyms that lexical matching misses.

\item[Runbook Reasoning]
The process by which an LLM agent consults a synthetic runbook (a structured troubleshooting guide for a specific fault type) to determine whether an alert is actionable, requires escalation, or can be safely aggregated as noise.

\item[Tool-Calling Agent]
An LLM agent augmented with a predefined set of functions (tools) it can invoke during reasoning: \texttt{get\_related\_alerts(service, time\_window)}, \texttt{get\_service\_dependency(service)}, and \texttt{get\_runbook(fault\_type)}. Tool outputs are fed back into the agent's context for subsequent reasoning steps.

\item[Chain-of-Thought (CoT) Prompting]
A prompting technique that instructs the LLM to generate intermediate reasoning steps before producing a final answer, improving accuracy on multi-step causal inference tasks~\cite{wei2022cot}.

\item[Fast Filter (Deduplication + Time-Window Pre-Grouping)]
A lightweight deterministic preprocessing stage that removes exact duplicate alerts and performs initial grouping by service and time window, reducing the alert volume passed to the LLM agent by an estimated 70--80\%.

\item[Synthetic Runbook]
A researcher-authored troubleshooting document for each fault type in the benchmark dataset (e.g., CPU stress, memory leak, network latency), containing symptoms, diagnosis steps, and remediation actions. Not a production SOP; used solely for experimental evaluation.

% =============================================================================
% EVALUATION METRICS
% =============================================================================

\item[Alert Reduction Ratio (ARR)]
The primary metric quantifying alert volume reduction:
\begin{equation}
\text{ARR} = \left(1 - \frac{N_{\text{groups}}}{N_{\text{raw\_alerts}}}\right) \times 100\%
\end{equation}
where $N_{\text{raw\_alerts}}$ is the number of warning alerts before aggregation and $N_{\text{groups}}$ is the number of aggregated groups presented to the engineer. Target: $\text{ARR} \ge 60\%$.

\item[Root Cause Preservation Rate (RCPR)]
The critical safety metric ensuring the root cause evidence is not lost during aggregation. It is defined at two granularities. At service granularity, which is the primary metric and spans every case that exposes logs ($M = 359$):
\begin{equation}
\text{RCPR}_{\text{svc}} = \frac{1}{M} \sum_{i=1}^{M} \mathbb{I}(s_{\text{root}}^{(i)} \in \text{PrimaryGroup}_i) \times 100\%
\end{equation}
At alert granularity, a secondary metric restricted to the $M'$ cases whose ground-truth \texttt{root\_cause.txt} indicator line is itself a warning-tier log ($M' = 4$):
\begin{equation}
\text{RCPR}_{\text{alert}} = \frac{1}{M'} \sum_{i=1}^{M'} \mathbb{I}(\text{RootCauseAlert}_i \in \text{PrimaryGroup}_i) \times 100\%
\end{equation}
where $\mathbb{I}$ is the indicator function. Target: $\text{RCPR} \ge 95\%$ (mandatory).

\item[Pairwise Grouping Precision ($P_p$)]
The fraction of alert pairs placed in the same group that truly share the same root cause:
\begin{equation}
P_p = \frac{|\{(a_x, a_y) : \text{same group} \land \text{same root cause}\}|}{|\{(a_x, a_y) : \text{same group}\}|}
\end{equation}

\item[Pairwise Grouping Recall ($R_p$)]
The fraction of alert pairs sharing the same root cause that are successfully grouped together:
\begin{equation}
R_p = \frac{|\{(a_x, a_y) : \text{same group} \land \text{same root cause}\}|}{|\{(a_x, a_y) : \text{same root cause}\}|}
\end{equation}

\item[Pairwise Grouping F1-Score ($F1_p$)]
The harmonic mean of pairwise precision and recall:
\begin{equation}
F1_p = 2 \times \frac{P_p \times R_p}{P_p + R_p}
\end{equation}
Target: $F1_p \ge 0.85$.

\item[Group Purity]
The average homogeneity of root causes within each generated group:
\begin{equation}
\text{Purity} = \frac{1}{N_{\text{groups}}} \sum_{g=1}^{N_{\text{groups}}} \frac{\max_{c} |\{a \in g : \text{root\_cause}(a) = c\}|}{|g|}
\end{equation}
Target: $\text{Purity} \ge 0.80$.

\item[Processing Latency]
The wall-clock time required to process 100 warning alerts through the aggregation pipeline, measured in milliseconds. For Arm 4, includes LLM inference time and token overhead.

% =============================================================================
% DATASET & BENCHMARK TERMS
% =============================================================================

\item[RCAEval]
A benchmark for root cause analysis in microservice systems comprising 735 failure cases across three systems (Online Boutique, Sock Shop, Train Ticket) with logs, metrics, traces, and root cause labels~\cite{pham2025rcaeval}. Licensed under MIT. Direct inspection shows that only \textbf{359 of the 735 cases ship raw log telemetry} (suite RE1 is metric-only) and that the log schema contains \texttt{timestamp}, \texttt{container\_name}, and \texttt{message} but \textbf{no severity column}.

\item[Severity Tier Assignment]
The three-tier procedure by which this work assigns the severity tier $\ell_i$, since RCAEval supplies none: \emph{Tier 1 (explicit)} parses a level token from the message text (Sock Shop, Train Ticket); \emph{Tier 2 (frequency-inferred)} flags a template as anomalous when its fault-window frequency exceeds its normal-window frequency by a ratio $R$ (Online Boutique); \emph{Tier 3 (ground-truth marker)} tags the \texttt{root\_cause.txt} indicator line as \texttt{ROOT\_CAUSE} for the eight annotated cases.

\item[LEMMA-RCA]
A large-scale multi-modal dataset for root cause analysis containing hundreds of microservice entities and four fault types (DDoS, storage failure, resource contention, noisy neighbor)~\cite{zheng2024lemma}. Licensed under CC-BY-NC-4.0 (non-commercial). Its log pipeline aggregates raw log lines into three-dimensional time series using golden-signal keywords that cover only \texttt{error}, \texttt{exception}, and \texttt{critical} --- notably \textbf{not} \texttt{warning} --- so it offers no independent warning tier.

\item[Fault Injection]
The controlled introduction of a failure (e.g., CPU stress, network latency, packet loss) into a microservice system at a known timestamp $T_{\text{inject}}$ to generate labeled failure cases for benchmarking.

\item[Normal Window]
The interval $[T_{\text{start}}, T_{\text{inject}})$ preceding fault injection, during which the system operates normally; alerts occurring here are candidates for \textit{Transient} noise. Note that the dataset fields \texttt{normal\_timesteps} and \texttt{faulty\_timesteps} are integer counts of sample points, not timestamp lists, and therefore define these windows only indirectly.

\item[Fault Window]
The interval $[T_{\text{inject}}, T_{\text{end}}]$ following fault injection, during which the system exhibits degradation; alerts here are labeled as \textit{Signal} (if from the root cause service) or \textit{Cascading} (if from a service reachable from the root cause service in the call graph).

\item[Ground Truth Root Cause Label]
The authoritative annotation of which service is the root cause of a failure case, provided by the RCAEval benchmark. Used to compute RCPR and pairwise metrics.

% =============================================================================
% BASELINE & INFRASTRUCTURE TERMS
% =============================================================================

\item[Prometheus Alertmanager]
The de facto standard alert routing and grouping component in the Cloud Native Computing Foundation (CNCF) ecosystem. Groups alerts by label matchers and time windows (\texttt{group\_wait}, \texttt{group\_interval}, \texttt{group\_by})~\cite{prometheus-alertmanager}.

\item[Group Wait (\texttt{group\_wait})]
The initial delay after the first alert in a group before sending a notification, allowing additional alerts to arrive and be grouped. Default: 30~seconds.

\item[Group Interval (\texttt{group\_interval})]
The minimum interval between subsequent notifications for the same group. Default: 5~minutes.

\item[Drain3]
An online log parsing algorithm that extracts log templates from raw log streams by clustering log messages based on token similarity, producing parameterized templates with variable placeholders~\cite{drain3}.

\item[Sentence-BERT (S-BERT)]
A modification of BERT that uses siamese and triplet network structures to derive semantically meaningful sentence embeddings, optimized for cosine similarity comparison~\cite{sbert}.

\item[HDBSCAN]
Hierarchical Density-Based Spatial Clustering of Applications with Noise; a density-based clustering algorithm that does not require specifying the number of clusters and can identify outliers (noise points)~\cite{hdbscan}.

% =============================================================================
% EXPERIMENTAL DESIGN TERMS
% =============================================================================

\item[Ablation Study]
An experimental procedure where individual components (e.g., \texttt{get\_service\_dependency} tool, fast filter, runbook) are removed from the full system to measure their individual contribution to performance metrics.

\item[Cross-System Generalization]
The ability of an aggregation method trained or tuned on one microservice system (e.g., Online Boutique) to maintain performance on unseen systems (e.g., Train Ticket) without retraining.

\item[Statistical Significance Test]
Paired Wilcoxon signed-rank test (non-parametric) used to compare metric distributions across the four arms on the same failure cases, with significance level $\alpha = 0.05$.

\item[Metric Gaming]
The phenomenon where a method optimizes a single metric (e.g., ARR) at the expense of operational utility (e.g., grouping all alerts into one group achieves ARR $\approx$ 100\% but RCPR = 0\%). Mitigated by the multi-metric framework (ARR, RCPR, F1, Purity).

% =============================================================================
% INDUSTRY STANDARDS & REFERENCES
% =============================================================================

\item[Google SRE Principles]
Foundational Site Reliability Engineering practices documented in \textit{Site Reliability Engineering} (Ch.~6, 10, 11) and \textit{The Site Reliability Workbook} (Ch.~5), emphasizing actionable alerts, maximum 2 incidents per 12-hour shift, and alerting on SLOs rather than raw metrics~\cite{google-sre-book, google-sre-workbook}.

\item[Observability Pillars]
The three canonical telemetry types: \textit{Logs} (immutable event records), \textit{Metrics} (aggregated numerical measurements), and \textit{Traces} (distributed request flow records)~\cite{grafana2025}.

\end{description}

% =============================================================================
% LOCAL BIBLIOGRAPHY MACROS (if not using main document's .bib)
% Include these in your main .bib or define via IEEEabrv.bib/IEEEfull.bib
% =============================================================================
% @string{srebook = "Site Reliability Engineering: How Google Runs Production Systems"}
% @string{sreworkbook = "The Site Reliability Workbook: Practical Ways to Implement SRE"}
% @article{pagerduty2020, ...}
% @article{grafana2025, ...}
% @inproceedings{kuang2024cola, ...}
% @inproceedings{pham2025rcaeval, ...}
% @article{zheng2024lemma, ...}
% @inproceedings{yu2025alertguardian, ...}
% @manual{prometheus-alertmanager, ...}
% @article{sbert, ...}
% @article{drain3, ...}
% @article{hdbscan, ...}
% @article{wei2022cot, ...}
% -----------------------------------------------------------------------------

\begin{description}

% =============================================================================
% CORE PROBLEM TERMS
% =============================================================================

\item[Alert Fatigue]
A cognitive state experienced by on-call engineers when the volume of alerts exceeds human cognitive capacity, leading to desensitization, missed critical alerts, and delayed incident response~\cite{pagerduty2020, grafana2025}.

\item[Alert Storm]
A sudden surge of alerts triggered by a single root cause propagating through dependent microservices, resulting in hundreds to thousands of alerts within minutes~\cite{kuang2024cola}.

\item[Noisy Alert Warning]
A warning-level alert generated by the monitoring system that is technically correct but does not require meaningful remediation action from an on-call engineer, classified into four categories: \textit{Duplicate}, \textit{Transient}, \textit{Unactionable}, and \textit{Cascading} (defined below).

\item[Early Warning Phase]
The pre-failure period during which a microservice exhibits performance degradation (e.g., increased latency, resource saturation) but has not yet violated Service Level Objectives (SLOs) or caused user-facing outages.

\item[Actionable Alert]
An alert that satisfies the Google SRE principle: ``every alert sent to an on-call engineer must require immediate human intervention''~\cite{google-sre-ch10, google-sre-ch11}.

% =============================================================================
% FOUR CATEGORIES OF NOISY ALERT WARNINGS
% =============================================================================

\item[Duplicate Warning Alert]
Multiple alerts originating from the same service with identical or near-identical message templates within a short time window ($\Delta t \le 60$~seconds), representing repeated observations of the same underlying condition rather than distinct incidents.

\item[Transient Warning Alert]
A warning alert that occurs during normal system operation (i.e., before fault injection time $T_{\text{inject}}$ in RCAEval) and resolves spontaneously without human intervention, typically caused by temporary network fluctuations or brief resource spikes.

\item[Unactionable Warning Alert]
A technically accurate warning alert for which no Standard Operating Procedure (SOP) or runbook exists, leaving the on-call engineer without a defined remediation path; the alert consumes attention but yields no operational value.

\item[Cascading Warning Alert]
A warning alert emitted by a downstream service that is a symptom of an upstream root cause; remediating the alerting service does not resolve the underlying issue, which originates in a different service (the root cause service).

% =============================================================================
% METHODOLOGY: FOUR EXPERIMENTAL ARMS
% =============================================================================

\item[Arm 1: Rule-Based Aggregation (Baseline)]
An alert grouping approach that clusters alerts based on static label matching (e.g., service name, alert name, severity) and fixed time windows (\texttt{group\_wait}, \texttt{group\_interval}), as implemented in Prometheus Alertmanager~\cite{prometheus-alertmanager}. Lacks awareness of service dependencies and semantic content.

\item[Arm 2: Semantic Similarity-Based Aggregation]
A method that parses raw logs into templates (via Drain3~\cite{drain3}), encodes templates into dense vector embeddings using Sentence-BERT~\cite{sbert}, and clusters alerts using density-based algorithms (HDBSCAN/DBSCAN) based on cosine similarity. Operates without topology knowledge.

\item[Arm 3: Temporal-Spatial Correlation Aggregation]
A topology-aware method that constructs a static service dependency graph (caller-callee adjacency matrix) from microservice architecture, then correlates alerts occurring within a sliding time window ($\Delta t \in \{30, 60, 120\}$~seconds) if a dependency path of length $k \le 2$ exists between the alerting services.

\item[Arm 4: LLM-Based Agentic Aggregation (Proposed)]
A hybrid two-stage approach: (1) a fast filter removes obvious duplicates and performs time-window pre-grouping; (2) an LLM agent (GPT-4o-mini / DeepSeek-V3) reasons over the remaining clusters using Chain-of-Thought prompting and three tools: \texttt{get\_related\_alerts}, \texttt{get\_service\_dependency}, and optionally \texttt{get\_runbook}. Limited to a maximum of two tool calls per grouping decision.

% =============================================================================
% TECHNICAL CONCEPTS & ARCHITECTURE
% =============================================================================

\item[Service Topology (Service Dependency Graph)]
A directed graph $G = (V, E)$ where vertices $V$ represent microservices and edges $E$ represent synchronous (HTTP/gRPC) or asynchronous (message queue) call relationships. Used to distinguish root cause services from downstream symptom services.

\item[Semantic Similarity]
A measure of meaning-based resemblance between two log messages, computed as the cosine similarity between their dense vector embeddings (typically 384-dimensional from \texttt{sentence-transformers/all-MiniLM-L6-v2}). Captures paraphrases and synonyms that lexical matching misses.

\item[Runbook Reasoning]
The process by which an LLM agent consults a synthetic runbook (a structured troubleshooting guide for a specific fault type) to determine whether an alert is actionable, requires escalation, or can be safely aggregated as noise.

\item[Tool-Calling Agent]
An LLM agent augmented with a predefined set of functions (tools) it can invoke during reasoning: \texttt{get\_related\_alerts(service, time\_window)}, \texttt{get\_service\_dependency(service)}, and \texttt{get\_runbook(fault\_type)}. Tool outputs are fed back into the agent's context for subsequent reasoning steps.

\item[Chain-of-Thought (CoT) Prompting]
A prompting technique that instructs the LLM to generate intermediate reasoning steps before producing a final answer, improving accuracy on multi-step causal inference tasks~\cite{wei2022cot}.

\item[Fast Filter (Deduplication + Time-Window Pre-Grouping)]
A lightweight deterministic preprocessing stage that removes exact duplicate alerts and performs initial grouping by service and time window, reducing the alert volume passed to the LLM agent by an estimated 70--80\%.

\item[Synthetic Runbook]
A researcher-authored troubleshooting document for each fault type in the benchmark dataset (e.g., CPU stress, memory leak, network latency), containing symptoms, diagnosis steps, and remediation actions. Not a production SOP; used solely for experimental evaluation.

% =============================================================================
% EVALUATION METRICS
% =============================================================================

\item[Alert Reduction Ratio (ARR)]
The primary metric quantifying alert volume reduction:
\begin{equation}
\text{ARR} = \left(1 - \frac{N_{\text{groups}}}{N_{\text{raw\_alerts}}}\right) \times 100\%
\end{equation}
where $N_{\text{raw\_alerts}}$ is the number of warning alerts before aggregation and $N_{\text{groups}}$ is the number of aggregated groups presented to the engineer. Target: $\text{ARR} \ge 60\%$.

\item[Root Cause Preservation Rate (RCPR)]
The critical safety metric ensuring the root cause evidence is not lost during aggregation. It is defined at two granularities. At service granularity, which is the primary metric and spans every case that exposes logs ($M = 359$):
\begin{equation}
\text{RCPR}_{\text{svc}} = \frac{1}{M} \sum_{i=1}^{M} \mathbb{I}(s_{\text{root}}^{(i)} \in \text{PrimaryGroup}_i) \times 100\%
\end{equation}
At alert granularity, a secondary metric restricted to the $M'$ cases whose ground-truth \texttt{root\_cause.txt} indicator line is itself a warning-tier log ($M' = 4$):
\begin{equation}
\text{RCPR}_{\text{alert}} = \frac{1}{M'} \sum_{i=1}^{M'} \mathbb{I}(\text{RootCauseAlert}_i \in \text{PrimaryGroup}_i) \times 100\%
\end{equation}
where $\mathbb{I}$ is the indicator function. Target: $\text{RCPR} \ge 95\%$ (mandatory).

\item[Pairwise Grouping Precision ($P_p$)]
The fraction of alert pairs placed in the same group that truly share the same root cause:
\begin{equation}
P_p = \frac{|\{(a_x, a_y) : \text{same group} \land \text{same root cause}\}|}{|\{(a_x, a_y) : \text{same group}\}|}
\end{equation}

\item[Pairwise Grouping Recall ($R_p$)]
The fraction of alert pairs sharing the same root cause that are successfully grouped together:
\begin{equation}
R_p = \frac{|\{(a_x, a_y) : \text{same group} \land \text{same root cause}\}|}{|\{(a_x, a_y) : \text{same root cause}\}|}
\end{equation}

\item[Pairwise Grouping F1-Score ($F1_p$)]
The harmonic mean of pairwise precision and recall:
\begin{equation}
F1_p = 2 \times \frac{P_p \times R_p}{P_p + R_p}
\end{equation}
Target: $F1_p \ge 0.85$.

\item[Group Purity]
The average homogeneity of root causes within each generated group:
\begin{equation}
\text{Purity} = \frac{1}{N_{\text{groups}}} \sum_{g=1}^{N_{\text{groups}}} \frac{\max_{c} |\{a \in g : \text{root\_cause}(a) = c\}|}{|g|}
\end{equation}
Target: $\text{Purity} \ge 0.80$.

\item[Processing Latency]
The wall-clock time required to process 100 warning alerts through the aggregation pipeline, measured in milliseconds. For Arm 4, includes LLM inference time and token overhead.

% =============================================================================
% DATASET & BENCHMARK TERMS
% =============================================================================

\item[RCAEval]
A benchmark for root cause analysis in microservice systems comprising 735 failure cases across three systems (Online Boutique, Sock Shop, Train Ticket) with logs, metrics, traces, and root cause labels~\cite{pham2025rcaeval}. Licensed under MIT. Direct inspection shows that only \textbf{359 of the 735 cases ship raw log telemetry} (suite RE1 is metric-only) and that the log schema contains \texttt{timestamp}, \texttt{container\_name}, and \texttt{message} but \textbf{no severity column}.

\item[Severity Tier Assignment]
The three-tier procedure by which this work assigns the severity tier $\ell_i$, since RCAEval supplies none: \emph{Tier 1 (explicit)} parses a level token from the message text (Sock Shop, Train Ticket); \emph{Tier 2 (frequency-inferred)} flags a template as anomalous when its fault-window frequency exceeds its normal-window frequency by a ratio $R$ (Online Boutique); \emph{Tier 3 (ground-truth marker)} tags the \texttt{root\_cause.txt} indicator line as \texttt{ROOT\_CAUSE} for the eight annotated cases.

\item[LEMMA-RCA]
A large-scale multi-modal dataset for root cause analysis containing hundreds of microservice entities and four fault types (DDoS, storage failure, resource contention, noisy neighbor)~\cite{zheng2024lemma}. Licensed under CC-BY-NC-4.0 (non-commercial). Its log pipeline aggregates raw log lines into three-dimensional time series using golden-signal keywords that cover only \texttt{error}, \texttt{exception}, and \texttt{critical} --- notably \textbf{not} \texttt{warning} --- so it offers no independent warning tier.

\item[Fault Injection]
The controlled introduction of a failure (e.g., CPU stress, network latency, packet loss) into a microservice system at a known timestamp $T_{\text{inject}}$ to generate labeled failure cases for benchmarking.

\item[Normal Window]
The interval $[T_{\text{start}}, T_{\text{inject}})$ preceding fault injection, during which the system operates normally; alerts occurring here are candidates for \textit{Transient} noise. Note that the dataset fields \texttt{normal\_timesteps} and \texttt{faulty\_timesteps} are integer counts of sample points, not timestamp lists, and therefore define these windows only indirectly.

\item[Fault Window]
The interval $[T_{\text{inject}}, T_{\text{end}}]$ following fault injection, during which the system exhibits degradation; alerts here are labeled as \textit{Signal} (if from the root cause service) or \textit{Cascading} (if from a service reachable from the root cause service in the call graph).

\item[Ground Truth Root Cause Label]
The authoritative annotation of which service is the root cause of a failure case, provided by the RCAEval benchmark. Used to compute RCPR and pairwise metrics.

% =============================================================================
% BASELINE & INFRASTRUCTURE TERMS
% =============================================================================

\item[Prometheus Alertmanager]
The de facto standard alert routing and grouping component in the Cloud Native Computing Foundation (CNCF) ecosystem. Groups alerts by label matchers and time windows (\texttt{group\_wait}, \texttt{group\_interval}, \texttt{group\_by})~\cite{prometheus-alertmanager}.

\item[Group Wait (\texttt{group\_wait})]
The initial delay after the first alert in a group before sending a notification, allowing additional alerts to arrive and be grouped. Default: 30~seconds.

\item[Group Interval (\texttt{group\_interval})]
The minimum interval between subsequent notifications for the same group. Default: 5~minutes.

\item[Drain3]
An online log parsing algorithm that extracts log templates from raw log streams by clustering log messages based on token similarity, producing parameterized templates with variable placeholders~\cite{drain3}.

\item[Sentence-BERT (S-BERT)]
A modification of BERT that uses siamese and triplet network structures to derive semantically meaningful sentence embeddings, optimized for cosine similarity comparison~\cite{sbert}.

\item[HDBSCAN]
Hierarchical Density-Based Spatial Clustering of Applications with Noise; a density-based clustering algorithm that does not require specifying the number of clusters and can identify outliers (noise points)~\cite{hdbscan}.

% =============================================================================
% EXPERIMENTAL DESIGN TERMS
% =============================================================================

\item[Ablation Study]
An experimental procedure where individual components (e.g., \texttt{get\_service\_dependency} tool, fast filter, runbook) are removed from the full system to measure their individual contribution to performance metrics.

\item[Cross-System Generalization]
The ability of an aggregation method trained or tuned on one microservice system (e.g., Online Boutique) to maintain performance on unseen systems (e.g., Train Ticket) without retraining.

\item[Statistical Significance Test]
Paired Wilcoxon signed-rank test (non-parametric) used to compare metric distributions across the four arms on the same failure cases, with significance level $\alpha = 0.05$.

\item[Metric Gaming]
The phenomenon where a method optimizes a single metric (e.g., ARR) at the expense of operational utility (e.g., grouping all alerts into one group achieves ARR $\approx$ 100\% but RCPR = 0\%). Mitigated by the multi-metric framework (ARR, RCPR, F1, Purity).

% =============================================================================
% INDUSTRY STANDARDS & REFERENCES
% =============================================================================

\item[Google SRE Principles]
Foundational Site Reliability Engineering practices documented in \textit{Site Reliability Engineering} (Ch.~6, 10, 11) and \textit{The Site Reliability Workbook} (Ch.~5), emphasizing actionable alerts, maximum 2 incidents per 12-hour shift, and alerting on SLOs rather than raw metrics~\cite{google-sre-book, google-sre-workbook}.

\item[Observability Pillars]
The three canonical telemetry types: \textit{Logs} (immutable event records), \textit{Metrics} (aggregated numerical measurements), and \textit{Traces} (distributed request flow records)~\cite{grafana2025}.

\end{description}

% =============================================================================
% LOCAL BIBLIOGRAPHY MACROS (if not using main document's .bib)
% Include these in your main .bib or define via IEEEabrv.bib/IEEEfull.bib
% =============================================================================
% @string{srebook = "Site Reliability Engineering: How Google Runs Production Systems"}
% @string{sreworkbook = "The Site Reliability Workbook: Practical Ways to Implement SRE"}
% @article{pagerduty2020, ...}
% @article{grafana2025, ...}
% @inproceedings{kuang2024cola, ...}
% @inproceedings{pham2025rcaeval, ...}
% @article{zheng2024lemma, ...}
% @inproceedings{yu2025alertguardian, ...}
% @manual{prometheus-alertmanager, ...}
% @article{sbert, ...}
% @article{drain3, ...}
% @article{hdbscan, ...}
% @article{wei2022cot, ...} in main IEEEtran document
% Compiles with: pdflatex main.tex (where main.tex includes this file)
% =============================================================================

% -----------------------------------------------------------------------------
% USAGE IN MAIN DOCUMENT:
% \section*{Glossary of Terms}
% \addcontentsline{toc}{section}{Glossary of Terms}
% % =============================================================================
% GLOSSARY OF TERMS — Alert Aggregation in Microservices
% Standalone file for % =============================================================================
% GLOSSARY OF TERMS — Alert Aggregation in Microservices
% Standalone file for \input{glossary} in main IEEEtran document
% Compiles with: pdflatex main.tex (where main.tex includes this file)
% =============================================================================

% -----------------------------------------------------------------------------
% USAGE IN MAIN DOCUMENT:
% \section*{Glossary of Terms}
% \addcontentsline{toc}{section}{Glossary of Terms}
% \input{glossary}
% -----------------------------------------------------------------------------

\begin{description}

% =============================================================================
% CORE PROBLEM TERMS
% =============================================================================

\item[Alert Fatigue]
A cognitive state experienced by on-call engineers when the volume of alerts exceeds human cognitive capacity, leading to desensitization, missed critical alerts, and delayed incident response~\cite{pagerduty2020, grafana2025}.

\item[Alert Storm]
A sudden surge of alerts triggered by a single root cause propagating through dependent microservices, resulting in hundreds to thousands of alerts within minutes~\cite{kuang2024cola}.

\item[Noisy Alert Warning]
A warning-level alert generated by the monitoring system that is technically correct but does not require meaningful remediation action from an on-call engineer, classified into four categories: \textit{Duplicate}, \textit{Transient}, \textit{Unactionable}, and \textit{Cascading} (defined below).

\item[Early Warning Phase]
The pre-failure period during which a microservice exhibits performance degradation (e.g., increased latency, resource saturation) but has not yet violated Service Level Objectives (SLOs) or caused user-facing outages.

\item[Actionable Alert]
An alert that satisfies the Google SRE principle: ``every alert sent to an on-call engineer must require immediate human intervention''~\cite{google-sre-ch10, google-sre-ch11}.

% =============================================================================
% FOUR CATEGORIES OF NOISY ALERT WARNINGS
% =============================================================================

\item[Duplicate Warning Alert]
Multiple alerts originating from the same service with identical or near-identical message templates within a short time window ($\Delta t \le 60$~seconds), representing repeated observations of the same underlying condition rather than distinct incidents.

\item[Transient Warning Alert]
A warning alert that occurs during normal system operation (i.e., before fault injection time $T_{\text{inject}}$ in RCAEval) and resolves spontaneously without human intervention, typically caused by temporary network fluctuations or brief resource spikes.

\item[Unactionable Warning Alert]
A technically accurate warning alert for which no Standard Operating Procedure (SOP) or runbook exists, leaving the on-call engineer without a defined remediation path; the alert consumes attention but yields no operational value.

\item[Cascading Warning Alert]
A warning alert emitted by a downstream service that is a symptom of an upstream root cause; remediating the alerting service does not resolve the underlying issue, which originates in a different service (the root cause service).

% =============================================================================
% METHODOLOGY: FOUR EXPERIMENTAL ARMS
% =============================================================================

\item[Arm 1: Rule-Based Aggregation (Baseline)]
An alert grouping approach that clusters alerts based on static label matching (e.g., service name, alert name, severity) and fixed time windows (\texttt{group\_wait}, \texttt{group\_interval}), as implemented in Prometheus Alertmanager~\cite{prometheus-alertmanager}. Lacks awareness of service dependencies and semantic content.

\item[Arm 2: Semantic Similarity-Based Aggregation]
A method that parses raw logs into templates (via Drain3~\cite{drain3}), encodes templates into dense vector embeddings using Sentence-BERT~\cite{sbert}, and clusters alerts using density-based algorithms (HDBSCAN/DBSCAN) based on cosine similarity. Operates without topology knowledge.

\item[Arm 3: Temporal-Spatial Correlation Aggregation]
A topology-aware method that constructs a static service dependency graph (caller-callee adjacency matrix) from microservice architecture, then correlates alerts occurring within a sliding time window ($\Delta t \in \{30, 60, 120\}$~seconds) if a dependency path of length $k \le 2$ exists between the alerting services.

\item[Arm 4: LLM-Based Agentic Aggregation (Proposed)]
A hybrid two-stage approach: (1) a fast filter removes obvious duplicates and performs time-window pre-grouping; (2) an LLM agent (GPT-4o-mini / DeepSeek-V3) reasons over the remaining clusters using Chain-of-Thought prompting and three tools: \texttt{get\_related\_alerts}, \texttt{get\_service\_dependency}, and optionally \texttt{get\_runbook}. Limited to a maximum of two tool calls per grouping decision.

% =============================================================================
% TECHNICAL CONCEPTS & ARCHITECTURE
% =============================================================================

\item[Service Topology (Service Dependency Graph)]
A directed graph $G = (V, E)$ where vertices $V$ represent microservices and edges $E$ represent synchronous (HTTP/gRPC) or asynchronous (message queue) call relationships. Used to distinguish root cause services from downstream symptom services.

\item[Semantic Similarity]
A measure of meaning-based resemblance between two log messages, computed as the cosine similarity between their dense vector embeddings (typically 384-dimensional from \texttt{sentence-transformers/all-MiniLM-L6-v2}). Captures paraphrases and synonyms that lexical matching misses.

\item[Runbook Reasoning]
The process by which an LLM agent consults a synthetic runbook (a structured troubleshooting guide for a specific fault type) to determine whether an alert is actionable, requires escalation, or can be safely aggregated as noise.

\item[Tool-Calling Agent]
An LLM agent augmented with a predefined set of functions (tools) it can invoke during reasoning: \texttt{get\_related\_alerts(service, time\_window)}, \texttt{get\_service\_dependency(service)}, and \texttt{get\_runbook(fault\_type)}. Tool outputs are fed back into the agent's context for subsequent reasoning steps.

\item[Chain-of-Thought (CoT) Prompting]
A prompting technique that instructs the LLM to generate intermediate reasoning steps before producing a final answer, improving accuracy on multi-step causal inference tasks~\cite{wei2022cot}.

\item[Fast Filter (Deduplication + Time-Window Pre-Grouping)]
A lightweight deterministic preprocessing stage that removes exact duplicate alerts and performs initial grouping by service and time window, reducing the alert volume passed to the LLM agent by an estimated 70--80\%.

\item[Synthetic Runbook]
A researcher-authored troubleshooting document for each fault type in the benchmark dataset (e.g., CPU stress, memory leak, network latency), containing symptoms, diagnosis steps, and remediation actions. Not a production SOP; used solely for experimental evaluation.

% =============================================================================
% EVALUATION METRICS
% =============================================================================

\item[Alert Reduction Ratio (ARR)]
The primary metric quantifying alert volume reduction:
\begin{equation}
\text{ARR} = \left(1 - \frac{N_{\text{groups}}}{N_{\text{raw\_alerts}}}\right) \times 100\%
\end{equation}
where $N_{\text{raw\_alerts}}$ is the number of warning alerts before aggregation and $N_{\text{groups}}$ is the number of aggregated groups presented to the engineer. Target: $\text{ARR} \ge 60\%$.

\item[Root Cause Preservation Rate (RCPR)]
The critical safety metric ensuring the root cause evidence is not lost during aggregation. It is defined at two granularities. At service granularity, which is the primary metric and spans every case that exposes logs ($M = 359$):
\begin{equation}
\text{RCPR}_{\text{svc}} = \frac{1}{M} \sum_{i=1}^{M} \mathbb{I}(s_{\text{root}}^{(i)} \in \text{PrimaryGroup}_i) \times 100\%
\end{equation}
At alert granularity, a secondary metric restricted to the $M'$ cases whose ground-truth \texttt{root\_cause.txt} indicator line is itself a warning-tier log ($M' = 4$):
\begin{equation}
\text{RCPR}_{\text{alert}} = \frac{1}{M'} \sum_{i=1}^{M'} \mathbb{I}(\text{RootCauseAlert}_i \in \text{PrimaryGroup}_i) \times 100\%
\end{equation}
where $\mathbb{I}$ is the indicator function. Target: $\text{RCPR} \ge 95\%$ (mandatory).

\item[Pairwise Grouping Precision ($P_p$)]
The fraction of alert pairs placed in the same group that truly share the same root cause:
\begin{equation}
P_p = \frac{|\{(a_x, a_y) : \text{same group} \land \text{same root cause}\}|}{|\{(a_x, a_y) : \text{same group}\}|}
\end{equation}

\item[Pairwise Grouping Recall ($R_p$)]
The fraction of alert pairs sharing the same root cause that are successfully grouped together:
\begin{equation}
R_p = \frac{|\{(a_x, a_y) : \text{same group} \land \text{same root cause}\}|}{|\{(a_x, a_y) : \text{same root cause}\}|}
\end{equation}

\item[Pairwise Grouping F1-Score ($F1_p$)]
The harmonic mean of pairwise precision and recall:
\begin{equation}
F1_p = 2 \times \frac{P_p \times R_p}{P_p + R_p}
\end{equation}
Target: $F1_p \ge 0.85$.

\item[Group Purity]
The average homogeneity of root causes within each generated group:
\begin{equation}
\text{Purity} = \frac{1}{N_{\text{groups}}} \sum_{g=1}^{N_{\text{groups}}} \frac{\max_{c} |\{a \in g : \text{root\_cause}(a) = c\}|}{|g|}
\end{equation}
Target: $\text{Purity} \ge 0.80$.

\item[Processing Latency]
The wall-clock time required to process 100 warning alerts through the aggregation pipeline, measured in milliseconds. For Arm 4, includes LLM inference time and token overhead.

% =============================================================================
% DATASET & BENCHMARK TERMS
% =============================================================================

\item[RCAEval]
A benchmark for root cause analysis in microservice systems comprising 735 failure cases across three systems (Online Boutique, Sock Shop, Train Ticket) with logs, metrics, traces, and root cause labels~\cite{pham2025rcaeval}. Licensed under MIT. Direct inspection shows that only \textbf{359 of the 735 cases ship raw log telemetry} (suite RE1 is metric-only) and that the log schema contains \texttt{timestamp}, \texttt{container\_name}, and \texttt{message} but \textbf{no severity column}.

\item[Severity Tier Assignment]
The three-tier procedure by which this work assigns the severity tier $\ell_i$, since RCAEval supplies none: \emph{Tier 1 (explicit)} parses a level token from the message text (Sock Shop, Train Ticket); \emph{Tier 2 (frequency-inferred)} flags a template as anomalous when its fault-window frequency exceeds its normal-window frequency by a ratio $R$ (Online Boutique); \emph{Tier 3 (ground-truth marker)} tags the \texttt{root\_cause.txt} indicator line as \texttt{ROOT\_CAUSE} for the eight annotated cases.

\item[LEMMA-RCA]
A large-scale multi-modal dataset for root cause analysis containing hundreds of microservice entities and four fault types (DDoS, storage failure, resource contention, noisy neighbor)~\cite{zheng2024lemma}. Licensed under CC-BY-NC-4.0 (non-commercial). Its log pipeline aggregates raw log lines into three-dimensional time series using golden-signal keywords that cover only \texttt{error}, \texttt{exception}, and \texttt{critical} --- notably \textbf{not} \texttt{warning} --- so it offers no independent warning tier.

\item[Fault Injection]
The controlled introduction of a failure (e.g., CPU stress, network latency, packet loss) into a microservice system at a known timestamp $T_{\text{inject}}$ to generate labeled failure cases for benchmarking.

\item[Normal Window]
The interval $[T_{\text{start}}, T_{\text{inject}})$ preceding fault injection, during which the system operates normally; alerts occurring here are candidates for \textit{Transient} noise. Note that the dataset fields \texttt{normal\_timesteps} and \texttt{faulty\_timesteps} are integer counts of sample points, not timestamp lists, and therefore define these windows only indirectly.

\item[Fault Window]
The interval $[T_{\text{inject}}, T_{\text{end}}]$ following fault injection, during which the system exhibits degradation; alerts here are labeled as \textit{Signal} (if from the root cause service) or \textit{Cascading} (if from a service reachable from the root cause service in the call graph).

\item[Ground Truth Root Cause Label]
The authoritative annotation of which service is the root cause of a failure case, provided by the RCAEval benchmark. Used to compute RCPR and pairwise metrics.

% =============================================================================
% BASELINE & INFRASTRUCTURE TERMS
% =============================================================================

\item[Prometheus Alertmanager]
The de facto standard alert routing and grouping component in the Cloud Native Computing Foundation (CNCF) ecosystem. Groups alerts by label matchers and time windows (\texttt{group\_wait}, \texttt{group\_interval}, \texttt{group\_by})~\cite{prometheus-alertmanager}.

\item[Group Wait (\texttt{group\_wait})]
The initial delay after the first alert in a group before sending a notification, allowing additional alerts to arrive and be grouped. Default: 30~seconds.

\item[Group Interval (\texttt{group\_interval})]
The minimum interval between subsequent notifications for the same group. Default: 5~minutes.

\item[Drain3]
An online log parsing algorithm that extracts log templates from raw log streams by clustering log messages based on token similarity, producing parameterized templates with variable placeholders~\cite{drain3}.

\item[Sentence-BERT (S-BERT)]
A modification of BERT that uses siamese and triplet network structures to derive semantically meaningful sentence embeddings, optimized for cosine similarity comparison~\cite{sbert}.

\item[HDBSCAN]
Hierarchical Density-Based Spatial Clustering of Applications with Noise; a density-based clustering algorithm that does not require specifying the number of clusters and can identify outliers (noise points)~\cite{hdbscan}.

% =============================================================================
% EXPERIMENTAL DESIGN TERMS
% =============================================================================

\item[Ablation Study]
An experimental procedure where individual components (e.g., \texttt{get\_service\_dependency} tool, fast filter, runbook) are removed from the full system to measure their individual contribution to performance metrics.

\item[Cross-System Generalization]
The ability of an aggregation method trained or tuned on one microservice system (e.g., Online Boutique) to maintain performance on unseen systems (e.g., Train Ticket) without retraining.

\item[Statistical Significance Test]
Paired Wilcoxon signed-rank test (non-parametric) used to compare metric distributions across the four arms on the same failure cases, with significance level $\alpha = 0.05$.

\item[Metric Gaming]
The phenomenon where a method optimizes a single metric (e.g., ARR) at the expense of operational utility (e.g., grouping all alerts into one group achieves ARR $\approx$ 100\% but RCPR = 0\%). Mitigated by the multi-metric framework (ARR, RCPR, F1, Purity).

% =============================================================================
% INDUSTRY STANDARDS & REFERENCES
% =============================================================================

\item[Google SRE Principles]
Foundational Site Reliability Engineering practices documented in \textit{Site Reliability Engineering} (Ch.~6, 10, 11) and \textit{The Site Reliability Workbook} (Ch.~5), emphasizing actionable alerts, maximum 2 incidents per 12-hour shift, and alerting on SLOs rather than raw metrics~\cite{google-sre-book, google-sre-workbook}.

\item[Observability Pillars]
The three canonical telemetry types: \textit{Logs} (immutable event records), \textit{Metrics} (aggregated numerical measurements), and \textit{Traces} (distributed request flow records)~\cite{grafana2025}.

\end{description}

% =============================================================================
% LOCAL BIBLIOGRAPHY MACROS (if not using main document's .bib)
% Include these in your main .bib or define via IEEEabrv.bib/IEEEfull.bib
% =============================================================================
% @string{srebook = "Site Reliability Engineering: How Google Runs Production Systems"}
% @string{sreworkbook = "The Site Reliability Workbook: Practical Ways to Implement SRE"}
% @article{pagerduty2020, ...}
% @article{grafana2025, ...}
% @inproceedings{kuang2024cola, ...}
% @inproceedings{pham2025rcaeval, ...}
% @article{zheng2024lemma, ...}
% @inproceedings{yu2025alertguardian, ...}
% @manual{prometheus-alertmanager, ...}
% @article{sbert, ...}
% @article{drain3, ...}
% @article{hdbscan, ...}
% @article{wei2022cot, ...} in main IEEEtran document
% Compiles with: pdflatex main.tex (where main.tex includes this file)
% =============================================================================

% -----------------------------------------------------------------------------
% USAGE IN MAIN DOCUMENT:
% \section*{Glossary of Terms}
% \addcontentsline{toc}{section}{Glossary of Terms}
% % =============================================================================
% GLOSSARY OF TERMS — Alert Aggregation in Microservices
% Standalone file for \input{glossary} in main IEEEtran document
% Compiles with: pdflatex main.tex (where main.tex includes this file)
% =============================================================================

% -----------------------------------------------------------------------------
% USAGE IN MAIN DOCUMENT:
% \section*{Glossary of Terms}
% \addcontentsline{toc}{section}{Glossary of Terms}
% \input{glossary}
% -----------------------------------------------------------------------------

\begin{description}

% =============================================================================
% CORE PROBLEM TERMS
% =============================================================================

\item[Alert Fatigue]
A cognitive state experienced by on-call engineers when the volume of alerts exceeds human cognitive capacity, leading to desensitization, missed critical alerts, and delayed incident response~\cite{pagerduty2020, grafana2025}.

\item[Alert Storm]
A sudden surge of alerts triggered by a single root cause propagating through dependent microservices, resulting in hundreds to thousands of alerts within minutes~\cite{kuang2024cola}.

\item[Noisy Alert Warning]
A warning-level alert generated by the monitoring system that is technically correct but does not require meaningful remediation action from an on-call engineer, classified into four categories: \textit{Duplicate}, \textit{Transient}, \textit{Unactionable}, and \textit{Cascading} (defined below).

\item[Early Warning Phase]
The pre-failure period during which a microservice exhibits performance degradation (e.g., increased latency, resource saturation) but has not yet violated Service Level Objectives (SLOs) or caused user-facing outages.

\item[Actionable Alert]
An alert that satisfies the Google SRE principle: ``every alert sent to an on-call engineer must require immediate human intervention''~\cite{google-sre-ch10, google-sre-ch11}.

% =============================================================================
% FOUR CATEGORIES OF NOISY ALERT WARNINGS
% =============================================================================

\item[Duplicate Warning Alert]
Multiple alerts originating from the same service with identical or near-identical message templates within a short time window ($\Delta t \le 60$~seconds), representing repeated observations of the same underlying condition rather than distinct incidents.

\item[Transient Warning Alert]
A warning alert that occurs during normal system operation (i.e., before fault injection time $T_{\text{inject}}$ in RCAEval) and resolves spontaneously without human intervention, typically caused by temporary network fluctuations or brief resource spikes.

\item[Unactionable Warning Alert]
A technically accurate warning alert for which no Standard Operating Procedure (SOP) or runbook exists, leaving the on-call engineer without a defined remediation path; the alert consumes attention but yields no operational value.

\item[Cascading Warning Alert]
A warning alert emitted by a downstream service that is a symptom of an upstream root cause; remediating the alerting service does not resolve the underlying issue, which originates in a different service (the root cause service).

% =============================================================================
% METHODOLOGY: FOUR EXPERIMENTAL ARMS
% =============================================================================

\item[Arm 1: Rule-Based Aggregation (Baseline)]
An alert grouping approach that clusters alerts based on static label matching (e.g., service name, alert name, severity) and fixed time windows (\texttt{group\_wait}, \texttt{group\_interval}), as implemented in Prometheus Alertmanager~\cite{prometheus-alertmanager}. Lacks awareness of service dependencies and semantic content.

\item[Arm 2: Semantic Similarity-Based Aggregation]
A method that parses raw logs into templates (via Drain3~\cite{drain3}), encodes templates into dense vector embeddings using Sentence-BERT~\cite{sbert}, and clusters alerts using density-based algorithms (HDBSCAN/DBSCAN) based on cosine similarity. Operates without topology knowledge.

\item[Arm 3: Temporal-Spatial Correlation Aggregation]
A topology-aware method that constructs a static service dependency graph (caller-callee adjacency matrix) from microservice architecture, then correlates alerts occurring within a sliding time window ($\Delta t \in \{30, 60, 120\}$~seconds) if a dependency path of length $k \le 2$ exists between the alerting services.

\item[Arm 4: LLM-Based Agentic Aggregation (Proposed)]
A hybrid two-stage approach: (1) a fast filter removes obvious duplicates and performs time-window pre-grouping; (2) an LLM agent (GPT-4o-mini / DeepSeek-V3) reasons over the remaining clusters using Chain-of-Thought prompting and three tools: \texttt{get\_related\_alerts}, \texttt{get\_service\_dependency}, and optionally \texttt{get\_runbook}. Limited to a maximum of two tool calls per grouping decision.

% =============================================================================
% TECHNICAL CONCEPTS & ARCHITECTURE
% =============================================================================

\item[Service Topology (Service Dependency Graph)]
A directed graph $G = (V, E)$ where vertices $V$ represent microservices and edges $E$ represent synchronous (HTTP/gRPC) or asynchronous (message queue) call relationships. Used to distinguish root cause services from downstream symptom services.

\item[Semantic Similarity]
A measure of meaning-based resemblance between two log messages, computed as the cosine similarity between their dense vector embeddings (typically 384-dimensional from \texttt{sentence-transformers/all-MiniLM-L6-v2}). Captures paraphrases and synonyms that lexical matching misses.

\item[Runbook Reasoning]
The process by which an LLM agent consults a synthetic runbook (a structured troubleshooting guide for a specific fault type) to determine whether an alert is actionable, requires escalation, or can be safely aggregated as noise.

\item[Tool-Calling Agent]
An LLM agent augmented with a predefined set of functions (tools) it can invoke during reasoning: \texttt{get\_related\_alerts(service, time\_window)}, \texttt{get\_service\_dependency(service)}, and \texttt{get\_runbook(fault\_type)}. Tool outputs are fed back into the agent's context for subsequent reasoning steps.

\item[Chain-of-Thought (CoT) Prompting]
A prompting technique that instructs the LLM to generate intermediate reasoning steps before producing a final answer, improving accuracy on multi-step causal inference tasks~\cite{wei2022cot}.

\item[Fast Filter (Deduplication + Time-Window Pre-Grouping)]
A lightweight deterministic preprocessing stage that removes exact duplicate alerts and performs initial grouping by service and time window, reducing the alert volume passed to the LLM agent by an estimated 70--80\%.

\item[Synthetic Runbook]
A researcher-authored troubleshooting document for each fault type in the benchmark dataset (e.g., CPU stress, memory leak, network latency), containing symptoms, diagnosis steps, and remediation actions. Not a production SOP; used solely for experimental evaluation.

% =============================================================================
% EVALUATION METRICS
% =============================================================================

\item[Alert Reduction Ratio (ARR)]
The primary metric quantifying alert volume reduction:
\begin{equation}
\text{ARR} = \left(1 - \frac{N_{\text{groups}}}{N_{\text{raw\_alerts}}}\right) \times 100\%
\end{equation}
where $N_{\text{raw\_alerts}}$ is the number of warning alerts before aggregation and $N_{\text{groups}}$ is the number of aggregated groups presented to the engineer. Target: $\text{ARR} \ge 60\%$.

\item[Root Cause Preservation Rate (RCPR)]
The critical safety metric ensuring the root cause evidence is not lost during aggregation. It is defined at two granularities. At service granularity, which is the primary metric and spans every case that exposes logs ($M = 359$):
\begin{equation}
\text{RCPR}_{\text{svc}} = \frac{1}{M} \sum_{i=1}^{M} \mathbb{I}(s_{\text{root}}^{(i)} \in \text{PrimaryGroup}_i) \times 100\%
\end{equation}
At alert granularity, a secondary metric restricted to the $M'$ cases whose ground-truth \texttt{root\_cause.txt} indicator line is itself a warning-tier log ($M' = 4$):
\begin{equation}
\text{RCPR}_{\text{alert}} = \frac{1}{M'} \sum_{i=1}^{M'} \mathbb{I}(\text{RootCauseAlert}_i \in \text{PrimaryGroup}_i) \times 100\%
\end{equation}
where $\mathbb{I}$ is the indicator function. Target: $\text{RCPR} \ge 95\%$ (mandatory).

\item[Pairwise Grouping Precision ($P_p$)]
The fraction of alert pairs placed in the same group that truly share the same root cause:
\begin{equation}
P_p = \frac{|\{(a_x, a_y) : \text{same group} \land \text{same root cause}\}|}{|\{(a_x, a_y) : \text{same group}\}|}
\end{equation}

\item[Pairwise Grouping Recall ($R_p$)]
The fraction of alert pairs sharing the same root cause that are successfully grouped together:
\begin{equation}
R_p = \frac{|\{(a_x, a_y) : \text{same group} \land \text{same root cause}\}|}{|\{(a_x, a_y) : \text{same root cause}\}|}
\end{equation}

\item[Pairwise Grouping F1-Score ($F1_p$)]
The harmonic mean of pairwise precision and recall:
\begin{equation}
F1_p = 2 \times \frac{P_p \times R_p}{P_p + R_p}
\end{equation}
Target: $F1_p \ge 0.85$.

\item[Group Purity]
The average homogeneity of root causes within each generated group:
\begin{equation}
\text{Purity} = \frac{1}{N_{\text{groups}}} \sum_{g=1}^{N_{\text{groups}}} \frac{\max_{c} |\{a \in g : \text{root\_cause}(a) = c\}|}{|g|}
\end{equation}
Target: $\text{Purity} \ge 0.80$.

\item[Processing Latency]
The wall-clock time required to process 100 warning alerts through the aggregation pipeline, measured in milliseconds. For Arm 4, includes LLM inference time and token overhead.

% =============================================================================
% DATASET & BENCHMARK TERMS
% =============================================================================

\item[RCAEval]
A benchmark for root cause analysis in microservice systems comprising 735 failure cases across three systems (Online Boutique, Sock Shop, Train Ticket) with logs, metrics, traces, and root cause labels~\cite{pham2025rcaeval}. Licensed under MIT. Direct inspection shows that only \textbf{359 of the 735 cases ship raw log telemetry} (suite RE1 is metric-only) and that the log schema contains \texttt{timestamp}, \texttt{container\_name}, and \texttt{message} but \textbf{no severity column}.

\item[Severity Tier Assignment]
The three-tier procedure by which this work assigns the severity tier $\ell_i$, since RCAEval supplies none: \emph{Tier 1 (explicit)} parses a level token from the message text (Sock Shop, Train Ticket); \emph{Tier 2 (frequency-inferred)} flags a template as anomalous when its fault-window frequency exceeds its normal-window frequency by a ratio $R$ (Online Boutique); \emph{Tier 3 (ground-truth marker)} tags the \texttt{root\_cause.txt} indicator line as \texttt{ROOT\_CAUSE} for the eight annotated cases.

\item[LEMMA-RCA]
A large-scale multi-modal dataset for root cause analysis containing hundreds of microservice entities and four fault types (DDoS, storage failure, resource contention, noisy neighbor)~\cite{zheng2024lemma}. Licensed under CC-BY-NC-4.0 (non-commercial). Its log pipeline aggregates raw log lines into three-dimensional time series using golden-signal keywords that cover only \texttt{error}, \texttt{exception}, and \texttt{critical} --- notably \textbf{not} \texttt{warning} --- so it offers no independent warning tier.

\item[Fault Injection]
The controlled introduction of a failure (e.g., CPU stress, network latency, packet loss) into a microservice system at a known timestamp $T_{\text{inject}}$ to generate labeled failure cases for benchmarking.

\item[Normal Window]
The interval $[T_{\text{start}}, T_{\text{inject}})$ preceding fault injection, during which the system operates normally; alerts occurring here are candidates for \textit{Transient} noise. Note that the dataset fields \texttt{normal\_timesteps} and \texttt{faulty\_timesteps} are integer counts of sample points, not timestamp lists, and therefore define these windows only indirectly.

\item[Fault Window]
The interval $[T_{\text{inject}}, T_{\text{end}}]$ following fault injection, during which the system exhibits degradation; alerts here are labeled as \textit{Signal} (if from the root cause service) or \textit{Cascading} (if from a service reachable from the root cause service in the call graph).

\item[Ground Truth Root Cause Label]
The authoritative annotation of which service is the root cause of a failure case, provided by the RCAEval benchmark. Used to compute RCPR and pairwise metrics.

% =============================================================================
% BASELINE & INFRASTRUCTURE TERMS
% =============================================================================

\item[Prometheus Alertmanager]
The de facto standard alert routing and grouping component in the Cloud Native Computing Foundation (CNCF) ecosystem. Groups alerts by label matchers and time windows (\texttt{group\_wait}, \texttt{group\_interval}, \texttt{group\_by})~\cite{prometheus-alertmanager}.

\item[Group Wait (\texttt{group\_wait})]
The initial delay after the first alert in a group before sending a notification, allowing additional alerts to arrive and be grouped. Default: 30~seconds.

\item[Group Interval (\texttt{group\_interval})]
The minimum interval between subsequent notifications for the same group. Default: 5~minutes.

\item[Drain3]
An online log parsing algorithm that extracts log templates from raw log streams by clustering log messages based on token similarity, producing parameterized templates with variable placeholders~\cite{drain3}.

\item[Sentence-BERT (S-BERT)]
A modification of BERT that uses siamese and triplet network structures to derive semantically meaningful sentence embeddings, optimized for cosine similarity comparison~\cite{sbert}.

\item[HDBSCAN]
Hierarchical Density-Based Spatial Clustering of Applications with Noise; a density-based clustering algorithm that does not require specifying the number of clusters and can identify outliers (noise points)~\cite{hdbscan}.

% =============================================================================
% EXPERIMENTAL DESIGN TERMS
% =============================================================================

\item[Ablation Study]
An experimental procedure where individual components (e.g., \texttt{get\_service\_dependency} tool, fast filter, runbook) are removed from the full system to measure their individual contribution to performance metrics.

\item[Cross-System Generalization]
The ability of an aggregation method trained or tuned on one microservice system (e.g., Online Boutique) to maintain performance on unseen systems (e.g., Train Ticket) without retraining.

\item[Statistical Significance Test]
Paired Wilcoxon signed-rank test (non-parametric) used to compare metric distributions across the four arms on the same failure cases, with significance level $\alpha = 0.05$.

\item[Metric Gaming]
The phenomenon where a method optimizes a single metric (e.g., ARR) at the expense of operational utility (e.g., grouping all alerts into one group achieves ARR $\approx$ 100\% but RCPR = 0\%). Mitigated by the multi-metric framework (ARR, RCPR, F1, Purity).

% =============================================================================
% INDUSTRY STANDARDS & REFERENCES
% =============================================================================

\item[Google SRE Principles]
Foundational Site Reliability Engineering practices documented in \textit{Site Reliability Engineering} (Ch.~6, 10, 11) and \textit{The Site Reliability Workbook} (Ch.~5), emphasizing actionable alerts, maximum 2 incidents per 12-hour shift, and alerting on SLOs rather than raw metrics~\cite{google-sre-book, google-sre-workbook}.

\item[Observability Pillars]
The three canonical telemetry types: \textit{Logs} (immutable event records), \textit{Metrics} (aggregated numerical measurements), and \textit{Traces} (distributed request flow records)~\cite{grafana2025}.

\end{description}

% =============================================================================
% LOCAL BIBLIOGRAPHY MACROS (if not using main document's .bib)
% Include these in your main .bib or define via IEEEabrv.bib/IEEEfull.bib
% =============================================================================
% @string{srebook = "Site Reliability Engineering: How Google Runs Production Systems"}
% @string{sreworkbook = "The Site Reliability Workbook: Practical Ways to Implement SRE"}
% @article{pagerduty2020, ...}
% @article{grafana2025, ...}
% @inproceedings{kuang2024cola, ...}
% @inproceedings{pham2025rcaeval, ...}
% @article{zheng2024lemma, ...}
% @inproceedings{yu2025alertguardian, ...}
% @manual{prometheus-alertmanager, ...}
% @article{sbert, ...}
% @article{drain3, ...}
% @article{hdbscan, ...}
% @article{wei2022cot, ...}
% -----------------------------------------------------------------------------

\begin{description}

% =============================================================================
% CORE PROBLEM TERMS
% =============================================================================

\item[Alert Fatigue]
A cognitive state experienced by on-call engineers when the volume of alerts exceeds human cognitive capacity, leading to desensitization, missed critical alerts, and delayed incident response~\cite{pagerduty2020, grafana2025}.

\item[Alert Storm]
A sudden surge of alerts triggered by a single root cause propagating through dependent microservices, resulting in hundreds to thousands of alerts within minutes~\cite{kuang2024cola}.

\item[Noisy Alert Warning]
A warning-level alert generated by the monitoring system that is technically correct but does not require meaningful remediation action from an on-call engineer, classified into four categories: \textit{Duplicate}, \textit{Transient}, \textit{Unactionable}, and \textit{Cascading} (defined below).

\item[Early Warning Phase]
The pre-failure period during which a microservice exhibits performance degradation (e.g., increased latency, resource saturation) but has not yet violated Service Level Objectives (SLOs) or caused user-facing outages.

\item[Actionable Alert]
An alert that satisfies the Google SRE principle: ``every alert sent to an on-call engineer must require immediate human intervention''~\cite{google-sre-ch10, google-sre-ch11}.

% =============================================================================
% FOUR CATEGORIES OF NOISY ALERT WARNINGS
% =============================================================================

\item[Duplicate Warning Alert]
Multiple alerts originating from the same service with identical or near-identical message templates within a short time window ($\Delta t \le 60$~seconds), representing repeated observations of the same underlying condition rather than distinct incidents.

\item[Transient Warning Alert]
A warning alert that occurs during normal system operation (i.e., before fault injection time $T_{\text{inject}}$ in RCAEval) and resolves spontaneously without human intervention, typically caused by temporary network fluctuations or brief resource spikes.

\item[Unactionable Warning Alert]
A technically accurate warning alert for which no Standard Operating Procedure (SOP) or runbook exists, leaving the on-call engineer without a defined remediation path; the alert consumes attention but yields no operational value.

\item[Cascading Warning Alert]
A warning alert emitted by a downstream service that is a symptom of an upstream root cause; remediating the alerting service does not resolve the underlying issue, which originates in a different service (the root cause service).

% =============================================================================
% METHODOLOGY: FOUR EXPERIMENTAL ARMS
% =============================================================================

\item[Arm 1: Rule-Based Aggregation (Baseline)]
An alert grouping approach that clusters alerts based on static label matching (e.g., service name, alert name, severity) and fixed time windows (\texttt{group\_wait}, \texttt{group\_interval}), as implemented in Prometheus Alertmanager~\cite{prometheus-alertmanager}. Lacks awareness of service dependencies and semantic content.

\item[Arm 2: Semantic Similarity-Based Aggregation]
A method that parses raw logs into templates (via Drain3~\cite{drain3}), encodes templates into dense vector embeddings using Sentence-BERT~\cite{sbert}, and clusters alerts using density-based algorithms (HDBSCAN/DBSCAN) based on cosine similarity. Operates without topology knowledge.

\item[Arm 3: Temporal-Spatial Correlation Aggregation]
A topology-aware method that constructs a static service dependency graph (caller-callee adjacency matrix) from microservice architecture, then correlates alerts occurring within a sliding time window ($\Delta t \in \{30, 60, 120\}$~seconds) if a dependency path of length $k \le 2$ exists between the alerting services.

\item[Arm 4: LLM-Based Agentic Aggregation (Proposed)]
A hybrid two-stage approach: (1) a fast filter removes obvious duplicates and performs time-window pre-grouping; (2) an LLM agent (GPT-4o-mini / DeepSeek-V3) reasons over the remaining clusters using Chain-of-Thought prompting and three tools: \texttt{get\_related\_alerts}, \texttt{get\_service\_dependency}, and optionally \texttt{get\_runbook}. Limited to a maximum of two tool calls per grouping decision.

% =============================================================================
% TECHNICAL CONCEPTS & ARCHITECTURE
% =============================================================================

\item[Service Topology (Service Dependency Graph)]
A directed graph $G = (V, E)$ where vertices $V$ represent microservices and edges $E$ represent synchronous (HTTP/gRPC) or asynchronous (message queue) call relationships. Used to distinguish root cause services from downstream symptom services.

\item[Semantic Similarity]
A measure of meaning-based resemblance between two log messages, computed as the cosine similarity between their dense vector embeddings (typically 384-dimensional from \texttt{sentence-transformers/all-MiniLM-L6-v2}). Captures paraphrases and synonyms that lexical matching misses.

\item[Runbook Reasoning]
The process by which an LLM agent consults a synthetic runbook (a structured troubleshooting guide for a specific fault type) to determine whether an alert is actionable, requires escalation, or can be safely aggregated as noise.

\item[Tool-Calling Agent]
An LLM agent augmented with a predefined set of functions (tools) it can invoke during reasoning: \texttt{get\_related\_alerts(service, time\_window)}, \texttt{get\_service\_dependency(service)}, and \texttt{get\_runbook(fault\_type)}. Tool outputs are fed back into the agent's context for subsequent reasoning steps.

\item[Chain-of-Thought (CoT) Prompting]
A prompting technique that instructs the LLM to generate intermediate reasoning steps before producing a final answer, improving accuracy on multi-step causal inference tasks~\cite{wei2022cot}.

\item[Fast Filter (Deduplication + Time-Window Pre-Grouping)]
A lightweight deterministic preprocessing stage that removes exact duplicate alerts and performs initial grouping by service and time window, reducing the alert volume passed to the LLM agent by an estimated 70--80\%.

\item[Synthetic Runbook]
A researcher-authored troubleshooting document for each fault type in the benchmark dataset (e.g., CPU stress, memory leak, network latency), containing symptoms, diagnosis steps, and remediation actions. Not a production SOP; used solely for experimental evaluation.

% =============================================================================
% EVALUATION METRICS
% =============================================================================

\item[Alert Reduction Ratio (ARR)]
The primary metric quantifying alert volume reduction:
\begin{equation}
\text{ARR} = \left(1 - \frac{N_{\text{groups}}}{N_{\text{raw\_alerts}}}\right) \times 100\%
\end{equation}
where $N_{\text{raw\_alerts}}$ is the number of warning alerts before aggregation and $N_{\text{groups}}$ is the number of aggregated groups presented to the engineer. Target: $\text{ARR} \ge 60\%$.

\item[Root Cause Preservation Rate (RCPR)]
The critical safety metric ensuring the root cause evidence is not lost during aggregation. It is defined at two granularities. At service granularity, which is the primary metric and spans every case that exposes logs ($M = 359$):
\begin{equation}
\text{RCPR}_{\text{svc}} = \frac{1}{M} \sum_{i=1}^{M} \mathbb{I}(s_{\text{root}}^{(i)} \in \text{PrimaryGroup}_i) \times 100\%
\end{equation}
At alert granularity, a secondary metric restricted to the $M'$ cases whose ground-truth \texttt{root\_cause.txt} indicator line is itself a warning-tier log ($M' = 4$):
\begin{equation}
\text{RCPR}_{\text{alert}} = \frac{1}{M'} \sum_{i=1}^{M'} \mathbb{I}(\text{RootCauseAlert}_i \in \text{PrimaryGroup}_i) \times 100\%
\end{equation}
where $\mathbb{I}$ is the indicator function. Target: $\text{RCPR} \ge 95\%$ (mandatory).

\item[Pairwise Grouping Precision ($P_p$)]
The fraction of alert pairs placed in the same group that truly share the same root cause:
\begin{equation}
P_p = \frac{|\{(a_x, a_y) : \text{same group} \land \text{same root cause}\}|}{|\{(a_x, a_y) : \text{same group}\}|}
\end{equation}

\item[Pairwise Grouping Recall ($R_p$)]
The fraction of alert pairs sharing the same root cause that are successfully grouped together:
\begin{equation}
R_p = \frac{|\{(a_x, a_y) : \text{same group} \land \text{same root cause}\}|}{|\{(a_x, a_y) : \text{same root cause}\}|}
\end{equation}

\item[Pairwise Grouping F1-Score ($F1_p$)]
The harmonic mean of pairwise precision and recall:
\begin{equation}
F1_p = 2 \times \frac{P_p \times R_p}{P_p + R_p}
\end{equation}
Target: $F1_p \ge 0.85$.

\item[Group Purity]
The average homogeneity of root causes within each generated group:
\begin{equation}
\text{Purity} = \frac{1}{N_{\text{groups}}} \sum_{g=1}^{N_{\text{groups}}} \frac{\max_{c} |\{a \in g : \text{root\_cause}(a) = c\}|}{|g|}
\end{equation}
Target: $\text{Purity} \ge 0.80$.

\item[Processing Latency]
The wall-clock time required to process 100 warning alerts through the aggregation pipeline, measured in milliseconds. For Arm 4, includes LLM inference time and token overhead.

% =============================================================================
% DATASET & BENCHMARK TERMS
% =============================================================================

\item[RCAEval]
A benchmark for root cause analysis in microservice systems comprising 735 failure cases across three systems (Online Boutique, Sock Shop, Train Ticket) with logs, metrics, traces, and root cause labels~\cite{pham2025rcaeval}. Licensed under MIT. Direct inspection shows that only \textbf{359 of the 735 cases ship raw log telemetry} (suite RE1 is metric-only) and that the log schema contains \texttt{timestamp}, \texttt{container\_name}, and \texttt{message} but \textbf{no severity column}.

\item[Severity Tier Assignment]
The three-tier procedure by which this work assigns the severity tier $\ell_i$, since RCAEval supplies none: \emph{Tier 1 (explicit)} parses a level token from the message text (Sock Shop, Train Ticket); \emph{Tier 2 (frequency-inferred)} flags a template as anomalous when its fault-window frequency exceeds its normal-window frequency by a ratio $R$ (Online Boutique); \emph{Tier 3 (ground-truth marker)} tags the \texttt{root\_cause.txt} indicator line as \texttt{ROOT\_CAUSE} for the eight annotated cases.

\item[LEMMA-RCA]
A large-scale multi-modal dataset for root cause analysis containing hundreds of microservice entities and four fault types (DDoS, storage failure, resource contention, noisy neighbor)~\cite{zheng2024lemma}. Licensed under CC-BY-NC-4.0 (non-commercial). Its log pipeline aggregates raw log lines into three-dimensional time series using golden-signal keywords that cover only \texttt{error}, \texttt{exception}, and \texttt{critical} --- notably \textbf{not} \texttt{warning} --- so it offers no independent warning tier.

\item[Fault Injection]
The controlled introduction of a failure (e.g., CPU stress, network latency, packet loss) into a microservice system at a known timestamp $T_{\text{inject}}$ to generate labeled failure cases for benchmarking.

\item[Normal Window]
The interval $[T_{\text{start}}, T_{\text{inject}})$ preceding fault injection, during which the system operates normally; alerts occurring here are candidates for \textit{Transient} noise. Note that the dataset fields \texttt{normal\_timesteps} and \texttt{faulty\_timesteps} are integer counts of sample points, not timestamp lists, and therefore define these windows only indirectly.

\item[Fault Window]
The interval $[T_{\text{inject}}, T_{\text{end}}]$ following fault injection, during which the system exhibits degradation; alerts here are labeled as \textit{Signal} (if from the root cause service) or \textit{Cascading} (if from a service reachable from the root cause service in the call graph).

\item[Ground Truth Root Cause Label]
The authoritative annotation of which service is the root cause of a failure case, provided by the RCAEval benchmark. Used to compute RCPR and pairwise metrics.

% =============================================================================
% BASELINE & INFRASTRUCTURE TERMS
% =============================================================================

\item[Prometheus Alertmanager]
The de facto standard alert routing and grouping component in the Cloud Native Computing Foundation (CNCF) ecosystem. Groups alerts by label matchers and time windows (\texttt{group\_wait}, \texttt{group\_interval}, \texttt{group\_by})~\cite{prometheus-alertmanager}.

\item[Group Wait (\texttt{group\_wait})]
The initial delay after the first alert in a group before sending a notification, allowing additional alerts to arrive and be grouped. Default: 30~seconds.

\item[Group Interval (\texttt{group\_interval})]
The minimum interval between subsequent notifications for the same group. Default: 5~minutes.

\item[Drain3]
An online log parsing algorithm that extracts log templates from raw log streams by clustering log messages based on token similarity, producing parameterized templates with variable placeholders~\cite{drain3}.

\item[Sentence-BERT (S-BERT)]
A modification of BERT that uses siamese and triplet network structures to derive semantically meaningful sentence embeddings, optimized for cosine similarity comparison~\cite{sbert}.

\item[HDBSCAN]
Hierarchical Density-Based Spatial Clustering of Applications with Noise; a density-based clustering algorithm that does not require specifying the number of clusters and can identify outliers (noise points)~\cite{hdbscan}.

% =============================================================================
% EXPERIMENTAL DESIGN TERMS
% =============================================================================

\item[Ablation Study]
An experimental procedure where individual components (e.g., \texttt{get\_service\_dependency} tool, fast filter, runbook) are removed from the full system to measure their individual contribution to performance metrics.

\item[Cross-System Generalization]
The ability of an aggregation method trained or tuned on one microservice system (e.g., Online Boutique) to maintain performance on unseen systems (e.g., Train Ticket) without retraining.

\item[Statistical Significance Test]
Paired Wilcoxon signed-rank test (non-parametric) used to compare metric distributions across the four arms on the same failure cases, with significance level $\alpha = 0.05$.

\item[Metric Gaming]
The phenomenon where a method optimizes a single metric (e.g., ARR) at the expense of operational utility (e.g., grouping all alerts into one group achieves ARR $\approx$ 100\% but RCPR = 0\%). Mitigated by the multi-metric framework (ARR, RCPR, F1, Purity).

% =============================================================================
% INDUSTRY STANDARDS & REFERENCES
% =============================================================================

\item[Google SRE Principles]
Foundational Site Reliability Engineering practices documented in \textit{Site Reliability Engineering} (Ch.~6, 10, 11) and \textit{The Site Reliability Workbook} (Ch.~5), emphasizing actionable alerts, maximum 2 incidents per 12-hour shift, and alerting on SLOs rather than raw metrics~\cite{google-sre-book, google-sre-workbook}.

\item[Observability Pillars]
The three canonical telemetry types: \textit{Logs} (immutable event records), \textit{Metrics} (aggregated numerical measurements), and \textit{Traces} (distributed request flow records)~\cite{grafana2025}.

\end{description}

% =============================================================================
% LOCAL BIBLIOGRAPHY MACROS (if not using main document's .bib)
% Include these in your main .bib or define via IEEEabrv.bib/IEEEfull.bib
% =============================================================================
% @string{srebook = "Site Reliability Engineering: How Google Runs Production Systems"}
% @string{sreworkbook = "The Site Reliability Workbook: Practical Ways to Implement SRE"}
% @article{pagerduty2020, ...}
% @article{grafana2025, ...}
% @inproceedings{kuang2024cola, ...}
% @inproceedings{pham2025rcaeval, ...}
% @article{zheng2024lemma, ...}
% @inproceedings{yu2025alertguardian, ...}
% @manual{prometheus-alertmanager, ...}
% @article{sbert, ...}
% @article{drain3, ...}
% @article{hdbscan, ...}
% @article{wei2022cot, ...}
% -----------------------------------------------------------------------------

\begin{description}

% =============================================================================
% CORE PROBLEM TERMS
% =============================================================================

\item[Alert Fatigue]
A cognitive state experienced by on-call engineers when the volume of alerts exceeds human cognitive capacity, leading to desensitization, missed critical alerts, and delayed incident response~\cite{pagerduty2020, grafana2025}.

\item[Alert Storm]
A sudden surge of alerts triggered by a single root cause propagating through dependent microservices, resulting in hundreds to thousands of alerts within minutes~\cite{kuang2024cola}.

\item[Noisy Alert Warning]
A warning-level alert generated by the monitoring system that is technically correct but does not require meaningful remediation action from an on-call engineer, classified into four categories: \textit{Duplicate}, \textit{Transient}, \textit{Unactionable}, and \textit{Cascading} (defined below).

\item[Early Warning Phase]
The pre-failure period during which a microservice exhibits performance degradation (e.g., increased latency, resource saturation) but has not yet violated Service Level Objectives (SLOs) or caused user-facing outages.

\item[Actionable Alert]
An alert that satisfies the Google SRE principle: ``every alert sent to an on-call engineer must require immediate human intervention''~\cite{google-sre-ch10, google-sre-ch11}.

% =============================================================================
% FOUR CATEGORIES OF NOISY ALERT WARNINGS
% =============================================================================

\item[Duplicate Warning Alert]
Multiple alerts originating from the same service with identical or near-identical message templates within a short time window ($\Delta t \le 60$~seconds), representing repeated observations of the same underlying condition rather than distinct incidents.

\item[Transient Warning Alert]
A warning alert that occurs during normal system operation (i.e., before fault injection time $T_{\text{inject}}$ in RCAEval) and resolves spontaneously without human intervention, typically caused by temporary network fluctuations or brief resource spikes.

\item[Unactionable Warning Alert]
A technically accurate warning alert for which no Standard Operating Procedure (SOP) or runbook exists, leaving the on-call engineer without a defined remediation path; the alert consumes attention but yields no operational value.

\item[Cascading Warning Alert]
A warning alert emitted by a downstream service that is a symptom of an upstream root cause; remediating the alerting service does not resolve the underlying issue, which originates in a different service (the root cause service).

% =============================================================================
% METHODOLOGY: FOUR EXPERIMENTAL ARMS
% =============================================================================

\item[Arm 1: Rule-Based Aggregation (Baseline)]
An alert grouping approach that clusters alerts based on static label matching (e.g., service name, alert name, severity) and fixed time windows (\texttt{group\_wait}, \texttt{group\_interval}), as implemented in Prometheus Alertmanager~\cite{prometheus-alertmanager}. Lacks awareness of service dependencies and semantic content.

\item[Arm 2: Semantic Similarity-Based Aggregation]
A method that parses raw logs into templates (via Drain3~\cite{drain3}), encodes templates into dense vector embeddings using Sentence-BERT~\cite{sbert}, and clusters alerts using density-based algorithms (HDBSCAN/DBSCAN) based on cosine similarity. Operates without topology knowledge.

\item[Arm 3: Temporal-Spatial Correlation Aggregation]
A topology-aware method that constructs a static service dependency graph (caller-callee adjacency matrix) from microservice architecture, then correlates alerts occurring within a sliding time window ($\Delta t \in \{30, 60, 120\}$~seconds) if a dependency path of length $k \le 2$ exists between the alerting services.

\item[Arm 4: LLM-Based Agentic Aggregation (Proposed)]
A hybrid two-stage approach: (1) a fast filter removes obvious duplicates and performs time-window pre-grouping; (2) an LLM agent (GPT-4o-mini / DeepSeek-V3) reasons over the remaining clusters using Chain-of-Thought prompting and three tools: \texttt{get\_related\_alerts}, \texttt{get\_service\_dependency}, and optionally \texttt{get\_runbook}. Limited to a maximum of two tool calls per grouping decision.

% =============================================================================
% TECHNICAL CONCEPTS & ARCHITECTURE
% =============================================================================

\item[Service Topology (Service Dependency Graph)]
A directed graph $G = (V, E)$ where vertices $V$ represent microservices and edges $E$ represent synchronous (HTTP/gRPC) or asynchronous (message queue) call relationships. Used to distinguish root cause services from downstream symptom services.

\item[Semantic Similarity]
A measure of meaning-based resemblance between two log messages, computed as the cosine similarity between their dense vector embeddings (typically 384-dimensional from \texttt{sentence-transformers/all-MiniLM-L6-v2}). Captures paraphrases and synonyms that lexical matching misses.

\item[Runbook Reasoning]
The process by which an LLM agent consults a synthetic runbook (a structured troubleshooting guide for a specific fault type) to determine whether an alert is actionable, requires escalation, or can be safely aggregated as noise.

\item[Tool-Calling Agent]
An LLM agent augmented with a predefined set of functions (tools) it can invoke during reasoning: \texttt{get\_related\_alerts(service, time\_window)}, \texttt{get\_service\_dependency(service)}, and \texttt{get\_runbook(fault\_type)}. Tool outputs are fed back into the agent's context for subsequent reasoning steps.

\item[Chain-of-Thought (CoT) Prompting]
A prompting technique that instructs the LLM to generate intermediate reasoning steps before producing a final answer, improving accuracy on multi-step causal inference tasks~\cite{wei2022cot}.

\item[Fast Filter (Deduplication + Time-Window Pre-Grouping)]
A lightweight deterministic preprocessing stage that removes exact duplicate alerts and performs initial grouping by service and time window, reducing the alert volume passed to the LLM agent by an estimated 70--80\%.

\item[Synthetic Runbook]
A researcher-authored troubleshooting document for each fault type in the benchmark dataset (e.g., CPU stress, memory leak, network latency), containing symptoms, diagnosis steps, and remediation actions. Not a production SOP; used solely for experimental evaluation.

% =============================================================================
% EVALUATION METRICS
% =============================================================================

\item[Alert Reduction Ratio (ARR)]
The primary metric quantifying alert volume reduction:
\begin{equation}
\text{ARR} = \left(1 - \frac{N_{\text{groups}}}{N_{\text{raw\_alerts}}}\right) \times 100\%
\end{equation}
where $N_{\text{raw\_alerts}}$ is the number of warning alerts before aggregation and $N_{\text{groups}}$ is the number of aggregated groups presented to the engineer. Target: $\text{ARR} \ge 60\%$.

\item[Root Cause Preservation Rate (RCPR)]
The critical safety metric ensuring the root cause evidence is not lost during aggregation. It is defined at two granularities. At service granularity, which is the primary metric and spans every case that exposes logs ($M = 359$):
\begin{equation}
\text{RCPR}_{\text{svc}} = \frac{1}{M} \sum_{i=1}^{M} \mathbb{I}(s_{\text{root}}^{(i)} \in \text{PrimaryGroup}_i) \times 100\%
\end{equation}
At alert granularity, a secondary metric restricted to the $M'$ cases whose ground-truth \texttt{root\_cause.txt} indicator line is itself a warning-tier log ($M' = 4$):
\begin{equation}
\text{RCPR}_{\text{alert}} = \frac{1}{M'} \sum_{i=1}^{M'} \mathbb{I}(\text{RootCauseAlert}_i \in \text{PrimaryGroup}_i) \times 100\%
\end{equation}
where $\mathbb{I}$ is the indicator function. Target: $\text{RCPR} \ge 95\%$ (mandatory).

\item[Pairwise Grouping Precision ($P_p$)]
The fraction of alert pairs placed in the same group that truly share the same root cause:
\begin{equation}
P_p = \frac{|\{(a_x, a_y) : \text{same group} \land \text{same root cause}\}|}{|\{(a_x, a_y) : \text{same group}\}|}
\end{equation}

\item[Pairwise Grouping Recall ($R_p$)]
The fraction of alert pairs sharing the same root cause that are successfully grouped together:
\begin{equation}
R_p = \frac{|\{(a_x, a_y) : \text{same group} \land \text{same root cause}\}|}{|\{(a_x, a_y) : \text{same root cause}\}|}
\end{equation}

\item[Pairwise Grouping F1-Score ($F1_p$)]
The harmonic mean of pairwise precision and recall:
\begin{equation}
F1_p = 2 \times \frac{P_p \times R_p}{P_p + R_p}
\end{equation}
Target: $F1_p \ge 0.85$.

\item[Group Purity]
The average homogeneity of root causes within each generated group:
\begin{equation}
\text{Purity} = \frac{1}{N_{\text{groups}}} \sum_{g=1}^{N_{\text{groups}}} \frac{\max_{c} |\{a \in g : \text{root\_cause}(a) = c\}|}{|g|}
\end{equation}
Target: $\text{Purity} \ge 0.80$.

\item[Processing Latency]
The wall-clock time required to process 100 warning alerts through the aggregation pipeline, measured in milliseconds. For Arm 4, includes LLM inference time and token overhead.

% =============================================================================
% DATASET & BENCHMARK TERMS
% =============================================================================

\item[RCAEval]
A benchmark for root cause analysis in microservice systems comprising 735 failure cases across three systems (Online Boutique, Sock Shop, Train Ticket) with logs, metrics, traces, and root cause labels~\cite{pham2025rcaeval}. Licensed under MIT. Direct inspection shows that only \textbf{359 of the 735 cases ship raw log telemetry} (suite RE1 is metric-only) and that the log schema contains \texttt{timestamp}, \texttt{container\_name}, and \texttt{message} but \textbf{no severity column}.

\item[Severity Tier Assignment]
The three-tier procedure by which this work assigns the severity tier $\ell_i$, since RCAEval supplies none: \emph{Tier 1 (explicit)} parses a level token from the message text (Sock Shop, Train Ticket); \emph{Tier 2 (frequency-inferred)} flags a template as anomalous when its fault-window frequency exceeds its normal-window frequency by a ratio $R$ (Online Boutique); \emph{Tier 3 (ground-truth marker)} tags the \texttt{root\_cause.txt} indicator line as \texttt{ROOT\_CAUSE} for the eight annotated cases.

\item[LEMMA-RCA]
A large-scale multi-modal dataset for root cause analysis containing hundreds of microservice entities and four fault types (DDoS, storage failure, resource contention, noisy neighbor)~\cite{zheng2024lemma}. Licensed under CC-BY-NC-4.0 (non-commercial). Its log pipeline aggregates raw log lines into three-dimensional time series using golden-signal keywords that cover only \texttt{error}, \texttt{exception}, and \texttt{critical} --- notably \textbf{not} \texttt{warning} --- so it offers no independent warning tier.

\item[Fault Injection]
The controlled introduction of a failure (e.g., CPU stress, network latency, packet loss) into a microservice system at a known timestamp $T_{\text{inject}}$ to generate labeled failure cases for benchmarking.

\item[Normal Window]
The interval $[T_{\text{start}}, T_{\text{inject}})$ preceding fault injection, during which the system operates normally; alerts occurring here are candidates for \textit{Transient} noise. Note that the dataset fields \texttt{normal\_timesteps} and \texttt{faulty\_timesteps} are integer counts of sample points, not timestamp lists, and therefore define these windows only indirectly.

\item[Fault Window]
The interval $[T_{\text{inject}}, T_{\text{end}}]$ following fault injection, during which the system exhibits degradation; alerts here are labeled as \textit{Signal} (if from the root cause service) or \textit{Cascading} (if from a service reachable from the root cause service in the call graph).

\item[Ground Truth Root Cause Label]
The authoritative annotation of which service is the root cause of a failure case, provided by the RCAEval benchmark. Used to compute RCPR and pairwise metrics.

% =============================================================================
% BASELINE & INFRASTRUCTURE TERMS
% =============================================================================

\item[Prometheus Alertmanager]
The de facto standard alert routing and grouping component in the Cloud Native Computing Foundation (CNCF) ecosystem. Groups alerts by label matchers and time windows (\texttt{group\_wait}, \texttt{group\_interval}, \texttt{group\_by})~\cite{prometheus-alertmanager}.

\item[Group Wait (\texttt{group\_wait})]
The initial delay after the first alert in a group before sending a notification, allowing additional alerts to arrive and be grouped. Default: 30~seconds.

\item[Group Interval (\texttt{group\_interval})]
The minimum interval between subsequent notifications for the same group. Default: 5~minutes.

\item[Drain3]
An online log parsing algorithm that extracts log templates from raw log streams by clustering log messages based on token similarity, producing parameterized templates with variable placeholders~\cite{drain3}.

\item[Sentence-BERT (S-BERT)]
A modification of BERT that uses siamese and triplet network structures to derive semantically meaningful sentence embeddings, optimized for cosine similarity comparison~\cite{sbert}.

\item[HDBSCAN]
Hierarchical Density-Based Spatial Clustering of Applications with Noise; a density-based clustering algorithm that does not require specifying the number of clusters and can identify outliers (noise points)~\cite{hdbscan}.

% =============================================================================
% EXPERIMENTAL DESIGN TERMS
% =============================================================================

\item[Ablation Study]
An experimental procedure where individual components (e.g., \texttt{get\_service\_dependency} tool, fast filter, runbook) are removed from the full system to measure their individual contribution to performance metrics.

\item[Cross-System Generalization]
The ability of an aggregation method trained or tuned on one microservice system (e.g., Online Boutique) to maintain performance on unseen systems (e.g., Train Ticket) without retraining.

\item[Statistical Significance Test]
Paired Wilcoxon signed-rank test (non-parametric) used to compare metric distributions across the four arms on the same failure cases, with significance level $\alpha = 0.05$.

\item[Metric Gaming]
The phenomenon where a method optimizes a single metric (e.g., ARR) at the expense of operational utility (e.g., grouping all alerts into one group achieves ARR $\approx$ 100\% but RCPR = 0\%). Mitigated by the multi-metric framework (ARR, RCPR, F1, Purity).

% =============================================================================
% INDUSTRY STANDARDS & REFERENCES
% =============================================================================

\item[Google SRE Principles]
Foundational Site Reliability Engineering practices documented in \textit{Site Reliability Engineering} (Ch.~6, 10, 11) and \textit{The Site Reliability Workbook} (Ch.~5), emphasizing actionable alerts, maximum 2 incidents per 12-hour shift, and alerting on SLOs rather than raw metrics~\cite{google-sre-book, google-sre-workbook}.

\item[Observability Pillars]
The three canonical telemetry types: \textit{Logs} (immutable event records), \textit{Metrics} (aggregated numerical measurements), and \textit{Traces} (distributed request flow records)~\cite{grafana2025}.

\end{description}

% =============================================================================
% LOCAL BIBLIOGRAPHY MACROS (if not using main document's .bib)
% Include these in your main .bib or define via IEEEabrv.bib/IEEEfull.bib
% =============================================================================
% @string{srebook = "Site Reliability Engineering: How Google Runs Production Systems"}
% @string{sreworkbook = "The Site Reliability Workbook: Practical Ways to Implement SRE"}
% @article{pagerduty2020, ...}
% @article{grafana2025, ...}
% @inproceedings{kuang2024cola, ...}
% @inproceedings{pham2025rcaeval, ...}
% @article{zheng2024lemma, ...}
% @inproceedings{yu2025alertguardian, ...}
% @manual{prometheus-alertmanager, ...}
% @article{sbert, ...}
% @article{drain3, ...}
% @article{hdbscan, ...}
% @article{wei2022cot, ...}
% -----------------------------------------------------------------------------

\begin{description}

% =============================================================================
% CORE PROBLEM TERMS
% =============================================================================

\item[Alert Fatigue]
A cognitive state experienced by on-call engineers when the volume of alerts exceeds human cognitive capacity, leading to desensitization, missed critical alerts, and delayed incident response~\cite{pagerduty2020, grafana2025}.

\item[Alert Storm]
A sudden surge of alerts triggered by a single root cause propagating through dependent microservices, resulting in hundreds to thousands of alerts within minutes~\cite{kuang2024cola}.

\item[Noisy Alert Warning]
A warning-level alert generated by the monitoring system that is technically correct but does not require meaningful remediation action from an on-call engineer, classified into four categories: \textit{Duplicate}, \textit{Transient}, \textit{Unactionable}, and \textit{Cascading} (defined below).

\item[Early Warning Phase]
The pre-failure period during which a microservice exhibits performance degradation (e.g., increased latency, resource saturation) but has not yet violated Service Level Objectives (SLOs) or caused user-facing outages.

\item[Actionable Alert]
An alert that satisfies the Google SRE principle: ``every alert sent to an on-call engineer must require immediate human intervention''~\cite{google-sre-ch10, google-sre-ch11}.

% =============================================================================
% FOUR CATEGORIES OF NOISY ALERT WARNINGS
% =============================================================================

\item[Duplicate Warning Alert]
Multiple alerts originating from the same service with identical or near-identical message templates within a short time window ($\Delta t \le 60$~seconds), representing repeated observations of the same underlying condition rather than distinct incidents.

\item[Transient Warning Alert]
A warning alert that occurs during normal system operation (i.e., before fault injection time $T_{\text{inject}}$ in RCAEval) and resolves spontaneously without human intervention, typically caused by temporary network fluctuations or brief resource spikes.

\item[Unactionable Warning Alert]
A technically accurate warning alert for which no Standard Operating Procedure (SOP) or runbook exists, leaving the on-call engineer without a defined remediation path; the alert consumes attention but yields no operational value.

\item[Cascading Warning Alert]
A warning alert emitted by a downstream service that is a symptom of an upstream root cause; remediating the alerting service does not resolve the underlying issue, which originates in a different service (the root cause service).

% =============================================================================
% METHODOLOGY: FOUR EXPERIMENTAL ARMS
% =============================================================================

\item[Arm 1: Rule-Based Aggregation (Baseline)]
An alert grouping approach that clusters alerts based on static label matching (e.g., service name, alert name, severity) and fixed time windows (\texttt{group\_wait}, \texttt{group\_interval}), as implemented in Prometheus Alertmanager~\cite{prometheus-alertmanager}. Lacks awareness of service dependencies and semantic content.

\item[Arm 2: Semantic Similarity-Based Aggregation]
A method that parses raw logs into templates (via Drain3~\cite{drain3}), encodes templates into dense vector embeddings using Sentence-BERT~\cite{sbert}, and clusters alerts using density-based algorithms (HDBSCAN/DBSCAN) based on cosine similarity. Operates without topology knowledge.

\item[Arm 3: Temporal-Spatial Correlation Aggregation]
A topology-aware method that constructs a static service dependency graph (caller-callee adjacency matrix) from microservice architecture, then correlates alerts occurring within a sliding time window ($\Delta t \in \{30, 60, 120\}$~seconds) if a dependency path of length $k \le 2$ exists between the alerting services.

\item[Arm 4: LLM-Based Agentic Aggregation (Proposed)]
A hybrid two-stage approach: (1) a fast filter removes obvious duplicates and performs time-window pre-grouping; (2) an LLM agent (GPT-4o-mini / DeepSeek-V3) reasons over the remaining clusters using Chain-of-Thought prompting and three tools: \texttt{get\_related\_alerts}, \texttt{get\_service\_dependency}, and optionally \texttt{get\_runbook}. Limited to a maximum of two tool calls per grouping decision.

% =============================================================================
% TECHNICAL CONCEPTS & ARCHITECTURE
% =============================================================================

\item[Service Topology (Service Dependency Graph)]
A directed graph $G = (V, E)$ where vertices $V$ represent microservices and edges $E$ represent synchronous (HTTP/gRPC) or asynchronous (message queue) call relationships. Used to distinguish root cause services from downstream symptom services.

\item[Semantic Similarity]
A measure of meaning-based resemblance between two log messages, computed as the cosine similarity between their dense vector embeddings (typically 384-dimensional from \texttt{sentence-transformers/all-MiniLM-L6-v2}). Captures paraphrases and synonyms that lexical matching misses.

\item[Runbook Reasoning]
The process by which an LLM agent consults a synthetic runbook (a structured troubleshooting guide for a specific fault type) to determine whether an alert is actionable, requires escalation, or can be safely aggregated as noise.

\item[Tool-Calling Agent]
An LLM agent augmented with a predefined set of functions (tools) it can invoke during reasoning: \texttt{get\_related\_alerts(service, time\_window)}, \texttt{get\_service\_dependency(service)}, and \texttt{get\_runbook(fault\_type)}. Tool outputs are fed back into the agent's context for subsequent reasoning steps.

\item[Chain-of-Thought (CoT) Prompting]
A prompting technique that instructs the LLM to generate intermediate reasoning steps before producing a final answer, improving accuracy on multi-step causal inference tasks~\cite{wei2022cot}.

\item[Fast Filter (Deduplication + Time-Window Pre-Grouping)]
A lightweight deterministic preprocessing stage that removes exact duplicate alerts and performs initial grouping by service and time window, reducing the alert volume passed to the LLM agent by an estimated 70--80\%.

\item[Synthetic Runbook]
A researcher-authored troubleshooting document for each fault type in the benchmark dataset (e.g., CPU stress, memory leak, network latency), containing symptoms, diagnosis steps, and remediation actions. Not a production SOP; used solely for experimental evaluation.

% =============================================================================
% EVALUATION METRICS
% =============================================================================

\item[Alert Reduction Ratio (ARR)]
The primary metric quantifying alert volume reduction:
\begin{equation}
\text{ARR} = \left(1 - \frac{N_{\text{groups}}}{N_{\text{raw\_alerts}}}\right) \times 100\%
\end{equation}
where $N_{\text{raw\_alerts}}$ is the number of warning alerts before aggregation and $N_{\text{groups}}$ is the number of aggregated groups presented to the engineer. Target: $\text{ARR} \ge 60\%$.

\item[Root Cause Preservation Rate (RCPR)]
The critical safety metric ensuring the root cause evidence is not lost during aggregation. It is defined at two granularities. At service granularity, which is the primary metric and spans every case that exposes logs ($M = 359$):
\begin{equation}
\text{RCPR}_{\text{svc}} = \frac{1}{M} \sum_{i=1}^{M} \mathbb{I}(s_{\text{root}}^{(i)} \in \text{PrimaryGroup}_i) \times 100\%
\end{equation}
At alert granularity, a secondary metric restricted to the $M'$ cases whose ground-truth \texttt{root\_cause.txt} indicator line is itself a warning-tier log ($M' = 4$):
\begin{equation}
\text{RCPR}_{\text{alert}} = \frac{1}{M'} \sum_{i=1}^{M'} \mathbb{I}(\text{RootCauseAlert}_i \in \text{PrimaryGroup}_i) \times 100\%
\end{equation}
where $\mathbb{I}$ is the indicator function. Target: $\text{RCPR} \ge 95\%$ (mandatory).

\item[Pairwise Grouping Precision ($P_p$)]
The fraction of alert pairs placed in the same group that truly share the same root cause:
\begin{equation}
P_p = \frac{|\{(a_x, a_y) : \text{same group} \land \text{same root cause}\}|}{|\{(a_x, a_y) : \text{same group}\}|}
\end{equation}

\item[Pairwise Grouping Recall ($R_p$)]
The fraction of alert pairs sharing the same root cause that are successfully grouped together:
\begin{equation}
R_p = \frac{|\{(a_x, a_y) : \text{same group} \land \text{same root cause}\}|}{|\{(a_x, a_y) : \text{same root cause}\}|}
\end{equation}

\item[Pairwise Grouping F1-Score ($F1_p$)]
The harmonic mean of pairwise precision and recall:
\begin{equation}
F1_p = 2 \times \frac{P_p \times R_p}{P_p + R_p}
\end{equation}
Target: $F1_p \ge 0.85$.

\item[Group Purity]
The average homogeneity of root causes within each generated group:
\begin{equation}
\text{Purity} = \frac{1}{N_{\text{groups}}} \sum_{g=1}^{N_{\text{groups}}} \frac{\max_{c} |\{a \in g : \text{root\_cause}(a) = c\}|}{|g|}
\end{equation}
Target: $\text{Purity} \ge 0.80$.

\item[Processing Latency]
The wall-clock time required to process 100 warning alerts through the aggregation pipeline, measured in milliseconds. For Arm 4, includes LLM inference time and token overhead.

% =============================================================================
% DATASET & BENCHMARK TERMS
% =============================================================================

\item[RCAEval]
A benchmark for root cause analysis in microservice systems comprising 735 failure cases across three systems (Online Boutique, Sock Shop, Train Ticket) with logs, metrics, traces, and root cause labels~\cite{pham2025rcaeval}. Licensed under MIT. Direct inspection shows that only \textbf{359 of the 735 cases ship raw log telemetry} (suite RE1 is metric-only) and that the log schema contains \texttt{timestamp}, \texttt{container\_name}, and \texttt{message} but \textbf{no severity column}.

\item[Severity Tier Assignment]
The three-tier procedure by which this work assigns the severity tier $\ell_i$, since RCAEval supplies none: \emph{Tier 1 (explicit)} parses a level token from the message text (Sock Shop, Train Ticket); \emph{Tier 2 (frequency-inferred)} flags a template as anomalous when its fault-window frequency exceeds its normal-window frequency by a ratio $R$ (Online Boutique); \emph{Tier 3 (ground-truth marker)} tags the \texttt{root\_cause.txt} indicator line as \texttt{ROOT\_CAUSE} for the eight annotated cases.

\item[LEMMA-RCA]
A large-scale multi-modal dataset for root cause analysis containing hundreds of microservice entities and four fault types (DDoS, storage failure, resource contention, noisy neighbor)~\cite{zheng2024lemma}. Licensed under CC-BY-NC-4.0 (non-commercial). Its log pipeline aggregates raw log lines into three-dimensional time series using golden-signal keywords that cover only \texttt{error}, \texttt{exception}, and \texttt{critical} --- notably \textbf{not} \texttt{warning} --- so it offers no independent warning tier.

\item[Fault Injection]
The controlled introduction of a failure (e.g., CPU stress, network latency, packet loss) into a microservice system at a known timestamp $T_{\text{inject}}$ to generate labeled failure cases for benchmarking.

\item[Normal Window]
The interval $[T_{\text{start}}, T_{\text{inject}})$ preceding fault injection, during which the system operates normally; alerts occurring here are candidates for \textit{Transient} noise. Note that the dataset fields \texttt{normal\_timesteps} and \texttt{faulty\_timesteps} are integer counts of sample points, not timestamp lists, and therefore define these windows only indirectly.

\item[Fault Window]
The interval $[T_{\text{inject}}, T_{\text{end}}]$ following fault injection, during which the system exhibits degradation; alerts here are labeled as \textit{Signal} (if from the root cause service) or \textit{Cascading} (if from a service reachable from the root cause service in the call graph).

\item[Ground Truth Root Cause Label]
The authoritative annotation of which service is the root cause of a failure case, provided by the RCAEval benchmark. Used to compute RCPR and pairwise metrics.

% =============================================================================
% BASELINE & INFRASTRUCTURE TERMS
% =============================================================================

\item[Prometheus Alertmanager]
The de facto standard alert routing and grouping component in the Cloud Native Computing Foundation (CNCF) ecosystem. Groups alerts by label matchers and time windows (\texttt{group\_wait}, \texttt{group\_interval}, \texttt{group\_by})~\cite{prometheus-alertmanager}.

\item[Group Wait (\texttt{group\_wait})]
The initial delay after the first alert in a group before sending a notification, allowing additional alerts to arrive and be grouped. Default: 30~seconds.

\item[Group Interval (\texttt{group\_interval})]
The minimum interval between subsequent notifications for the same group. Default: 5~minutes.

\item[Drain3]
An online log parsing algorithm that extracts log templates from raw log streams by clustering log messages based on token similarity, producing parameterized templates with variable placeholders~\cite{drain3}.

\item[Sentence-BERT (S-BERT)]
A modification of BERT that uses siamese and triplet network structures to derive semantically meaningful sentence embeddings, optimized for cosine similarity comparison~\cite{sbert}.

\item[HDBSCAN]
Hierarchical Density-Based Spatial Clustering of Applications with Noise; a density-based clustering algorithm that does not require specifying the number of clusters and can identify outliers (noise points)~\cite{hdbscan}.

% =============================================================================
% EXPERIMENTAL DESIGN TERMS
% =============================================================================

\item[Ablation Study]
An experimental procedure where individual components (e.g., \texttt{get\_service\_dependency} tool, fast filter, runbook) are removed from the full system to measure their individual contribution to performance metrics.

\item[Cross-System Generalization]
The ability of an aggregation method trained or tuned on one microservice system (e.g., Online Boutique) to maintain performance on unseen systems (e.g., Train Ticket) without retraining.

\item[Statistical Significance Test]
Paired Wilcoxon signed-rank test (non-parametric) used to compare metric distributions across the four arms on the same failure cases, with significance level $\alpha = 0.05$.

\item[Metric Gaming]
The phenomenon where a method optimizes a single metric (e.g., ARR) at the expense of operational utility (e.g., grouping all alerts into one group achieves ARR $\approx$ 100\% but RCPR = 0\%). Mitigated by the multi-metric framework (ARR, RCPR, F1, Purity).

% =============================================================================
% INDUSTRY STANDARDS & REFERENCES
% =============================================================================

\item[Google SRE Principles]
Foundational Site Reliability Engineering practices documented in \textit{Site Reliability Engineering} (Ch.~6, 10, 11) and \textit{The Site Reliability Workbook} (Ch.~5), emphasizing actionable alerts, maximum 2 incidents per 12-hour shift, and alerting on SLOs rather than raw metrics~\cite{google-sre-book, google-sre-workbook}.

\item[Observability Pillars]
The three canonical telemetry types: \textit{Logs} (immutable event records), \textit{Metrics} (aggregated numerical measurements), and \textit{Traces} (distributed request flow records)~\cite{grafana2025}.

\end{description}

% =============================================================================
% LOCAL BIBLIOGRAPHY MACROS (if not using main document's .bib)
% Include these in your main .bib or define via IEEEabrv.bib/IEEEfull.bib
% =============================================================================
% @string{srebook = "Site Reliability Engineering: How Google Runs Production Systems"}
% @string{sreworkbook = "The Site Reliability Workbook: Practical Ways to Implement SRE"}
% @article{pagerduty2020, ...}
% @article{grafana2025, ...}
% @inproceedings{kuang2024cola, ...}
% @inproceedings{pham2025rcaeval, ...}
% @article{zheng2024lemma, ...}
% @inproceedings{yu2025alertguardian, ...}
% @manual{prometheus-alertmanager, ...}
% @article{sbert, ...}
% @article{drain3, ...}
% @article{hdbscan, ...}
% @article{wei2022cot, ...}